\chapter{从整体积分恢复局部结构：微分、变差与卷积}\label{ch:04}

积分把局部信息压缩成总体平均，分析中的许多反问题却沿相反方向行进：怎样从越来越小的平均恢复
一点的值，怎样从区间增量恢复函数的导数，怎样将粗糙对象平滑并在原来的范数空间中逼近它，又怎样
从一族有限时间演化读出瞬时生成元？形式微分对这些问题帮助有限；跳跃、奇异质量、边界泄漏和
无界算子的定义域都会迫使我们把缺失的条件写出来。

第一条恢复链从覆盖几何开始。重叠小球使局部异常无法直接求和，Vitali 选择和
Hardy--Littlewood 极大估计把问题转化为一次质量计算，继而得到 Lebesgue 点和测度密度。第二条链
留在一维：总变差把所有分割增量累计成测度，Stieltjes 积分同时容纳密度、原子与 Cantor 型奇异
部分，绝对连续性则准确刻画何时普通导数足以恢复原函数。第三条链把微分移到测试函数上；
mollifier、分布和 Sobolev 空间让不光滑函数仍可参与微分方程的能量估计。第四条链从平移与卷积
抽象出近似恒等和强连续半群，并在这一框架下证明生成元的稠密性与闭性，以及正实参数的 Laplace
预解式。

热核、Gaussian 测度与 Markov 性也可分别用概率论和 Fourier 分析解释，相应理论将在后续章节展开。
Besicovitch 微分定理在本章只作选读说明；Rellich--Kondrachov 定理以选读形式证明；对
Hille--Yosida 理论则只讨论从给定半群导出的必要结论。

\section{积分的微分：从平均值恢复点值}\label{sec:04-01}

\subsection{Vitali 覆盖与 Hardy--Littlewood 极大函数}\label{subsec:4-1-1}
\index{Hardy--Littlewood 极大函数}\index{Vitali 覆盖}

设 $f$ 是局部可积函数。若只看一个固定半径 $r$，球平均
\[
  \frac{1}{m(B(x,r))}\int_{B(x,r)}|f(y)|\,\dd y
\]
并不难处理。困难来自“同时看所有半径”：一个点只要在某个尺度上平均过大，就应当算作坏点。
每个坏点都带来一个见证球，而这些球可能层层重叠。把它们的体积直接相加，会把同一份函数质量
重复计算任意多次。我们需要先从见证球中挑出互不相交的一族，同时又不丢掉原来的坏点。

\begin{definition}[中心极大函数]\label{def:4-1-1-1}
对 $f\in L^1_{\mathrm{loc}}(\R^n)$，定义其\emph{中心 Hardy--Littlewood 极大函数}
\[
  Mf(x)=\sup_{r>0}\frac{1}{m(B(x,r))}
                \int_{B(x,r)}|f(y)|\,\dd y .
\]
这里及下文 $B(x,r)$ 表示 Euclidean 开球，并记 $\omega_n=m(B(0,1))$。
\end{definition}

这个定义中的上确界确实给出可测函数。固定 $r>0$，令
$I_r(x)=\int_{B(x,r)}|f|$。当 $x_k\to x$ 时，可把所有相关球限制在某个有界集合 $Q$ 内；
$|f|\1_Q\in L^1$。若 $h_k=|x_k-x|<r$，则
\[
 B(x_k,r)\mathbin\triangle B(x,r)
 \subset B(x,r+h_k)\setminus\overline{B(x,r-h_k)},
\]
因而
\[
 m\bigl(B(x_k,r)\mathbin\triangle B(x,r)\bigr)
 \le \omega_n\bigl((r+h_k)^n-(r-h_k)^n\bigr)\longrightarrow0.
\]
由积分的绝对连续性，
\[
 |I_r(x_k)-I_r(x)|
 \le \int_{B(x_k,r)\triangle B(x,r)}|f(y)|\,\dd y\longrightarrow0.
\]
因此每个固定半径的平均关于 $x$ 连续，连续函数族的上确界 $Mf$ 下半连续；特别地，
$\{Mf>\lambda\}$ 是开集。另一方面，若 $r_k\to r>0$，则
\[
 B(x,r_k)\triangle B(x,r)
 \subset B\bigl(x,r+|r_k-r|\bigr)
       \setminus\overline{B\bigl(x,r-|r_k-r|\bigr)},
\]
所以同一估计证明球积分关于 $r$ 连续。因此定义中的上确界也可只取
$r\in\Q_{>0}$。

\begin{definition}[Vitali 球族]\label{def:4-1-1-2}
称球族 $\mathcal V$ 以 Vitali 意义覆盖 $E\subset\R^n$，若对每个 $x\in E$ 和
$\delta>0$，$\mathcal V$ 中都有一个以 $x$ 为中心、半径小于 $\delta$ 的球。
\end{definition}

例如，开集 $G$ 内所有满足 $B(x,r)\Subset G$ 的中心球构成 $G$ 的 Vitali 覆盖。
相反，只允许一个固定半径的中心球族不满足“半径任意小”的要求。
下面的选择引理只要求球的半径有统一上界，并不要求半径可以任意小。任意小半径条件属于微分基的
局部性质，二者应予区分。

\begin{definition}[弱 $L^1$ 空间]\label{def:4-1-1-3}
可测函数 $g$ 属于弱 $L^1$，记作 $g\in L^{1,\infty}$，若
\[
  \|g\|_{L^{1,\infty}}
  :=\sup_{\lambda>0}\lambda\,m\{x:|g(x)|>\lambda\}<\infty.
\]
这个泛函通常称为弱 $L^1$ 准范数；它一般不满足范数的三角不等式。
\end{definition}

先看一个能直接算出的极大函数。

\begin{example}[球的示性函数]\label{ex:4-1-1-1}
令 $f=\1_{B(0,1)}$。若 $x\in B(0,1)$，任意充分小的 $B(x,r)$ 都包含在单位球内，故
$Mf(x)=1$。若 $|x|>2$，取半径 $|x|+1$ 的中心球，它包含 $B(0,1)$，于是
\[
  Mf(x)\ge \frac{m(B(0,1))}{m(B(x,|x|+1))}=(|x|+1)^{-n}.
\]
反之，任何与 $B(0,1)$ 相交的 $B(x,r)$ 都满足 $r>|x|-1$，从而
\[
 \frac{m(B(x,r)\cap B(0,1))}{m(B(x,r))}
 \le (|x|-1)^{-n}.
\]
不相交的球给出平均值 $0$，所以
\[
 (|x|+1)^{-n}\le Mf(x)\le(|x|-1)^{-n}\qquad(|x|>2).
\]
特别地，$Mf(x)$ 在无穷远与 $|x|^{-n}$ 同阶；用极坐标看，
$\int_{|x|>2}|x|^{-n}\,\dd x$ 含有 $\int_2^\infty r^{-1}\,\dd r$，因而发散。
所以极大算子不可能把 $L^1$ 有界地映入 $L^1$。
\end{example}

\begin{example}[高度重叠的球族]\label{ex:4-1-1-2}
固定 $N\ge1$，取
\[
 B_k=B(0,2^{-k}),\qquad 1\le k\le N,
 \qquad f_N=\1_{B(0,2^{-N})}.
\]
这些球全部含有 $B_N$，因而
\[
 \sum_{k=1}^N\int_{B_k}f_N\,\dd x
 =N\,m(B_N),
 \qquad
 \int_{\bigcup_{k=1}^NB_k}f_N\,\dd x=m(B_N).
\]
直接对见证球求和已把同一份质量重复计算 $N$ 次。甚至把每个球复制有限多次也不改变
它们的并，却能任意增大体积之和。覆盖估计若不先做选择，便不可能只依赖被覆盖集合。
\end{example}

\begin{lemma}[$5r$ 覆盖引理]\label{lemma:4-1-1-1}
设 $\mathcal V$ 是 $\R^n$ 中半径有上界的球族。存在可数个两两不交的球
$B_j\in\mathcal V$，使得
\[
  \bigcup_{B\in\mathcal V}B\subset\bigcup_{j=1}^{\infty}5B_j.
\]
而且每个 $B\in\mathcal V$ 都与某个满足
$\operatorname{rad}(B_j)\ge\tfrac12\operatorname{rad}(B)$ 的 $B_j$ 相交。
\end{lemma}

\begin{proof}
先设 $\mathcal V$ 有限。按半径不增的次序检查各球：若当前球与已经选出的球都不交，就选入；
否则舍去。每个舍去的球 $B=B(x,r)$ 都与一个较早选出的
$B_j=B(x_j,r_j)$ 相交，并有 $r_j\ge r$。

一般情形取所有半径的一个上界 $R$，并按
\[
 \mathcal V_k=\{B\in\mathcal V:2^{-k-1}R<\operatorname{rad}(B)\le2^{-k}R\},
 \qquad k=0,1,\ldots,
\]
分层。从 $k=0$ 起，依次在 $\mathcal V_k$ 中选取一个极大的两两不交子族，同时要求它与此前层中
已选球都不交。这样的极大子族可由集合包含关系下的极大原理取得。任一两两不交的正半径球族都是
可数的：给每个球选一个不同的有理坐标点即可。因此全部选出的球可编号为 $(B_j)$。

若 $B\in\mathcal V_k$ 未被选入，极大性说明它与第 $k$ 层或更早一层的某个已选球 $B_j$ 相交。
若 $r,r_j$ 为二者半径，则 $r_j>2^{-k-1}R\ge r/2$。取交点 $w$。对任意 $z\in B$，
\[
 |z-x_j|\le |z-x|+|x-w|+|w-x_j|<2r+r_j\le5r_j,
\]
故 $B\subset5B_j$。已选球自身当然也包含在相应的 $5B_j$ 中，结论成立。
\end{proof}

有限球族的贪心选择甚至给出更好的放大常数；这里保留 $5$，是为了让有限与一般球族共用同一陈述。

\begin{theorem}[Hardy--Littlewood 弱型 $(1,1)$]\label{theorem:4-1-1-2}
对每个 $f\in L^1(\R^n)$ 和 $\lambda>0$，有
\[
 m\{Mf>\lambda\}\le \frac{5^n}{\lambda}\|f\|_1.
\]
特别地，$M:L^1(\R^n)\to L^{1,\infty}(\R^n)$。
\end{theorem}

\begin{proof}
记 $E_\lambda=\{Mf>\lambda\}$。上文已经证明它是开集。先取紧集
$K\subset E_\lambda$。对每个 $x\in K$，按极大函数的定义选 $r_x>0$，使
\[
  \int_{B(x,r_x)}|f|>\lambda m(B(x,r_x)).
\]
这些见证球覆盖 $K$，紧性给出有限子覆盖。对这有限族应用
\cref{lemma:4-1-1-1}，取得两两不交的 $B_1,\ldots,B_N$，且
$K\subset\bigcup_{j=1}^N5B_j$。于是
\[
 \begin{aligned}
 m(K)&\le\sum_{j=1}^N m(5B_j)
      =5^n\sum_{j=1}^Nm(B_j)\\
 &<\frac{5^n}{\lambda}\sum_{j=1}^N\int_{B_j}|f|
 \le\frac{5^n}{\lambda}\|f\|_1,
 \end{aligned}
\]
最后一步正是选择不交球的用途。Lebesgue 测度对开集内正则，故对所有紧集
$K\subset E_\lambda$ 取上确界即得所述估计。
\end{proof}

\begin{remark}[一维覆盖常数的改进]\label{ex:4-1-1-3}
在 $\R$ 上，球就是区间。若改用 rising sun lemma，可按开集分量选择区间，得到比 $5$ 更好的常数。
球覆盖证明则可直接用于随后的 $\R^n$ 微分定理；一维改进见本小节习题。
\end{remark}

\begin{proposition}[强型 $(p,p)$（选读）]\label{proposition:4-1-1-3}
若 $1<p\le\infty$，则存在只依赖 $n,p$ 的常数 $C_{n,p}$，使
\[
  \|Mf\|_p\le C_{n,p}\|f\|_p\qquad(f\in L^p(\R^n)).
\]
\end{proposition}

\begin{proof}
$p=\infty$ 时，每个球平均不超过 $\|f\|_\infty$。设 $1<p<\infty$。对固定
$\lambda>0$ 分解
 \[
 f=f_\lambda^{(0)}+f_\lambda^{(1)},\qquad
 f_\lambda^{(0)}=f\1_{\{|f|\le\lambda/2\}},\qquad
 f_\lambda^{(1)}=f\1_{\{|f|>\lambda/2\}}.
 \]
由于
\[
 \int_{\{|f|>\lambda/2\}}|f|
 \le (2/\lambda)^{p-1}\|f\|_p^p,
\]
$f_\lambda^{(1)}\in L^1$，所以可以对它使用弱型估计。
由 $M f_\lambda^{(0)}\le\lambda/2$ 及 $M(f+g)\le Mf+Mg$，
\[
 \{Mf>\lambda\}\subset
 \{M f_\lambda^{(1)}>\lambda/2\}.
\]
弱型估计给出
\[
 m\{Mf>\lambda\}
 \le\frac{2\cdot5^n}{\lambda}
       \int_{\{|f|>\lambda/2\}}|f(x)|\,\dd x.
\]
层蛋糕公式和 Tonelli 定理于是给出
\[
 \begin{aligned}
 \|Mf\|_p^p
 &=p\int_0^\infty\lambda^{p-1}m\{Mf>\lambda\}\,\dd\lambda\\
 &\le2\cdot5^np\int_{\R^n}|f(x)|
     \left(\int_0^{2|f(x)|}\lambda^{p-2}\,\dd\lambda\right)\dd x\\
 &=\frac{2^p5^np}{p-1}\|f\|_p^p.
 \end{aligned}
\]
取 $p$ 次方根即得结论。
\end{proof}

\begin{proposition}[局部极大函数]\label{proposition:4-1-1-4}
设 $K\subset\R^n$ 有界可测，$r_0>0$，并记
$K_{r_0}=\{y:\dist(y,K)<r_0\}$。若 $f\in L^1(K_{r_0})$，则
对 $x\in K$ 定义
\[
 M_{r_0}f(x)=\sup_{0<r<r_0}\frac1{m(B(x,r))}\int_{B(x,r)}|f(y)|\,\dd y
\]
满足
\[
 m\{x\in K:M_{r_0}f(x)>\lambda\}
 \le\frac{5^n}{\lambda}\int_{K_{r_0}}|f|.
\]
\end{proposition}

\begin{proof}
把 $f$ 在 $K_{r_0}$ 外延为零，所得函数记为 $F\in L^1(\R^n)$。若 $x\in K$ 且
$0<r<r_0$，则 $B(x,r)\subset K_{r_0}$，所以相应的 $f$ 平均等于 $F$ 平均。于是
$M_{r_0}f(x)\le MF(x)$，结论由全局弱型估计得到。
\end{proof}

若 $K\Subset\Omega$，取 $r_0<\dist(K,\Omega^c)$ 就得到完全在 $\Omega$ 内取平均的版本。
这与先把 $f$ 在 $\Omega$ 外延为零、再允许任意半径的算子不同：后者会把边界外的零也纳入平均。

\begin{figure}[htbp]
\centering
\begin{tikzpicture}[scale=.82]
  \foreach \x/\y/\r in {0/0/1.15,.7/.15/.78,-.65/.25/.63,.15/.7/.52,
                           1.35/.15/.46,-1.2/-.2/.40,.55/-.65/.34}
    \draw[BookInk!30] (\x,\y) circle (\r);
  \draw[BookAccent,very thick] (-.65,.25) circle (.63);
  \draw[BookAccent,very thick] (1.35,.15) circle (.46);
  \draw[BookAccent,very thick] (.55,-.65) circle (.34);
  \begin{scope}[xshift=6.1cm]
    \draw[BookAccent!55,thick,dashed] (-.65,.25) circle (3.15);
    \draw[BookAccent!55,thick,dashed] (1.35,.15) circle (2.30);
    \draw[BookAccent!55,thick,dashed] (.55,-.65) circle (1.70);
    \fill[BookAccent] (-.65,.25) circle (2pt);
    \fill[BookAccent] (1.35,.15) circle (2pt);
    \fill[BookAccent] (.55,-.65) circle (2pt);
  \end{scope}
  \draw[book arrow] (2.1,0)--(3.8,0)
    node[midway,above,book label]{选择不交球并放大};
\end{tikzpicture}
\caption{左侧见证球高度重叠；选择出的三个实线球互不相交。右侧的五倍球覆盖原球族的并，
而积分质量只在实线球上累计一次。}
\label{fig:4-1-vitali-5r}
\end{figure}

弱型不等式的常数 $5^n$ 来自球放大后的体积比，而不是某个隐藏的解析恒等式。它和上面使用的
示性函数点态估计都把“大于阈值”的集合转化为总质量除以阈值；不过前者还要先解决所有尺度见证球的重叠。
卷积近似、Poisson 平均以及调和函数的边界值问题也会遇到同时控制所有尺度的困难。概率论中的最大
不等式具有相似的尾界形式，但其对象是一族随机变量，而不是一族 Euclidean 球平均。

\begin{exercise}
证明有限球族按半径递减贪心选择后，其余每个球都包含在某个已选球的三倍球中。
\end{exercise}
\begin{exercise}
补全 $I_r(x)$ 连续性的证明，并由此证明 $Mf$ 下半连续。
\end{exercise}
\begin{exercise}
用 \cref{ex:4-1-1-1} 证明不存在常数 $C$ 使 $\|Mf\|_1\le C\|f\|_1$ 对所有
$f\in L^1(\R^n)$ 成立。
\end{exercise}
\begin{exercise}
重新计算 $M\1_{B(0,1)}$：分别用“包住单位球”的球与“只要相交便有半径下界”证明
\[
 (|x|+1)^{-n}\le M\1_{B(0,1)}(x)\le(|x|-1)^{-n}
 \qquad(|x|>2).
\]
\end{exercise}
\begin{exercise}
证明一维 rising sun lemma 的下列形式。若 $F\in C([a,b])$，则
\[
 G=\{x\in(a,b):F(x)<\max_{x\le y\le b}F(y)\}
\]
是可数个不交开区间之并；对每个不碰 $b$ 的分量 $(\alpha,\beta)$，有
$F(\alpha)=F(\beta)$。把它用于
$F(x)=\int_a^x|f(t)|\dd t-\lambda x$，推导一侧区间极大函数的弱型估计，并比较球覆盖证明中的常数。
\end{exercise}

极大函数已经把“某个半径上的平均异常”压缩成一个可测坏集，并控制了它的大小。下一步是先对连续函数
验证小球平均确实回到点值，再让稠密逼近和弱型估计共同处理一般局部可积函数。

\subsection{Lebesgue 微分定理与密度点}\label{subsec:4-1-2}
\index{Lebesgue 微分定理}\index{Lebesgue 点}\index{密度点}

积分会在固定尺度上平滑局部细节；Lebesgue 微分定理说明，当平均半径趋于零时，原函数仍可在几乎
每一点恢复。连续点上的结论直接来自连续性。对一般 $L^1_{\mathrm{loc}}$ 函数，则需要把连续函数的
计算通过稠密逼近与极大估计传递过去。

\begin{definition}[Lebesgue 点]\label{def:4-1-2-1}
设 $f\in L^1_{\mathrm{loc}}(\R^n)$ 是一个选定的可测代表。若 $f(x)$ 有限且
\[
 \lim_{r\downarrow0}\frac1{m(B(x,r))}
       \int_{B(x,r)}|f(y)-f(x)|\,\dd y=0,
\]
则称 $x$ 为 $f$ 的一个 \emph{Lebesgue 点}。
\end{definition}

\begin{definition}[密度点]\label{def:4-1-2-2}
设 $E\subset\R^n$ 可测。若
\[
 \lim_{r\downarrow0}\frac{m(E\cap B(x,r))}{m(B(x,r))}=1,
\]
则称 $x$ 为 $E$ 的密度点。若同一极限为 $0$，称 $x$ 为 $E$ 的密度零点。
\end{definition}

\begin{definition}[平均算子]\label{def:4-1-2-3}
对 $r>0$，定义
\[
 A_rf(x)=\frac1{m(B(0,r))}\int_{B(x,r)}f(y)\,\dd y.
\]
积分只依赖 $f$ 的局部 $L^1$ 等价类。
\end{definition}

\begin{example}[连续点]\label{ex:4-1-2-1}
若 $f$ 在 $x$ 连续，则给定 $\varepsilon>0$，存在 $\delta>0$，使得只要
$|y-x|<\delta$，就有 $|f(y)-f(x)|<\varepsilon$。因此 $0<r<\delta$ 时
\[
 \frac1{m(B(x,r))}\int_{B(x,r)}|f(y)-f(x)|\,\dd y\le\varepsilon.
\]
所以每个连续点都是 Lebesgue 点。这是后面逼近论证中负责提供正确点值的部分。
\end{example}

\begin{theorem}[Lebesgue 微分定理]\label{theorem:4-1-2-1}
若 $f\in L^1_{\mathrm{loc}}(\R^n)$，则 $f$ 的几乎每个点都是 Lebesgue 点。特别地，
\[
  A_rf(x)\longrightarrow f(x)\qquad\text{对几乎每个 }x
\]
成立。同样的结论对以 $x$ 为中心、边平行于坐标轴的立方体平均成立。
\end{theorem}

\begin{proof}
对 $f$ 的无限值点重新赋值为零不会改变其局部 $L^1$ 等价类，因此不妨设 $f$ 处处有限。固定正整数
$R$，令
\[
 K=B(0,R),\qquad F=f\1_{B(0,R+1)}\in L^1(\R^n).
\]
当 $x\in K$ 且 $0<r<1$ 时，$B(x,r)\subset B(0,R+1)$，故在这个球上 $F=f$。
定义
\[
 \Omega_f(x)=\limsup_{r\downarrow0}\frac1{m(B(x,r))}
             \int_{B(x,r)}|f(y)-f(x)|\,\dd y.
\]
固定 $x$ 时，球积分关于 $r$ 连续，所以可在趋于零的正有理半径上取上极限；由此
$\Omega_f$ 可测。这里对每个固定的有理 $r$，函数
$(x,y)\mapsto\1_{\{|y-x|<r\}}|f(y)-f(x)|$ 可测，参数积分的可测性保证相应球平均关于 $x$ 可测。

由 \cref{theorem:3-3-3-1} 中 $C_c(\R^n)$ 在 $L^1(\R^n)$ 中的稠密性，给定 $\eta>0$ 可取
$g\in C_c(\R^n)$ 使 $\|F-g\|_1<\eta$。对 $x\in K$ 及 $r<1$，三角不等式给出
\[
 \begin{aligned}
 \frac1{m(B(x,r))}\int_{B(x,r)}|f(y)-f(x)|\,\dd y
 &\le \frac1{m(B(x,r))}\int_{B(x,r)}|F(y)-g(y)|\,\dd y\\
 &\quad+\frac1{m(B(x,r))}\int_{B(x,r)}|g(y)-g(x)|\,\dd y\\
 &\quad+|g(x)-F(x)|.
 \end{aligned}
\]
中间一项因 $g$ 在 $x$ 连续而随 $r\downarrow0$ 趋于零，故
\[
  \Omega_f(x)\le M(F-g)(x)+|F(x)-g(x)|\qquad(x\in K).
\]
于是对 $\lambda>0$，由 \cref{theorem:4-1-1-2}，以及点态估计
\[
  \1_{\{|F-g|>\lambda/2\}}\le \frac{2}{\lambda}|F-g|,
\]
\[
 \begin{aligned}
 m\{x\in K:\Omega_f(x)>\lambda\}
 &\le m\{M(F-g)>\lambda/2\}
      +m\{|F-g|>\lambda/2\}\\
 &\le \frac{2(5^n+1)}{\lambda}\|F-g\|_1
 <\frac{2(5^n+1)}{\lambda}\eta.
 \end{aligned}
\]
$\eta$ 可任意小，故这个集合测度为零。对 $\lambda=1/k$ 取可数并，得到
$\Omega_f=0$ 几乎处处于 $K$。再令 $R$ 遍历正整数，得到几乎每点都是 Lebesgue 点。

最后，
\[
 |A_rf(x)-f(x)|\le\frac1{m(B(x,r))}
       \int_{B(x,r)}|f(y)-f(x)|\,\dd y,
\]
所以普通平均收敛是 Lebesgue 点结论的推论。对中心立方体 $Q(x,r)=x+[-r,r]^n$，有
$Q(x,r)\subset B(x,\sqrt n\,r)$，并且
$m(B(x,\sqrt n\,r))/m(Q(x,r))=\omega_n n^{n/2}/2^n$。因此
\[
 \frac1{m(Q(x,r))}\int_{Q(x,r)}|f(y)-f(x)|\,\dd y
 \le \frac{\omega_n n^{n/2}}{2^n}
 \frac1{m(B(x,\sqrt n r))}
 \int_{B(x,\sqrt n r)}|f(y)-f(x)|\,\dd y,
\]
右端趋于零，故立方体版结论成立。
\end{proof}

上述证明中，连续函数确定候选极限，$C_c$ 稠密性提供逼近，而极大弱型估计统一控制所有小半径上的
逼近误差。稠密性本身并不足以推出逐点结论。

\begin{corollary}[Lebesgue 密度定理]\label{corollary:4-1-2-2}
若 $E\subset\R^n$ 可测，则 $E$ 中几乎每点是 $E$ 的密度点，$E^c$ 中几乎每点是
$E$ 的密度零点。
\end{corollary}

\begin{proof}
$\1_E\in L^1_{\mathrm{loc}}$，并且
\[
 A_r\1_E(x)=\frac{m(E\cap B(x,r))}{m(B(x,r))}.
\]
对 $\1_E$ 应用 \cref{theorem:4-1-2-1}。在 $E$ 上其值为 $1$，在 $E^c$ 上其值为 $0$，
即得结论。
\end{proof}

\begin{example}[区间端点]\label{ex:4-1-2-2}
对 $E=(0,1)\subset\R$，若 $x\in(0,1)$，充分小的 $(x-r,x+r)$ 包含在 $E$ 中，密度为 $1$；
若 $x\notin[0,1]$，密度为 $0$。在 $x=0$，当 $0<r<1$ 时
\[
 \frac{m((0,1)\cap(-r,r))}{2r}=\frac12,
\]
$x=1$ 同样如此。两个端点正是零测的异常集。
\end{example}

\begin{example}[肥 Cantor 集]\label{ex:4-1-2-3}
从 $F_0=[0,1]$ 开始，在第 $k$ 步从 $F_{k-1}$ 的每个区间分量中央删去一个长度为 $4^{-k}$
的开区间，并令 $F=\bigcap_{k\ge1}F_k$。第 $k$ 步删去 $2^{k-1}$ 个区间，故删去的总长度为
\[
 \sum_{k=1}^\infty 2^{k-1}4^{-k}=\frac12,
\]
从而 $m(F)=1/2$。第 $k$ 级的 $2^k$ 个分量等长，其共同长度为
\[
 \ell_k=2^{-k}\left(1-\sum_{j=1}^k2^{j-1}4^{-j}\right)
 \le2^{-k}\longrightarrow0.
\]
若某个非退化开区间
$I$ 包含在 $F$ 中，它便包含在每个 $F_k$ 的某个分量内，这与分量长度趋于零矛盾。所以 $F$ 没有内点。
它又是闭集，因此每个 $x\in F$ 都是拓扑边界点。可是
\cref{corollary:4-1-2-2} 说明，$F$ 中几乎每个 $x$ 满足
\[
 \frac{m(F\cap(x-r,x+r))}{2r}\longrightarrow1.
\]
拓扑问题问“每个邻域是否碰到补集”，密度问题问“补集占邻域的多少”；这个例子使二者不能混为一谈。
\end{example}

\begin{proposition}[精确代表的约定]\label{proposition:4-1-2-4}
设 $f\in L^1_{\mathrm{loc}}(\R^n)$。在有限极限存在处令
\[
 \widetilde f(x)=\lim_{r\downarrow0}A_rf(x),
\]
在其他点令 $\widetilde f(x)=0$。则 $\widetilde f=f$ 几乎处处；而在极限存在处，
$\widetilde f(x)$ 只由 $f$ 的 $L^1_{\mathrm{loc}}$ 等价类决定。
\end{proposition}

\begin{proof}
若 $f=g$ 几乎处处，则对每个球 $B(x,r)$，二者积分相等，故 $A_rf(x)=A_rg(x)$；所有存在的
平均极限也相等。由 \cref{theorem:4-1-2-1}，对几乎每个 $x$，平均极限存在且等于所选代表
$f(x)$，所以 $\widetilde f=f$ 几乎处处。异常集上取零只是约定；积分数据并未在那里选出唯一点值。
\end{proof}

\begin{figure}[htbp]
\centering
\begin{tikzpicture}[x=1cm,y=1cm]
  \draw[->,BookInk] (-.2,0)--(8.1,0) node[right]{$x$};
  \draw[->,BookInk] (0,-.2)--(0,3.1) node[above]{$y$};
  \draw[BookInk!65,thick,smooth] plot coordinates
   {(0,1.2) (.6,1.5) (1.2,.8) (1.8,2.5) (2.4,.9) (3,1.8) (3.6,.6)
    (4.2,2.6) (4.8,1.0) (5.4,2.1) (6,1.2) (6.6,1.8) (7.4,1.4)};
  \draw[BookAccent,very thick,smooth] plot coordinates
   {(0,1.25) (.8,1.3) (1.6,1.35) (2.4,1.4) (3.2,1.45) (4,1.5)
    (4.8,1.55) (5.6,1.6) (6.4,1.65) (7.4,1.7)};
  \fill[BookWarm!18] (3.45,0) rectangle (4.55,2.85);
  \draw[BookWarm,decorate,decoration={brace,mirror,amplitude=4pt}]
    (3.45,-.05)--(4.55,-.05) node[midway,below=5pt]{$B(x,r)$};
  \node[BookAccent] at (6.7,2.45) {$g$：连续近似};
  \node[BookInk!75] at (6.8,.65) {$f$};
\end{tikzpicture}
\caption{连续近似 $g$ 的平均振荡随半径缩小而消失；$f-g$ 在所有半径上的误差由
$M(f-g)$ 统一控制。}
\label{fig:4-1-differentiation-approximation}
\end{figure}

微分理论还给出 Riemann 可积性的精确判据：决定因素不是不连续点是否有限或稠密，而是其集合是否
为零测集。

对有界 $f:[a,b]\to\R$ 定义其在 $x$ 的振荡
\[
 \omega_f(x)=\inf_{\delta>0}
 \sup\{|f(y)-f(z)|:y,z\in[a,b],\ |y-x|<\delta,\ |z-x|<\delta\}.
\]
$f$ 在 $x$ 连续当且仅当 $\omega_f(x)=0$。此外
$D_k=\{x:\omega_f(x)\ge1/k\}$ 是闭集：若 $\omega_f(x)<1/k$，存在一个邻域，使其中任一点
都继承同一个更大的振荡控制，故该邻域不碰 $D_k$。

\begin{theorem}[Riemann 可积判据]\label{theorem:4-1-2-3}
有界函数 $f:[a,b]\to\R$ Riemann 可积，当且仅当它的不连续点集合具有 Lebesgue 测度零。
\end{theorem}

\begin{proof}
先设 $f$ Riemann 可积。固定 $k\ge1$ 与 $\varepsilon>0$，取分割
$P:a=x_0<\cdots<x_N=b$ 使
\[
 U(f,P)-L(f,P)=\sum_{j=1}^N\operatorname{osc}(f;[x_{j-1},x_j])
                       (x_j-x_{j-1})<\frac{\varepsilon}{2k},
\]
其中 $\operatorname{osc}(f;I)=\sup_I f-\inf_I f$。若
$x\in D_k$ 且 $x$ 不是分割点，则含 $x$ 的分割小区间的振荡至少为 $1/k$。所有这类小区间的
总长度因此不超过
\[
 k\bigl(U(f,P)-L(f,P)\bigr)<\varepsilon/2.
\]
详细地，若 $x\in(x_{j-1},x_j)$，取
$0<\delta<\min\{x-x_{j-1},x_j-x\}$。由 $\omega_f(x)\ge1/k$，
\[
 \operatorname{osc}(f;[x_{j-1},x_j])
 \ge \sup_{\substack{|y-x|<\delta\\|z-x|<\delta}}|f(y)-f(z)|
 \ge\frac1k.
\]
因而把不含分割点的 $D_k$ 覆盖起来的小区间总长度确实至多为
$k(U(f,P)-L(f,P))$。
有限个分割点可由总长度小于 $\varepsilon/2$ 的有限开区间覆盖。所以 $D_k$ 的外测度小于
$\varepsilon$，即 $m(D_k)=0$。不连续点集为 $\bigcup_{k\ge1}D_k$，故它是零测集。

反之，设不连续点集 $D$ 测度为零，并令 $B=\sup_{[a,b]}|f|$。若 $B=0$，结论成立；以下设
$B>0$。给定 $\varepsilon>0$，由 Lebesgue 外正则性取开集 $G\supset D$，使
$m(G)<\varepsilon/(4B)$，并令 $K=[a,b]\setminus G$。对每个 $x\in K$，$f$ 在 $x$ 连续，
故可取 $\delta_x>0$，使当 $|y-x|<2\delta_x$ 时
\[
 |f(y)-f(x)|<\frac{\varepsilon}{4(b-a)}.
\]
有限个 $B(x_i,\delta_{x_i})$ 覆盖紧集 $K$。取分割 $P$ 的网格小于
$\min_i\delta_{x_i}$。若分割小区间 $I$ 与 $K$ 相交，取
$y\in I\cap K$，再选 $i$ 使 $y\in B(x_i,\delta_{x_i})$。对每个 $z\in I$，
$|z-x_i|\le |z-y|+|y-x_i|<2\delta_{x_i}$，故 $I\subset B(x_i,2\delta_{x_i})$。
于是这类 $I$ 上
\[
 \operatorname{osc}(f;I)<\frac{\varepsilon}{2(b-a)}.
\]
其余小区间不碰 $K$，因而包含于 $G$，它们的总长度至多 $m(G)$。所以
\[
 \begin{aligned}
 U(f,P)-L(f,P)
 &\le \frac{\varepsilon}{2(b-a)}(b-a)+2B\,m(G)\\
 &<\varepsilon.
 \end{aligned}
\]
这里第一类小区间的长度和至多 $b-a$；第二类小区间两两仅在端点相交，且其并包含于 $G$，
故它们的长度和不超过 $m(G)$；上式中的两个估计因此都已显式核对。
Darboux 判据遂给出 $f$ Riemann 可积。
\end{proof}

\begin{example}[Thomae 函数与 Dirichlet 函数]
在 $[0,1]$ 上令
\[
 t(x)=\begin{cases}q^{-1},&x=p/q\text{ 为既约分数},\\0,&x\notin\Q,
 \end{cases}
 \qquad d(x)=\1_{\Q}(x).
\]
$t$ 在每个有理点不连续，因为无理数列趋近该点而函数值恒为零。若 $x$ 无理，给定
$\varepsilon>0$，分母不超过 $1/\varepsilon$ 的既约分数只有有限多个；取一个避开这些点的
$x$ 的小邻域，其中 $t<\varepsilon$，故 $t$ 在 $x$ 连续。因此 $t$ 的不连续点集可数，
由 \cref{theorem:4-1-2-3} 它 Riemann 可积；每个非退化区间含无理数，所以所有 Darboux 下和为
$0$，其积分也为 $0$。函数 $d$ 在每个区间内同时取 $0$ 与 $1$，故每个分割的上、下和分别为
$1$ 和 $0$；它处处不连续，其不连续集测度为 $1$，因而不 Riemann 可积。这一对照也说明
“坏点稠密”并非判据。
\end{example}

球平均的形状可以换成上述中心立方体，也可以在有适当覆盖性质的微分基中推广；任意一列缩向点的
可测集合则不受本定理保证。下一小节把被平均的函数换成测度，检验局部质量比究竟恢复哪一部分信息。

\begin{exercise}
直接由定义证明连续点是 Lebesgue 点，并证明 Lebesgue 点处 $A_rf(x)\to f(x)$。
\end{exercise}
\begin{exercise}
计算 $E=[0,1]\cup[2,3]$ 的每个点的密度（只需讨论区间内部、端点和外部）。
\end{exercise}
\begin{exercise}
在 Riemann 可积判据的必要性证明中，详细证明若 $x\in D_k$ 位于分割小区间内部，则该区间
振荡至少为 $1/k$。
\end{exercise}
\begin{exercise}
对 $f=\1_E$ 应用 Lebesgue 微分定理，完整推出 Lebesgue 密度定理；须分别说明 $x\in E$ 与
$x\notin E$ 时平均值的极限。
\end{exercise}
\begin{exercise}
独立完成 Riemann 可积判据的充分性方向：把不连续点放进小测度开集，并在其补集上用有限连续性
邻域控制 Darboux 上、下和之差。
\end{exercise}

\subsection{测度的局部密度与抽象微分}\label{subsec:4-1-3}
\index{测度的密度}\index{Lebesgue 分解}

函数导数是增量之比的极限。若把函数换成测度，最直接的候选便是
\[
 \frac{\nu(B(x,r))}{m(B(x,r))}.
\]
三类测度说明了这一比值的不同表现。若 $\nu=fm$ 且 $f$ 连续，该比值是 $f$ 的球平均，因而趋于
$f(x)$；若 $\nu=\delta_0$，它在 $0$ 发散，在其他点最终为零；若质量集中在 Cantor 集上，则它在
Lebesgue 几乎处处为零。由此可见，局部体积比恢复绝对连续部分，却不能以普通函数记录奇异质量。

\begin{definition}[测度密度比]\label{def:4-1-3-1}
对 $\R^n$ 上的局部有限正 Borel 测度 $\nu$，定义
\[
 D_r\nu(x)=\frac{\nu(B(x,r))}{m(B(x,r))}.
\]
若 $\lim_{r\downarrow0}D_r\nu(x)$ 存在，则称它为 $\nu$ 相对于 Lebesgue 测度的对称密度。
对局部有限有符号或复测度，同一定义用于其有限取值的球质量。
\end{definition}

\begin{definition}[中心微分基与相对微分]\label{def:4-1-3-2}
设 $\mu$ 是局部有限正 Radon 测度。一个中心微分基 $\mathcal D$ 为每个 $x$ 指定一族含 $x$ 的有界
Borel 集 $\mathcal D(x)$，且其中集合的直径可以任意小。记 $V\to x$ 表示
$V\in\mathcal D(x)$、$\mu(V)>0$ 且 $\diam V\to0$。

若 $\nu$ 是局部有限有符号 Radon 测度且 $\nu\ll\mu$，则沿该微分基考察相对比值
\[
 \frac{\nu(V)}{\mu(V)}.
\]
其极限称为 $\nu$ 关于 $\mu$ 沿 $\mathcal D$ 的相对微分。分母所用的测度、集合缩向点的方式以及
“几乎处处”所相对的测度都是定义的一部分。
\end{definition}

\begin{definition}[上、下密度]\label{def:4-1-3-3}
对正测度 $\nu$，定义
\[
 \overline D\nu(x)=\limsup_{r\downarrow0}D_r\nu(x),\qquad
 \underline D\nu(x)=\liminf_{r\downarrow0}D_r\nu(x).
\]
二者允许取 $+\infty$。
\end{definition}

\begin{example}[连续密度]\label{ex:4-1-3-1}
若 $\nu=fm$，其中 $f\ge0$ 连续，则
\[
 D_r\nu(x)=A_rf(x)\longrightarrow f(x)
\]
在每个点成立。若只假设 $f\in L^1_{\mathrm{loc}}$，同一计算仍成立，只是极限由
\cref{theorem:4-1-2-1} 保证在几乎每一点成立。
\end{example}

\begin{example}[Dirac 测度]\label{ex:4-1-3-2}
对 $x\ne0$，当 $0<r<|x|$ 时 $\delta_0(B(x,r))=0$，故
$D_r\delta_0(x)$ 最终为零。在 $x=0$，
\[
 D_r\delta_0(0)=\frac1{\omega_nr^n}\longrightarrow+\infty.
\]
因此 $D_r\delta_0\to0$ 对 $m$-几乎每一点成立，而 $\delta_0(\R^n)=1$。密度为零几乎处处
不能推出原测度为零。
\end{example}

\begin{example}[Cantor 测度]\label{ex:4-1-3-3}
令 $C\subset[0,1]$ 为三分 Cantor 集，$\mu_C$ 为标准 Cantor 概率：每个第 $k$ 级基本区间
质量为 $2^{-k}$。记 $\alpha=\log2/\log3$。若
$3^{-(k+1)}<r\le3^{-k}$ 且 $x\in C$，含 $x$ 的第 $k+1$ 级基本区间长度小于 $r$，故整个
区间包含在 $B(x,r)$ 中，故
\[
 \mu_C(B(x,r))\ge2^{-(k+1)}\ge\frac12r^\alpha.
\]
另一方面，两个不同的第 $k$ 级基本区间之间的距离至少为 $3^{-k}$，而
$B(x,r)$ 的长度不超过 $2\cdot3^{-k}$，所以它至多与三个第 $k$ 级基本区间相交。从而
\[
 \mu_C(B(x,r))\le3\cdot2^{-k}<6r^\alpha,
\]
最后一步用了 $r>3^{-(k+1)}$。因此可以显式取 $c=1/2$ 和 $C=6$，使
\[
 c r^\alpha\le\mu_C(B(x,r))\le C r^\alpha
 \qquad(x\in C,\ 0<r<1).
\]
这里使用了 $2^{-k}=(3^{-k})^\alpha$。于是对 $x\in C$，
\[
 \frac14r^{\alpha-1}
 \le\frac{\mu_C(B(x,r))}{m(B(x,r))}
 \le3r^{\alpha-1},
 \qquad
 \frac{\mu_C(B(x,r))}{m(B(x,r))}\longrightarrow+\infty.
\]
可是 $m(C)=0$，且 $C$ 闭；对 $m$-几乎每个 $x$，事实上对每个 $x\notin C$，充分小的球不碰
$C$，所以该比值最终为零。同一测度在 $m$-几乎处处不可见，却在自身支撑上的每一点高度集中。
\end{example}

\begin{figure}[htbp]
\centering
\begin{tikzpicture}[x=1.05cm,y=.82cm]
  \draw[->,BookInk!75] (0,0)--(8.1,0)
    node[right,book label] {$s=-\log r$（向右即 $r\downarrow0$）};
  \draw[->,BookInk!75] (0,-.25)--(0,5.2)
    node[above,book label] {$\log D_r$};
  \draw[BookAccent,very thick] (.5,1.25)--(7.4,1.25)
    node[right,book label] {$fm$：$\log f(x)$};
  \draw[BookWarning,very thick] (.5,.35)--(7.3,4.7)
    node[right,book label] {$\delta_0$ 在 $0$：斜率 $n$};
  \draw[BookWarm,very thick,dashed] (.5,.7)--(7.3,3.2)
    node[right,book label] {$\mu_C$ 在 $C$：斜率 $1-\alpha$};
  \node[book node,align=left] at (3.8,4.55)
    {连续密度趋于有限值；\\奇异质量在自身支撑上发散};
\end{tikzpicture}
\caption{三种局部质量比置于同一对数尺度。水平线按 $f(x)>0$ 绘制；
在 Dirac 或 Cantor 支撑之外，相应比值对充分小的 $r$ 恒为零。图展示渐近阶，具体结论仍由
例~\ref{ex:4-1-3-1} 至例~\ref{ex:4-1-3-3} 的计算给出。}
\label{fig:4-1-density-scale}
\end{figure}

\begin{theorem}[绝对连续测度微分]\label{theorem:4-1-3-1}
若 $f\in L^1_{\mathrm{loc}}(\R^n)$ 且 $\nu=fm$，则
\[
 \lim_{r\downarrow0}\frac{\nu(B(x,r))}{m(B(x,r))}=f(x)
\]
对 $m$-几乎每个 $x$ 成立。
\end{theorem}

\begin{proof}
对每个球，
\[
 \frac{\nu(B(x,r))}{m(B(x,r))}
 =\frac1{m(B(x,r))}\int_{B(x,r)}f(y)\,\dd y=A_rf(x).
\]
结论正是 \cref{theorem:4-1-2-1}。
\end{proof}

为了处理奇异部分，先记录弱型证明并不需要被平均的质量具有密度。若 $\sigma$ 是有限正 Radon 测度，
定义
\[
 M\sigma(x)=\sup_{r>0}\frac{\sigma(B(x,r))}{m(B(x,r))}.
\]
固定 $r$ 时 $x\mapsto\sigma(B(x,r))$ 下半连续：若 $x_j\to x$，则对每个 $0<s<r$，
$B(x,s)\subset B(x_j,r)$ 当 $j$ 充分大，故
\[
 \liminf_j\sigma(B(x_j,r))\ge\sigma(B(x,s));
\]
再令 $s\uparrow r$，便知每个固定半径的质量函数下半连续；取上确界后
$M\sigma$ 仍下半连续，因而可测且其严格上水平集为开集。为把函数情形的弱型估计移到这里，记
$E_\lambda=\{M\sigma>\lambda\}$，并固定紧集 $K\subset E_\lambda$。对每个 $x\in K$ 选取
见证球 $B(x,r_x)$，使
\[
  \sigma(B(x,r_x))>\lambda m(B(x,r_x)).
\]
由紧性可先取有限子覆盖，再对这个有限球族应用
\cref{lemma:4-1-1-1}，得到两两不交的见证球
$B_1,\ldots,B_N$，且 $K\subset\bigcup_{j=1}^N5B_j$。于是
\[
 \begin{aligned}
 m(K)&\le\sum_{j=1}^Nm(5B_j)
      =5^n\sum_{j=1}^Nm(B_j)\\
 &<\frac{5^n}{\lambda}\sum_{j=1}^N\sigma(B_j)
 \le\frac{5^n}{\lambda}\sigma(\R^n).
 \end{aligned}
\]
最后对 $E_\lambda$ 的紧子集取上确界；$E_\lambda$ 是开集，Lebesgue 测度的内正则性给出
\begin{equation}\label{eq:4-1-measure-maximal}
 m\{M\sigma>\lambda\}\le\frac{5^n}{\lambda}\sigma(\R^n).
\end{equation}
需要注意，若奇异测度集中在零测集 $N$ 上，取开集 $U\supset N$ 并不能使
$\nu_s(U)$ 变小，因为 $U$ 仍含有全部奇异质量。下面的证明将改为截去承载绝大部分质量的紧集，
再对剩余测度应用极大估计。

\begin{theorem}[Lebesgue 分解的局部图像]\label{theorem:4-1-3-2}
设 $\nu$ 是 $\R^n$ 上的局部有限正 Radon 测度，且
\[
 \nu=fm+\nu_s,\qquad \nu_s\perp m,\qquad f\in L^1_{\mathrm{loc}}(\R^n)
\]
是它的 Lebesgue 分解。则
\[
 D_r\nu(x)\longrightarrow f(x)
\]
对 $m$-几乎每个 $x$ 成立。特别地，$D_r\nu_s(x)\to0$ 对 $m$-几乎每个 $x$ 成立。
\end{theorem}

\begin{proof}
绝对连续部分由 \cref{theorem:4-1-3-1} 处理，只需证明关于 $\nu_s$ 的最后一句。
取 Borel 零测集 $N$，使 $\nu_s(\R^n\setminus N)=0$。固定正整数 $R$，令
$W=B(0,R+1)$ 及 $\sigma=\nu_s|_W$。这是有限正 Radon 测度。给定 $\varepsilon>0$，
内正则性给出紧集 $K\subset N\cap W$，使
\[
 \sigma(W\setminus K)<\varepsilon.
\]
令 $\tau=\sigma|_{W\setminus K}$，则 $\tau(\R^n)<\varepsilon$。

若 $x\in B(0,R)\setminus K$，由于 $K$ 闭且 $B(0,R)$ 到 $W^c$ 的距离为 $1$，存在
$r_x>0$，使 $0<r<r_x$ 时 $B(x,r)\subset W\setminus K$。若 $K\ne\varnothing$，可取
\[
 r_x=\frac12\min\{1,\dist(x,K)\}>0.
\]
若 $K=\varnothing$，则直接取 $r_x=1/2$。
因此
\[
 \nu_s(B(x,r))=\sigma(B(x,r))=\tau(B(x,r))
 \qquad(0<r<r_x).
\]
对 $\lambda>0$，这说明
\[
 \left\{x\in B(0,R)\setminus K:
   \limsup_{r\downarrow0}D_r\nu_s(x)>\lambda\right\}
 \subset\{M\tau>\lambda\}.
\]
由于 $m(K)=0$，由 \eqref{eq:4-1-measure-maximal} 得
\[
 m\left\{x\in B(0,R):
   \limsup_{r\downarrow0}D_r\nu_s(x)>\lambda\right\}
 \le\frac{5^n\varepsilon}{\lambda}.
\]
令 $\varepsilon\downarrow0$，左侧集合为零测集。再对 $\lambda=1/k$ 取可数并，得到
$\limsup_{r\downarrow0}D_r\nu_s=0$ 几乎处处于 $B(0,R)$。密度比非负，所以其下极限不小于
$0$，从而极限为 $0$。最后令 $R\uparrow\infty$，得到全空间结论。

在绝对连续部分与奇异部分各自的满测集合之交上，
$D_r\nu=D_r(fm)+D_r\nu_s\to f+0$，证明完成。
\end{proof}

\begin{proposition}[相对密度与 Radon--Nikodym 导数]\label{proposition:4-1-3-3}
设 $\mathcal D$ 是关于局部有限正 Radon 测度 $\mu$ 的一个中心微分基，并假定它能够微分
$L^1_{\mathrm{loc}}(\mu)$：即对每个 $h\in L^1_{\mathrm{loc}}(\mu)$，沿 $V\in\mathcal D(x)$
缩向 $x$ 时
\[
 \frac1{\mu(V)}\int_Vh\,\dd\mu\longrightarrow h(x)
\]
对 $\mu$-几乎每个 $x$ 成立。若 $\nu=h\mu\ll\mu$，则沿同一微分基
\[
 \frac{\nu(V)}{\mu(V)}\longrightarrow h(x)
 =\frac{\dd\nu}{\dd\mu}(x)
\]
对 $\mu$-几乎每个 $x$ 成立。
\end{proposition}

\begin{proof}
在 $\mu(V)>0$ 的基集合上，Radon--Nikodym 表示直接给出
\[
 \frac{\nu(V)}{\mu(V)}
 =\frac1{\mu(V)}\int_Vh\,\dd\mu.
\]
右端由题设的微分基性质趋于 $h(x)$。因此，除 Radon--Nikodym 表示外，还必须假设
$\mathcal D$ 能够微分 $L^1_{\mathrm{loc}}(\mu)$。
\end{proof}

\begin{remark}[Besicovitch 微分定理（选读）]\label{proposition:4-1-3-4}
在 $\R^n$ 中，对任意局部有限正 Radon 测度 $\mu$，中心球基能够微分
$L^1_{\mathrm{loc}}(\mu)$ 中的函数。因此 \cref{proposition:4-1-3-3} 的相对密度结论对中心球成立。
\end{remark}

其证明需要 Besicovitch 覆盖定理：带指定中心的球族可分成至多 $N(n)$ 个不交子族，并仍覆盖中心集。
该覆盖结论产生相对于 $\mu$ 的极大弱型估计，再与连续逼近结合即可完成证明。这里不展开该覆盖定理。
Lebesgue 测度版的 $5r$ 论证不能直接用于一般 $\mu$，因为五倍球的 $\mu$-质量与原球质量之间通常
没有统一比值。本章后续结论不使用这条选读定理；它只记录一般相对微分所需的额外覆盖输入。

\begin{remark}[混合测度的局部密度]
取 $0\le p\le1$，令
\[
 \nu=p\delta_0+(1-p)\1_{(0,1)}m.
\]
它的总质量为 $1$。对 $m$-几乎每个 $x$，局部质量比只恢复
$(1-p)\1_{(0,1)}(x)$；原子质量 $p$ 在 $x=0$ 的比值中爆炸，却在 Lebesgue 几乎处处消失。
因此，总质量为一并不蕴含测度具有概率密度；后者要求该测度绝对连续于 $m$。
\end{remark}

局部质量比既可解释概率密度，也可用来研究数值密度估计；不过有限样本产生的是带带宽的近似比值，
不能把本节的几乎处处极限直接当作有限样本误差界。另一方面，一维有界变差函数的导数测度会同时含有
绝对连续和奇异部分，下一节将给出相应分解。

\begin{figure}[htbp]
\centering
\begin{tikzpicture}[node distance=13mm and 20mm]
  \node[book strong node] (nu) {$\nu$};
  \node[book node,below left=of nu] (ac) {铺在体积上的部分\\$fm$};
  \node[book warning node,below right=of nu] (sing) {零测骨架上的部分\\$\nu_s$};
  \node[book warm node,below=of ac] (ratio) {一般点的小球比值\\$\longrightarrow f(x)$};
  \node[book node,below=of sing] (zero) {$m$-a.e. 比值\\$\longrightarrow0$};
  \draw[book arrow] (nu)--(ac);
  \draw[book arrow] (nu)--(sing);
  \draw[book accent arrow] (ac)--(ratio);
  \draw[book dashed arrow] (sing)--(zero);
\end{tikzpicture}
\caption{Lebesgue 分解的局部图像。小球体积比恢复 $fm$ 的密度；奇异质量并未消失，
只是对 Lebesgue 几乎每个中心不可见。}
\label{fig:4-1-local-density-decomposition}
\end{figure}

\begin{exercise}
证明对 $x\ne0$，$D_r\delta_0(x)$ 最终为零；计算在 $x=0$ 的发散阶。
\end{exercise}
\begin{exercise}
把 \eqref{eq:4-1-measure-maximal} 的证明完整写出，并指出只在何处使用 $\sigma$ 的有限性。
\end{exercise}
\begin{exercise}
设 $S^{n-1}$ 上的表面测度作为 $\R^n$ 中的测度。证明它相对于 $n$ 维 Lebesgue 测度的密度
在 $m$-几乎处处为零，并说明球面上的中心点为何不与此矛盾。
\end{exercise}
\begin{exercise}
若 $\nu=fm$、$f\in L^1_{\mathrm{loc}}(\R^n)$，把密度比逐行化成球平均，并由 Lebesgue
微分定理推出 $D_r\nu(x)\to f(x)$ 对 $m$-几乎每个 $x$ 成立。
\end{exercise}
\begin{exercise}
给出一个非零奇异测度，使 $D_r\nu(x)\to0$ 对 $m$-几乎每点成立；分别用其总质量与局部比值解释
为何“密度几乎处处为零”不推出测度为零。
\end{exercise}

局部质量比已经把 Radon--Nikodym 密度变成可观察的极限，同时也暴露了奇异部分。实线上的奇异质量
常常可以编码进一个单调或有界变差函数；下一节将从分割上的总变差开始，把绝对连续、跳跃与奇异连续
三种行为逐一拆开。

\section{一维变差、Stieltjes 积分与绝对连续}\label{sec:04-02}

第~3 章在 Radon--Nikodym 定理的一维应用中已经引入有界变差与绝对连续函数。本节回顾这些定义，
并系统建立 Jordan 分解、Stieltjes 测度以及 Lebesgue 微积分基本定理之间的联系。

\subsection{单调函数、有界变差与 Jordan 分解}\label{subsec:4-2-1}
\index{有界变差}\index{Jordan 分解!函数版}

对 $C^1$ 曲线，弧长可以由速度积分计算。若曲线带有尖角，这个公式仍有希望；若它还带有跳跃，
普通导数却不能记录跳跃的大小。此外，函数可以在越来越短的区间上来回振荡，首尾增量很小，
沿途累计的路程却很大。因此应先问一个只涉及有限观测的问题：在任意有限分割上，把相邻函数值的
绝对增量相加，所得数能否由同一个常数控制？

有界变差的概念同时源于曲线长度与 Fourier 级数。分割定义既描述有限折线对总路程的逼近，也自然
导出增量测度，从而把函数的 Jordan 分解与测度的 Jordan 分解联系起来。

\begin{definition}[总变差]\label{def:4-2-1-1}
设 $f:[a,b]\to\R$。对分割
$P=\{a=x_0<x_1<\cdots <x_m=b\}$，记
\[
 V(f;P)=\sum_{j=1}^m|f(x_j)-f(x_{j-1})|,
 \qquad
 V_a^b(f)=\sup_P V(f;P).
\]
若 $V_a^b(f)<\infty$，称 $f$ 在 $[a,b]$ 上有界变差，记作
$f\in BV([a,b])$。同样的记号 $V_s^t(f)$ 用于任意子区间 $[s,t]\subset[a,b]$。
\end{definition}

把一点插入分割只会增大变差和。由此，对 $a\le s\le t\le u\le b$，
\begin{equation}\label{eq:4-2-variation-additivity}
 V_s^u(f)=V_s^t(f)+V_t^u(f).
\end{equation}
事实上，合并 $[s,t]$ 与 $[t,u]$ 上的两个分割给出“$\ge$”；反过来，把 $t$ 插进
$[s,u]$ 的任意分割，再按 $t$ 左右分组，便得到“$\le$”。这项简单的可加性是后面把变差
变成测度的第一个信号。
为了明确上确界的取法，给定 $\varepsilon>0$，取两个分割 $P_1,P_2$ 使
\[
 V(f;P_1)>V_s^t(f)-\frac\varepsilon2,\qquad
 V(f;P_2)>V_t^u(f)-\frac\varepsilon2.
\]
将它们在 $t$ 处合并，得
\[
 V_s^u(f)\ge V(f;P_1)+V(f;P_2)
 >V_s^t(f)+V_t^u(f)-\varepsilon.
\]
令 $\varepsilon\downarrow0$ 得到“$\ge$”。反向上，对 $[s,u]$ 的任意分割 $P$，令
$P'=P\cup\{t\}$，则
\[
 V(f;P)\le V(f;P')
 =V(f;P'\cap[s,t])+V(f;P'\cap[t,u])
 \le V_s^t(f)+V_t^u(f).
\]
对 $P$ 取上确界即得“$\le$”。

\begin{definition}[变差函数]\label{def:4-2-1-2}
若 $f\in BV([a,b])$，定义
\[
 v_f(x)=V_a^x(f),\qquad a\le x\le b.
\]
\end{definition}

由 \eqref{eq:4-2-variation-additivity}，$v_f$ 递增；并且对 $x<y$，
\begin{equation}\label{eq:4-2-increment-dominated}
 |f(y)-f(x)|\le V_x^y(f)=v_f(y)-v_f(x).
\end{equation}
所以 $v_f$ 记录的不是净位移，而是已经走过的总路程。

\begin{example}[阶梯函数]\label{ex:4-2-1-1}
设 $a<x_1<\cdots <x_N\le b$，并令
\[
 F(x)=\sum_{k=1}^N c_k\1_{[x_k,b]}(x).
\]
在不含跳点的区间上 $F$ 恒定，而经过 $x_k$ 时增量为 $c_k$。把所有跳点加入任意分割后，
任意分割 $P={t_j}$ 的第 $j$ 个增量可写成
\[
 F(t_j)-F(t_{j-1})
 =\sum_{\{k:t_{j-1}<x_k\le t_j\}}c_k.
\]
这些指标集两两不交，所以
\[
 V(F;P)
 \le\sum_j\sum_{\{k:t_{j-1}<x_k\le t_j\}}|c_k|
 =\sum_{k=1}^N|c_k|.
\]
取包含全部 $x_k$ 的分割时，每个非零增量正好是一个 $c_k$，上界达到。故
\[
 V_a^b(F)=\sum_{k=1}^N|c_k|.
\]
与 $F$ 对应的增量测度是 $\sum_kc_k\delta_{x_k}$。如果把一个跳点放在 $a$，则所有
$(s,t]$（$a\le s<t$）都不包含该跳跃；因此左端原子必须另行指定。
\end{example}

\begin{proposition}[单调函数的正则性]\label{proposition:4-2-1-1}
单调函数在每个内点都有有限的左、右极限；它的不连续点都是跳跃点，并且至多可数。
特别地，单调函数 Borel 可测。
\end{proposition}

\begin{proof}
只证非减函数 $g$；对非增函数考察 $-g$。若 $a<x<b$，则
\[
 g(x-):=\sup_{a\le t<x}g(t),\qquad
 g(x+):=\inf_{x<t\le b}g(t)
\]
都是有限实数，并且单调性直接给
$g(t)\to g(x-)$（$t\uparrow x$）与 $g(t)\to g(x+)$（$t\downarrow x$）。
因此若 $g$ 在 $x$ 不连续，就有 $g(x-)<g(x)$ 或 $g(x)<g(x+)$；其函数值中出现一个
正长度的空隙。

对每个不连续点 $x$，从开区间 $(g(x-),g(x+))$ 中选一个有理数；若 $x<y$，单调性给
$g(x+)\le g(y-)$，所以不同点对应的开区间两两不交。每个这样的区间含有不同的有理数，
故内点中的不连续点集可数；两个端点至多再贡献两个单侧跳跃，因而全部不连续点仍至多可数。
最后，对任意 $c\in\R$，集合
$\{x:g(x)>c\}$ 是空集、全区间或一个端点取舍未定的区间，因而是 Borel 集。
这些半直线生成 $\Borel(\R)$，所以 $g$ Borel 可测。
\end{proof}

\begin{theorem}[Jordan 分解（函数版）]\label{theorem:4-2-1-2}
函数 $f:[a,b]\to\R$ 属于 $BV([a,b])$，当且仅当它可以写成两个非减函数之差。
若 $f\in BV([a,b])$，则可取
\[
 p=\frac{v_f+f}{2},\qquad q=\frac{v_f-f}{2},\qquad f=p-q.
\]
若再有 $f(a)=0$，则 $p(a)=q(a)=0$；在所有满足
$f=\widetilde p-\widetilde q$ 且
$\widetilde p(a)=\widetilde q(a)=0$ 的非减分解中，这一对逐点最小：
$\widetilde p\ge p$、$\widetilde q\ge q$。
\end{theorem}

\begin{proof}
设 $f\in BV([a,b])$。由 \eqref{eq:4-2-increment-dominated}，对 $x<y$，
\[
 (v_f\pm f)(y)-(v_f\pm f)(x)
 =v_f(y)-v_f(x)\pm(f(y)-f(x))\ge0.
\]
故 $p,q$ 非减，且 $p-q=f$。

反过来，若 $f=p-q$，其中 $p,q$ 非减，则对每个分割 $P$，
\[
 V(f;P)
 \le\sum_j\bigl(p(x_j)-p(x_{j-1})\bigr)
   +\sum_j\bigl(q(x_j)-q(x_{j-1})\bigr)
 =p(b)-p(a)+q(b)-q(a).
\]
所以 $f\in BV([a,b])$。

最后设 $f(a)=0$，并给定规范化分解 $f=\widetilde p-\widetilde q$。对 $[a,x]$ 上任意分割，
上面的估计给
\[
 V(f;P)\le \widetilde p(x)+\widetilde q(x).
\]
取上确界得 $v_f(x)\le\widetilde p(x)+\widetilde q(x)$。与
$f(x)=\widetilde p(x)-\widetilde q(x)$ 相加、相减，分别得到
\[
 p(x)=\frac{v_f(x)+f(x)}2\le\widetilde p(x),\qquad
 q(x)=\frac{v_f(x)-f(x)}2\le\widetilde q(x).
\]
这正是所述最小性。
\end{proof}

上一证明还给出 BV 函数在每点的左右极限：它是两个单调函数之差。特别地，BV 函数的跳跃点
至多可数。不过，点态 BV 不是按几乎处处相等取商的空间。若 $a<x_0<b$ 且
$f=\1_{\{x_0\}}$，则 $V_a^b(f)=2$；把 $x_0$ 处的值改为 $0$ 后，函数变成零函数，变差也变成
$0$。因此“选一个右连续版本”是额外操作，不是像 $L^p$ 中那样免费的代表元替换。

\begin{proposition}[BV 的代数性质]\label{proposition:4-2-1-3}
$BV([a,b])$ 是向量空间，并且
\[
 V_a^b(\alpha f)=|\alpha|V_a^b(f),\qquad
 V_a^b(f+g)\le V_a^b(f)+V_a^b(g).
\]
\end{proposition}

\begin{proof}
两个等式分别来自
$|\alpha\Delta f|=|\alpha|\,|\Delta f|$ 与
$|\Delta(f+g)|\le|\Delta f|+|\Delta g|$：先对每个分割求和，再取上确界。
\end{proof}

Helly 选择定理断言：统一有界且总变差统一有界的函数列具有点态收敛子列。其证明从两个单调
Jordan 分量出发，先在可数稠密集上作对角抽取，再在连续点利用单调性夹逼，并处理可数多个跳点。
该定理将在讨论弱收敛时证明。

要让区间增量唯一编码测度，我们固定如下约定。

\begin{definition}[右连续规范与左端原子]\label{def:4-2-1-3}
与 Stieltjes 测度对应的函数总取右连续的 BV 代表，并规定 $F(a)=0$；其增量只编码
$(a,b]$ 上的质量。若要编码 $[a,b]$ 上的任意测度，必须另给数
$\nu(\{a\})$，或者在区间外规定 $F(a-)=0$ 并允许 $F(a)\ne0$。这两种端点约定不可混用。
\end{definition}

只知道所有 $F(t)-F(s)$ 不能恢复左端原子：$0$ 与 $c\delta_a$ 在每个
$(s,t]$（$a\le s<t\le b$）上都给出零质量。相同的数据也不能区分 $F$ 与 $F+C$。

\begin{lemma}[右连续性传给变差函数]\label{lemma:4-2-1-right-variation}
若 $F\in BV([a,b])$、$F(a)=0$ 且 $F$ 右连续，则 $v_F$ 以及最小 Jordan 分量
$p=(v_F+F)/2$、$q=(v_F-F)/2$ 都右连续。
\end{lemma}

\begin{proof}
$p,q$ 非减，所以每点有右极限。固定 $x<b$。由 $F=p-q$ 的右连续性，
\[
 p(x+)-p(x)=q(x+)-q(x)=:d\ge0.
\]
若 $d>0$，令 $r=d\1_{(x,b]}$。函数 $p-r$ 与 $q-r$ 仍非减：越过 $x$ 时恰好删去共同的
右跳跃，其他增量不变。它们在 $a$ 仍为零，并且
$F=(p-r)-(q-r)$，却在 $x$ 右侧严格小于最小分量 $p,q$，与
\cref{theorem:4-2-1-2} 的最小性矛盾。因此 $d=0$。这对每个 $x<b$ 成立，故 $p,q$ 右连续，
$v_F=p+q$ 也右连续。
\end{proof}

\begin{theorem}[BV 函数与有限有符号 Stieltjes 测度]\label{theorem:4-2-1-4}
右连续、满足 $F(a)=0$ 的 $BV([a,b])$ 函数，与 $(a,b]$ 上有限有符号 Borel 测度一一对应；
对应由
\[
 \nu_F((s,t])=F(t)-F(s),\qquad a\le s<t\le b,
\]
刻画，并且
\begin{equation}\label{eq:4-2-variation-measure}
 |\nu_F|((a,b])=V_a^b(F).
\end{equation}
反向对应为 $F_\nu(x)=\nu((a,x])$。
\end{theorem}

\begin{proof}
先设 $F$ 如定理所述，并取最小 Jordan 分解 $F=p-q$。由上一引理，$p,q$ 右连续、非减且
在 $a$ 为零。\Cref{theorem:2-3-3-2,cor:stieltjes-compact-interval} 分别给出
$(a,b]$ 上有限正测度 $\mu_p,\mu_q$，使
\[
 \mu_p((s,t])=p(t)-p(s),\qquad
 \mu_q((s,t])=q(t)-q(s).
\]
令 $\nu_F=\mu_p-\mu_q$，便得到所需增量。半开区间构成生成 Borel $\sigma$-代数的
$\pi$-系统。若两个有限有符号测度有相同增量，写成 $\nu_1^+-\nu_1^-$ 与
$\nu_2^+-\nu_2^-$ 后，两个有限正测度
$\nu_1^++\nu_2^-$、$\nu_2^++\nu_1^-$ 在该 $\pi$-系统上相等；有限正测度唯一性说明它们
处处相等，故 $\nu_1=\nu_2$。这证明 $\nu_F$ 唯一。

先证 \eqref{eq:4-2-variation-measure} 的一个方向。任意分割 $P$ 给出
\[
 \sum_j|F(x_j)-F(x_{j-1})|
 =\sum_j|\nu_F((x_{j-1},x_j])|
 \le\sum_j|\nu_F|((x_{j-1},x_j])
 =|\nu_F|((a,b]),
\]
故 $V_a^b(F)\le|\nu_F|((a,b])$。

反向不等式把有符号测度的正负区域化为有限分割上的估计。取 $\nu_F$ 的 Hahn 分解
$(P_+,P_-)$。给定 $\varepsilon>0$，有限 Borel 测度 $|\nu_F|$ 在紧区间上正则，故可取互不相交
紧集 $K_+\subset P_+$、$K_-\subset P_-$，使
\[
 |\nu_F|\bigl((a,b]\setminus(K_+\cup K_-)\bigr)<\varepsilon.
\]
再取互不相交的相对开集 $U_+\supset K_+$、$U_-\supset K_-$，并使
$|\nu_F|(U_\pm\setminus K_\pm)<\varepsilon$。由紧性，可在各 $U_\pm$ 内用有限个半开区间覆盖
$K_\pm$；细分全部端点后，得到两族互不相交半开区间，其并记作 $A_+,A_-$。令
$h=\1_{A_+}-\1_{A_-}$。在 $K_+\cup K_-$ 上，$h$ 等于
$\dd\nu_F/\dd|\nu_F|$ 的符号，而所有误差落在
$(a,b]\setminus(K_+\cup K_-)$ 与 $U_\pm\setminus K_\pm$ 中。因此
\[
 \begin{aligned}
 \int h\,\dd\nu_F
 &=\int h\frac{\dd\nu_F}{\dd|\nu_F|}\,\dd|\nu_F|\\
 &\ge |\nu_F|(K_+\cup K_-)
      -|\nu_F|\bigl((A_+\cup A_-)\setminus(K_+\cup K_-)\bigr)\\
 &\ge |\nu_F|((a,b])
      -|\nu_F|\bigl((a,b]\setminus(K_+\cup K_-)\bigr)
      -|\nu_F|(U_+\setminus K_+)-|\nu_F|(U_-\setminus K_-)\\
 &>|\nu_F|((a,b])-3\varepsilon.
 \end{aligned}
\]
也就是
\[
 |\nu_F|((a,b])-3\varepsilon
 \le \int h\,\dd\nu_F.
\]
把 $A_+,A_-$ 的全部端点并入同一分割，并把各区间的符号记为 $\epsilon_j\in\{-1,0,1\}$，则
\[
 \int h\,\dd\nu_F
 =\sum_j\epsilon_j\bigl(F(x_j)-F(x_{j-1})\bigr)
 \le\sum_j|F(x_j)-F(x_{j-1})|
 \le V_a^b(F).
\]
令 $\varepsilon\downarrow0$，得到 $|\nu_F|((a,b])\le V_a^b(F)$。

反过来，给定有限有符号测度 $\nu$，令 $F(x)=\nu((a,x])$。若 $x_n\downarrow x$，则
$(a,x_n]\downarrow(a,x]$，有限正负部分的从上连续性给 $F(x_n)\to F(x)$，故 $F$ 右连续。
对任意分割，
\[
 \sum_j|F(x_j)-F(x_{j-1})|
 \le |\nu|((a,b]),
\]
所以 $F\in BV$ 且 $F(a)=0$。区间增量等于 $\nu((s,t])$；唯一性给
$\nu_F=\nu$，而已证的等式给 $V(F)=|\nu|((a,b])$。两项构造因此互逆。
\end{proof}

\begin{example}[不定积分的变差]\label{ex:4-2-1-2}
设 $g\in L^1([a,b])$，$F(x)=\int_a^xg(t)\,\dd t$。对应增量测度为
$g\,m$，而 $|g\,m|=|g|\,m$。因此 \eqref{eq:4-2-variation-measure} 给出
\[
 V_a^b(F)=\int_a^b|g(t)|\,\dd t.
\]
这也说明只写 $V(F)\le\int|g|$ 会漏掉一半信息。
\end{example}

\begin{example}[Cantor 函数]\label{ex:4-2-1-3}
令 $K$ 为三分 Cantor 集。若
$x\in K$ 有三进制展开 $x=0.\varepsilon_1\varepsilon_2\cdots{}_3$，其中
$\varepsilon_j\in\{0,2\}$，置
\[
 C(x)=0.(\varepsilon_1/2)(\varepsilon_2/2)\cdots{}_2,
\]
并在每个被删去的三分开区间上作常值延拓。端点的两种展开给同一连续延拓。于是 $C$ 连续非减，
$C(0)=0,C(1)=1$，故 $V_0^1(C)=1$。它在 $[0,1]\setminus K$ 上局部常值，而 $K$ 的
Lebesgue 测度为零，所以 $C'=0$ a.e.；然而总变差仍为 $1$。普通 a.e. 导数没有记录这份奇异连续增长。
\end{example}

\begin{figure}[htbp]
\centering
\begin{tikzpicture}[x=1.0cm,y=.75cm]
  \begin{scope}
    \draw[->] (0,0)--(5.4,0) node[right]{$x$};
    \draw[->] (0,-1.25)--(0,1.55) node[above]{$f$};
    \draw[BookAccent,thick] plot[smooth] coordinates
      {(0,.1) (.7,.85) (1.4,-.55) (2.1,1.05) (3,.25) (3.8,1.2) (4.8,.6)};
    \node[book label] at (2.5,-1.05){有升有降的 $f$};
  \end{scope}
  \begin{scope}[xshift=6.3cm]
    \draw[->] (0,0)--(5.4,0) node[right]{$x$};
    \draw[->] (0,0)--(0,3.2) node[above]{$v_f$};
    \draw[BookWarm,thick] plot[smooth] coordinates
      {(0,.1) (.7,.65) (1.4,1.25) (2.1,1.9) (3,2.25) (3.8,2.75) (4.8,3)};
    \draw[BookAccent,dashed,thick] plot[smooth] coordinates
      {(0,.05) (.7,.55) (1.4,.62) (2.1,1.45) (3,1.5) (3.8,2.1) (4.8,2.25)};
    \draw[BookWarning,dashed,thick] plot[smooth] coordinates
      {(0,.02) (.7,.12) (1.4,.68) (2.1,.72) (3,1.05) (3.8,1.08) (4.8,1.35)};
    \node[book label] at (2.6,.32){$v_f$（实线），$p,q$（虚线）};
  \end{scope}
  \draw[book accent arrow] (5.1,1.15)--node[above]{$v_f(x)=V_a^x(f)$}(6.1,1.15);
\end{tikzpicture}
\caption{变差函数把正负振荡都记成增长；$p=(v_f+f)/2$ 与 $q=(v_f-f)/2$ 再把增长分到两个单调分量。}
\label{fig:variation-jordan}
\end{figure}

\begin{remark}
BV 不推出连续，更不推出绝对连续；Cantor 函数又说明连续 BV 仍可能不能由 a.e. 导数恢复。
测度对应保留跳跃、绝对连续和奇异连续三类增量，后两小节将依次把它们用于积分与微积分基本定理。
\end{remark}

\begin{exercise}
证明 $V_s^u(f)=V_s^t(f)+V_t^u(f)$，并据此计算任意有限阶梯函数的总变差。
\end{exercise}
\begin{exercise}
利用 Jordan 分解证明每个 BV 函数都有左右极限，并证明其不连续点至多可数。
\end{exercise}
\begin{exercise}
若 $f(a)=0$ 且 $f=p-q$ 是最小 Jordan 分解，证明
$V_a^b(f)=V_a^b(p)+V_a^b(q)=p(b)+q(b)$。说明为何最后一个等号必须先规范端点。
\end{exercise}
区间增量现在已经成为测度。下一步是检验经典分割和是否收敛，以及它与测度积分是否真的给出同一个数。

\subsection{Stieltjes 测度、分割和与积分}\label{subsec:4-2-2}
\index{Stieltjes 积分}\index{标记分割网}

普通 Riemann 和用区间长度作权重。若权重来自一个递增函数 $F$，第 $j$ 段更自然的“宽度”是
$F(x_j)-F(x_{j-1})$。这允许某个几何长度趋于零的区间仍携带固定权重---那正是 $F$ 的跳跃。
然而，若被积函数也在同一点跳跃，不同标记可能产生不同的极限。下面先计算原子、绝对连续密度与
Cantor 型奇异测度三种典型情形，再给出分割和稳定收敛的条件。

Stieltjes 积分的早期问题来自连分数与矩问题；现代测度论把分割增量扩张为 Borel 测度。标记分割网
仍保留有限和，并精确区分沿某一列分割收敛与对所有充分细分割一致稳定这两种要求。

\begin{example}[跳跃积分]\label{ex:4-2-2-1}
设 $a<x_k\le b$，$\sum_{k\ge1}|a_k|<\infty$，并令
\[
 F(x)=\sum_{k\ge1}a_k\1_{[x_k,b]}(x).
\]
级数一致绝对收敛，故 $F$ 右连续；对每个半开区间直接相减可得
\[
 \nu_F=\sum_{k\ge1}a_k\delta_{x_k}.
\]
因而对连续 $f$，
\begin{equation}\label{eq:4-2-jump-integral}
 \int_{(a,b]}f\,\dd F=\sum_{k\ge1}a_kf(x_k).
\end{equation}
右侧绝对收敛，因为
$\sum_k|a_kf(x_k)|\le\|f\|_\infty\sum_k|a_k|$。有限阶梯情形相应地给出有限加权和
$\sum_k a_kf(x_k)$。
\end{example}

\begin{example}[光滑密度]\label{ex:4-2-2-2}
若 $g\in L^1([a,b])$、$g\ge0$，且
$F(x)=\int_a^xg(t)\,\dd t$，则对半开区间
\[
 \mu_F((s,t])=\int_s^tg(x)\,\dd x.
\]
测度唯一性给 $\dd\mu_F=g\,\dd x$，所以每个 $fg\in L^1$ 都满足
\[
 \int_{(a,b]}f\,\dd F=\int_a^bf(x)g(x)\,\dd x.
\]
\end{example}

\begin{example}[Cantor 柱状近似]\label{ex:4-2-2-3}
令 $K_n$ 是 Cantor 构造第 $n$ 步剩下的 $2^n$ 个闭区间之并，并置
\[
 \dd\mu_n=\left(\frac32\right)^n\1_{K_n}\,\dd x.
\]
每个长度 $3^{-n}$ 的柱子质量为 $2^{-n}$，故 $\mu_n([0,1])=1$。令 $\mu_C=\dd C$ 为
Cantor 函数的 Stieltjes 测度；它在每个第 $n$ 级柱子上也有质量 $2^{-n}$。若这些柱子的左端点为
$\ell_1,\dots,\ell_{2^n}$，连续函数 $f$ 的模连续性记作 $\omega_f$，则
\[
 \left|\int f\,\dd\mu_n-2^{-n}\sum_{j=1}^{2^n}f(\ell_j)\right|
 \le\omega_f(3^{-n}),
\]
而把 $\mu_n$ 换成 $\mu_C$ 后同一估计仍成立。因此
\begin{equation}\label{eq:4-2-cantor-column-limit}
 \left|\int f\,\dd\mu_n-\int f\,\dd C\right|
 \le2\omega_f(3^{-n})\longrightarrow0.
\end{equation}
这些绝对连续柱状密度越来越高、支撑总长度 $(2/3)^n$ 越来越小，极限却不是原子测度，而是集中在
Cantor 零测集上的奇异连续测度。公式 \eqref{eq:4-2-cantor-column-limit} 给出了相应积分的定量收敛。
\end{example}

上面三例中的积分暂时只是第 3 章已经定义的测度积分，增量测度来自
\cref{theorem:4-2-1-4}。现在把两种 Stieltjes 语言正式并列起来。

\begin{definition}[Stieltjes 测度]\label{def:4-2-2-1}
若 $F:[a,b]\to\R$ 右连续、非减且 $F(a)=0$，记 $\mu_F$ 为 $(a,b]$ 上满足
\[
 \mu_F((s,t])=F(t)-F(s)
\]
的有限正测度。若 $F$ 是右连续、规范化的 BV 函数，则以 $\nu_F$ 表示
\cref{theorem:4-2-1-4} 给出的有限有符号测度。若积分域写成 $[a,b]$，左端原子
$\nu(\{a\})$ 必须另外加入。
\end{definition}

\begin{definition}[标记分割网]\label{def:4-2-2-2}
标记分割 $\mathcal P=(P,\xi)$ 由
\[
 P=\{a=x_0<\cdots <x_m=b\},\qquad
 \xi_j\in[x_{j-1},x_j]
\]
组成。若 $Q$ 细化 $P$，规定 $(P,\xi)\preceq(Q,\eta)$；不要求新旧标记相容。
任意两个底层分割的端点之并是共同细化，所以这些指标构成有向预序。定义
\[
 S_F(f;P,\xi)
 =\sum_{j=1}^mf(\xi_j)\bigl(F(x_j)-F(x_{j-1})\bigr).
\]
若这个网在值域中收敛，称其极限为 Riemann--Stieltjes 积分。
\end{definition}

\begin{definition}[Lebesgue--Stieltjes 积分]\label{def:4-2-2-3}
对 $\mu_F$ 可积的标量函数或 Bochner 可积的 Banach 值函数 $f$，定义
\[
 \int_{(a,b]}f\,\dd F:=\int_{(a,b]}f\,\dd\mu_F.
\]
BV 情形令 $u=\dd\nu_F/\dd|\nu_F|$。若 Banach 值函数 $f$ 强可测且
\[
 \int_{(a,b]}\|f\|\,\dd|\nu_F|<\infty,
\]
定义
\[
 \int_{(a,b]}f\,\dd F
 :=\int_{(a,b]}f\,\dd\nu_F
 :=\int_{(a,b]}u(x)f(x)\,\dd|\nu_F|(x),
\]
右端是关于正测度 $|\nu_F|$ 的 Bochner 积分。等价地，也可对 $\nu_F$ 的 Jordan 正负部分
分别积分后作差。于是
\[
 \left\|\int f\,\dd\nu_F\right\|
 \leq\int\|f\|\,\dd|\nu_F|.
\]
连续 Banach 值函数在紧区间上有界，而 $|\nu_F|$ 有限，所以自动满足上述可积条件。
\end{definition}

不要求细化分割继承原标记是必要的。若只沿某一列特定标记取极限，极限可能依赖标记选择。
网的 Cauchy 性要求所有足够细的分割及其任意标记给出相近的数值。

\begin{theorem}[分割和收敛]\label{theorem:4-2-2-1}
设 $F$ 右连续、非减且 $F(a)=0$，$B$ 为 Banach 空间，$f\in C([a,b];B)$。则
$S_F(f;P,\xi)$ 构成 $B$ 中的 Cauchy 网，因而存在极限，并且
\[
 \left\|\int_a^bf\,\dd F\right\|
 \le\|f\|_\infty\bigl(F(b)-F(a)\bigr).
\]
\end{theorem}

\begin{proof}
记
\[
 \omega_f(\delta)=\sup\{\|f(x)-f(y)\|:|x-y|\le\delta\}.
\]
紧区间上的一致连续性给 $\omega_f(\delta)\to0$。给定 $\varepsilon>0$；若 $F(b)=F(a)$，
所有和均为零。否则选 $\delta>0$ 使
$2\omega_f(\delta)(F(b)-F(a))<\varepsilon$，再取网格小于 $\delta$ 的底分割
$P_0=\{x_j\}$，并在各段固定标记 $\zeta_j$。

设 $(Q,\eta)$ 是 $P_0$ 的任意细化。把 $Q$ 的小区间按其所在的
$[x_{j-1},x_j]$ 分组。在第 $j$ 组中，每个新标记与 $\zeta_j$ 的距离不超过 $\delta$，而所有
$F$ 增量之和恰为 $F(x_j)-F(x_{j-1})$。所以
\[
 \left\|S_F(f;Q,\eta)-S_F(f;P_0,\zeta)\right\|
 \le\omega_f(\delta)\sum_j\bigl(F(x_j)-F(x_{j-1})\bigr)
 =\omega_f(\delta)(F(b)-F(a)).
\]
两个任意细化分别与固定和比较，距离小于 $\varepsilon$。故该网 Cauchy；$B$ 的完备性给出极限。
最后每个和满足
\[
 \|S_F(f;P,\xi)\|
 \le\|f\|_\infty\sum_j(F(x_j)-F(x_{j-1}))
 =\|f\|_\infty(F(b)-F(a)),
\]
令网取极限即得估计。
\end{proof}

\begin{theorem}[Riemann--Stieltjes 与 Lebesgue--Stieltjes 一致]\label{theorem:4-2-2-2}
在 \cref{theorem:4-2-2-1} 的假设下，标记分割网的极限等于 Bochner 积分
\[
 \int_{(a,b]} f\,\dd\mu_F.
\]
特别地，标量连续函数的两种 Stieltjes 积分一致。
\end{theorem}

\begin{proof}
给标记分割 $(P,\xi)$ 配上阶梯函数
\[
 s_{P,\xi}(x)=\sum_{j=1}^mf(\xi_j)\1_{(x_{j-1},x_j]}(x).
\]
有限简单函数的 Bochner 积分定义与 Stieltjes 增量给
\[
 \int_{(a,b]}s_{P,\xi}\,\dd\mu_F=S_F(f;P,\xi).
\]
若 $P$ 的网格小于 $\delta$，则
$\|s_{P,\xi}(x)-f(x)\|\le\omega_f(\delta)$ 对所有 $x\in(a,b]$ 成立。由 Bochner 积分基本估计
\cref{theorem:3-4-2-2}，
\[
 \left\|S_F(f;P,\xi)-\int_{(a,b]}f\,\dd\mu_F\right\|
 \le\omega_f(\delta)\mu_F((a,b])
 =\omega_f(\delta)(F(b)-F(a)).
\]
右端与标记无关并趋于零，故网极限就是测度积分。
\end{proof}

\begin{proposition}[BV Stieltjes 积分估计]\label{proposition:4-2-2-3}
若 $F$ 右连续、$F(a)=0$ 且 $F\in BV([a,b])$，则对连续 Banach 值 $f$，
\[
 \left\|\int_{(a,b]}f\,\dd F\right\|
 \le\|f\|_\infty V_a^b(F).
\]
在实标量情形，$L_F:C([a,b])\to\R$，
$L_F(f)=\int f\,\dd F$ 的算子范数恰为 $V_a^b(F)$。
\end{proposition}

\begin{proof}
由有限有符号测度积分估计和 \eqref{eq:4-2-variation-measure}，
\[
 \left\|\int f\,\dd\nu_F\right\|
 \le\int\|f\|\,\dd|\nu_F|
 \le\|f\|_\infty V_a^b(F).
\]
所以 $\|L_F\|\le V_a^b(F)$。

为证反向，取 $\nu_F$ 的 Hahn 分解 $(P_+,P_-)$。给定 $\varepsilon>0$，由正则性取互不相交紧集
$K_+\subset P_+$、$K_-\subset P_-$，使
$|\nu_F|((a,b]\setminus(K_+\cup K_-))<\varepsilon$。若两者都非空，令
\[
 \varphi(x)=
 \frac{\dist(x,K_-)-\dist(x,K_+)}
      {\dist(x,K_-)+\dist(x,K_+)};
\]
它连续、$|\varphi|\le1$，并在 $K_+$ 上等于 $1$、在 $K_-$ 上等于 $-1$。若某个紧集为空，
取相应常数符号即可。于是
\[
 L_F(\varphi)=\int\varphi\,\dd\nu_F
 \ge |\nu_F|(K_+\cup K_-)-\varepsilon
 \ge V_a^b(F)-2\varepsilon.
\]
令 $\varepsilon\downarrow0$，得到 $\|L_F\|\ge V_a^b(F)$。
\end{proof}

如果 $F$ 是一般 BV 函数，分割和可由其最小 Jordan 分量之差处理；误差估计中的总质量相应换成
$V(F)$。被积函数的连续性仍是分割网收敛的重要条件。

\begin{remark}[共同跳跃使标记网失败]
在 $[-1,1]$ 上令 $F=H=\1_{[0,1]}$。则 $\dd F=\delta_0$，按右连续代表，
Lebesgue--Stieltjes 积分给
$\int H\,\dd F=H(0)=1$。但每个足够细的分割都有一个小区间跨过 $0$；若该段标记取在
$0$ 左侧，相应分割和的原子项为 $0$，取在 $0$ 或右侧则为 $1$。这两族任意细指标始终相差 $1$，
所以标记网不收敛。现代测度积分仍有定义，经典分割和却要求更强的相容性。
\end{remark}

BV 函数的部分积分公式必须区分当前值与左极限。下面给出公式及其有限阶梯模型；一般 BV 情形需要
进一步的近似论证。

\begin{remark}[BV 部分积分公式的端点形式]\label{proposition:4-2-2-4}
若 $F,G$ 右连续且有界变差，预期的公式是
\begin{equation}\label{eq:4-2-bv-ibp-preview}
 F(b)G(b)-F(a)G(a)
 =\int_{(a,b]}F(x-)\,\dd G(x)
  +\int_{(a,b]}G(x)\,\dd F(x).
\end{equation}
对共同分割上的阶梯函数，这不是极限论证，而是逐段恒等式
\[
 F(x_j)G(x_j)-F(x_{j-1})G(x_{j-1})
 =F(x_{j-1})\Delta_jG+G(x_j)\Delta_jF
\]
的望远镜求和。若二者在 $x$ 共同跳跃，把两个积分都写成当前值会多算
$\Delta F(x)\Delta G(x)$；把二者都写成左极限则少算同一项。因此，公式
\eqref{eq:4-2-bv-ibp-preview} 必须采用一项左极限与一项当前值。下一小节将对绝对连续
函数证明不涉及跳跃项的部分积分公式。
\end{remark}

\begin{figure}[htbp]
\centering
\begin{tikzpicture}[x=1.05cm,y=.72cm]
  \draw[->] (0,0)--(10.2,0) node[right]{$x$};
  \foreach \x in {0,1.1,2.1,3.4,4.5,5.7,6.8,8.1,9.4}
    \draw[BookInk!55] (\x,-.12)--(\x,.12);
  \draw[BookAccent,thick]
    (0,.3)--(1.1,.3)--(1.1,.65)--(2.1,.65)--(2.1,.9)--(3.4,.9)--
    (3.4,2.6)--(4.5,2.6)--(4.5,2.85)--(5.7,2.85)--(5.7,3.25)--
    (6.8,3.25)--(6.8,3.5)--(8.1,3.5)--(8.1,4.0)--(9.4,4.0);
  \draw[BookWarning,very thick] (3.4,.9)--(3.4,2.6);
  \node[book warning node] at (3.4,3.3){跳跃给原子柱};
  \draw[decorate,decoration={brace,mirror,amplitude=4pt},BookWarm]
    (5.7,-.35)--(6.8,-.35) node[midway,below=5pt]{$\Delta F$ 是权重};
  \node[book label] at (7.5,1.35){细化 $x$--分割};
\end{tikzpicture}
\caption{Stieltjes 和的权重由 $F$ 的纵向增量决定；几何上很窄的区间也可能承载一个固定原子。}
\label{fig:stieltjes-tagged-net}
\end{figure}

离散谱产生 \eqref{eq:4-2-jump-integral} 型的加权和，矩问题考察测试函数
$1,x,x^2,\ldots$ 的 Stieltjes 积分，紧区间上的测度表示则反向使用
\cref{proposition:4-2-2-3} 的线性泛函估计。这些应用将在后文展开。

\begin{exercise}
不引用定理结论，按任意共同细化逐项证明标记分割和构成 Cauchy 网。
\end{exercise}
\begin{exercise}
对 $F=2\1_{[1/3,1]}-\1_{[2/3,1]}$ 计算 $V_0^1(F)$，并求
$\int_0^1x^2\,\dd F$。
\end{exercise}
\begin{exercise}
证明 \cref{proposition:4-2-2-3} 中的范数等式；特别说明紧集分离函数为何连续且模不超过 $1$。
\end{exercise}
\begin{exercise}
若改用左连续累计函数与区间 $[s,t)$，核对原子落点和部分积分公式中的左右值。构造一个例子说明
只改连续性约定、不改半开区间方向会产生错误。
\end{exercise}

Stieltjes 测度同时容纳原子、普通密度和 Cantor 质量。下一步要回答更具体的问题：哪些累计函数
完全由普通 $L^1$ 密度产生？

\subsection{绝对连续函数与 Lebesgue 基本定理}\label{subsec:4-2-3}
\index{绝对连续}\index{Lebesgue 基本定理}\index{Luzin N 性质}

Cantor 函数连续、单调且几乎处处导数为零，却从 $0$ 增长到 $1$。因此，连续性与几乎处处可微性
都不足以保证由导数恢复函数。绝对连续性所排除的正是如下行为：在总长度任意小的不交区间族上，
函数增量的总和仍保持为一个固定的正数。

\begin{definition}[绝对连续]\label{def:4-2-3-1}
函数 $F:[a,b]\to\R$ 称为绝对连续，如果对每个 $\varepsilon>0$，存在 $\delta>0$，使得任意有限族
两两不交的开区间 $(\alpha_k,\beta_k)\subset[a,b]$ 只要满足
\[
 \sum_k(\beta_k-\alpha_k)<\delta,
\]
就有
\[
 \sum_k|F(\beta_k)-F(\alpha_k)|<\varepsilon.
\]
记所有这类函数为 $AC([a,b])$。
\end{definition}

只取一个区间便知绝对连续函数一致连续。更重要的是，绝对连续性不仅控制各区间的首尾差，还控制
其中所有振荡。

\begin{lemma}[短区间族的局部变差控制]\label{lemma:4-2-3-local-variation}
若 $F\in AC([a,b])$，则对每个 $\varepsilon>0$，存在 $\delta>0$，使任意有限族两两不交区间
$(\alpha_k,\beta_k)$ 满足总长度小于 $\delta$ 时，
\[
 \sum_k V_{\alpha_k}^{\beta_k}(F)\le\varepsilon.
\]
特别地，$F\in BV([a,b])$。
\end{lemma}

\begin{proof}
先取绝对连续定义中对应 $\varepsilon$ 的 $\delta$。在每个
$[\alpha_k,\beta_k]$ 内任取有限分割；所有相邻开子区间仍两两不交，总长度等于
$\sum_k(\beta_k-\alpha_k)<\delta$。绝对连续性说明这些子区间上的绝对增量总和小于
$\varepsilon$。更明确地，若区间族共有 $N$ 个，则对任意 $\eta>0$ 可为每个 $k$ 选分割 $P_k$，使
\[
 V(F;P_k)>V_{\alpha_k}^{\beta_k}(F)-\frac\eta N.
\]
将全部 $P_k$ 的相邻子区间一起代入绝对连续性，得
\[
 \sum_{k=1}^NV_{\alpha_k}^{\beta_k}(F)-\eta
 <\sum_{k=1}^NV(F;P_k)<\varepsilon.
\]
令 $\eta\downarrow0$ 得到所述估计。

为证全局 BV，只需在绝对连续定义中固定例如 $\varepsilon=1$，并取相应 $\delta$。把
$[a,b]$ 分成 $N$ 个长度严格小于 $\delta$ 的基区间。任意分割加入这些基端点后，每个基区间内
的变差和小于 $1$，故整个变差和小于 $N$。对原分割取上确界，得到 $V_a^b(F)\le N<\infty$。
\end{proof}

\begin{definition}[Luzin $N$ 性质]\label{def:4-2-3-2}
函数 $F:[a,b]\to\R$ 具有 Luzin $N$ 性质，如果每个 Lebesgue 零测集 $E\subset[a,b]$
的像 $F(E)$ 仍是 Lebesgue 零测集。
\end{definition}

\begin{proposition}[绝对连续函数的基本性质]\label{proposition:4-2-3-ac-basic}
绝对连续函数连续、属于 $BV([a,b])$，并具有 Luzin $N$ 性质。
\end{proposition}

\begin{proof}
连续性与 BV 已在上面证明。设 $m(E)=0$，给定 $\varepsilon>0$，取
\cref{lemma:4-2-3-local-variation} 对应的 $\delta$。用开集
$G\supset E$ 覆盖，使 $m(G)<\delta$，并把 $G\cap[a,b]$ 写成至多可数个两两不交相对开区间
$J_k$。若某个分量含有 $a$ 或 $b$，先从区间内部截去该端点，再借 $F$ 的连续性令截点趋向端点；
因此上一引理同样适用于任意有限个相对开分量，并给出
$\sum_{k=1}^nV_{J_k}(F)\le\varepsilon$；令 $n\to\infty$，同一估计对全体成立。

由于 $F$ 连续，$F(J_k)$ 包含在一个长度不超过
$\operatorname{osc}_{J_k}F\le V_{J_k}(F)$ 的区间中。这些像区间覆盖 $F(E)$，故
\[
 m^*(F(E))\le\sum_k\operatorname{osc}_{J_k}F
 \le\sum_kV_{J_k}(F)\le\varepsilon.
\]
$\varepsilon$ 任意，遂有 $m(F(E))=0$。
\end{proof}

\begin{definition}[不定积分]\label{def:4-2-3-3}
对实值 $f\in L^1([a,b])$，定义其从 $a$ 起算的不定积分为
\[
 I_f(x)=\int_a^xf(t)\,\dd t.
\]
\end{definition}

\begin{theorem}[Lebesgue 微积分基本定理 I]\label{theorem:4-2-3-1}
若 $f\in L^1([a,b])$ 为实值函数，$F=I_f$，则 $F\in AC([a,b])$，
\[
 F'(x)=f(x)\quad\text{a.e.},
 \qquad
 V_a^b(F)=\int_a^b|f(t)|\,\dd t.
\]
\end{theorem}

\begin{proof}
先证积分对小测度集合一致绝对连续。给定 $\varepsilon>0$，选 $M>0$ 使
\[
 \int_{\{|f|>M\}}|f|<\frac\varepsilon2,
\]
再取 $\delta=\varepsilon/(2M)$。若可测集 $A$ 满足 $m(A)<\delta$，则
\begin{equation}\label{eq:4-2-integral-absolute-continuity}
 \int_A|f|
 \le M m(A)+\int_{\{|f|>M\}}|f|<\varepsilon.
\end{equation}
对定义中的两两不交区间族，令 $A$ 为其并；于是
\[
 \sum_k|F(\beta_k)-F(\alpha_k)|
 \le\sum_k\int_{\alpha_k}^{\beta_k}|f|
 =\int_A|f|<\varepsilon.
\]
所以 $F$ 绝对连续。

把 $f$ 在 $[a,b]$ 外延拓为零。由 Lebesgue 微分定理
\cref{theorem:4-1-2-1}，几乎每个内点 $x$ 是 $f$ 的 Lebesgue 点。对这样的 $x$ 及
$h>0$，
\[
 \left|\frac{F(x+h)-F(x)}h-f(x)\right|
 \le\frac1h\int_x^{x+h}|f(t)-f(x)|\,\dd t
 \le\frac2{2h}\int_{x-h}^{x+h}|f(t)-f(x)|\,\dd t\longrightarrow0.
\]
对 $h>0$，左差商满足
\[
 \left|\frac{F(x)-F(x-h)}h-f(x)\right|
 \le \frac1h\int_{x-h}^{x}|f(t)-f(x)|\,\dd t
 \le \frac2{2h}\int_{x-h}^{x+h}|f(t)-f(x)|\,\dd t,
\]
右端仍由 Lebesgue 点条件趋于零。故双侧差商都趋于 $f(x)$，即
$F'(x)=f(x)$ 于几乎每个 $x$。

对任意分割，三角不等式给
$V(F;P)\le\int_a^b|f|$，所以 $V_a^b(F)\le\int|f|$。为证反向，令
$\lambda=|f|m$，并在 $f\ne0$ 处置 $h=\sgn f$，在 $f=0$ 处置 $h=0$。记
$P_+=\{h=1\}$、$P_-=\{h=-1\}$。有限 Radon 测度 $\lambda$ 的内正则性给出互不相交的紧集
$K_\pm\subset P_\pm$，使
\[
 \lambda\bigl((P_+\cup P_-)\setminus(K_+\cup K_-)\bigr)<\frac\varepsilon2.
\]
由紧集分离可取互不相交的开集 $U_\pm\supset K_\pm$。在各 $U_\pm$ 内用有限多个半开区间覆盖
$K_\pm$，把全部端点细分成一个共同分割 $P=\{x_j\}$；令落在 $U_+$ 中的小区间系数为 $1$，
落在 $U_-$ 中的系数为 $-1$，其余为 $0$。所得阶梯函数
\[
 \phi=\sum_j\epsilon_j\1_{(x_{j-1},x_j]},\qquad
 \epsilon_j\in\{-1,0,1\},
\]
在 $K_+\cup K_-$ 上等于 $h$，而在其他地方 $|h-\phi|\le2$。又因
$\lambda(\{h=0\})=0$，
\[
 \int|h-\phi|\,\dd\lambda
 \le2\lambda\bigl((P_+\cup P_-)\setminus(K_+\cup K_-)\bigr)<\varepsilon.
\]
于是
\[
 \int_a^b|f|-\varepsilon
 <\int_a^b\phi f
 =\sum_j\epsilon_j\int_{x_{j-1}}^{x_j}f
 =\sum_j\epsilon_j(F(x_j)-F(x_{j-1}))
 \le V(F;P)\le V_a^b(F).
\]
令 $\varepsilon\downarrow0$，得到反向不等式。
\end{proof}

\begin{theorem}[Lebesgue 微积分基本定理 II]\label{theorem:4-2-3-2}
若 $F\in AC([a,b])$，则 $F'$ 几乎处处存在、$F'\in L^1([a,b])$，并且对每个
$x\in[a,b]$，
\begin{equation}\label{eq:4-2-ftc-ii}
 F(x)=F(a)+\int_a^xF'(t)\,\dd t.
\end{equation}
\end{theorem}

\begin{proof}
减去常数后可设 $F(a)=0$。由
\cref{lemma:4-2-3-local-variation}，$F\in BV([a,b])$，故
\cref{theorem:4-2-1-4} 给出有限有符号增量测度 $\nu_F$。

我们证明 $|\nu_F|\ll m$。给定 $\varepsilon>0$，取局部变差控制中的 $\delta$。若开集
$G\subset[a,b]$ 满足 $m(G)<\delta$，把它写成两两不交相对开区间 $J_k$ 的并。绝对连续函数连续，
所以 $\nu_F$ 与 $|\nu_F|$ 都没有原子，相对开分量的端点取舍不改变其测度；由变差测度等式在
每个子区间上的版本，
\[
 |\nu_F|(J_k)
 =\sup_{[s,t]\Subset J_k}|\nu_F|((s,t])
 =\sup_{[s,t]\Subset J_k}V_s^t(F)=:V_{J_k}(F),
 \qquad
 |\nu_F|(G)=\sum_kV_{J_k}(F).
\]
任意有限部分不超过 $\varepsilon$，测度从下连续给
$|\nu_F|(G)\le\varepsilon$。若 $m(E)=0$，可把 $E$ 放进这类开集 $G$，故
$|\nu_F|(E)=0$。于是 $|\nu_F|\ll m$。

对 $\nu_F$ 的正负部分分别应用 Radon--Nikodym 定理
\cref{theorem:3-3-1-1}，得到 $f\in L^1([a,b])$ 使
$\nu_F=fm$。区间增量公式给每个 $x$ 上的恒等式
\[
 F(x)-F(a)=\nu_F((a,x])=\int_a^xf(t)\,\dd t.
\]
现在对右侧应用 \cref{theorem:4-2-3-1}，得到 $F'=f$ a.e.；特别地 $F'\in L^1$，并且上式就是
\eqref{eq:4-2-ftc-ii}。注意，先得到增量测度的 $L^1$ 密度，再由 Lebesgue 点推出普通导数；
对称密度的存在不能替代这一步骤。
\end{proof}

\begin{example}[Lipschitz 与绝对连续]\label{ex:4-2-3-1}
若 $F$ 是 $L$-Lipschitz，则任意不交区间族满足
\[
 \sum_k|F(\beta_k)-F(\alpha_k)|
 \le L\sum_k(\beta_k-\alpha_k),
\]
故 $F$ 绝对连续。反向不成立：
\[
 \sqrt x=\int_0^x\frac{1}{2\sqrt t}\,\dd t,
\]
而 $(2\sqrt t)^{-1}\in L^1(0,1)$，所以 $\sqrt x$ 在 $[0,1]$ 上绝对连续；但
$\sqrt x/x=x^{-1/2}\to\infty$，不存在全局 Lipschitz 常数。
\end{example}

\begin{proposition}[绝对连续曲线的长度公式]\label{proposition:4-2-3-curve-length}
设 $\gamma:[a,b]\to\R^m$ 的每个坐标都绝对连续，并定义
\[
 \ell(\gamma)=\sup_P\sum_j|\gamma(x_j)-\gamma(x_{j-1})|.
\]
则 $\gamma'$ a.e. 存在、$|\gamma'|\in L^1$，并且
\[
 \ell(\gamma)=\int_a^b|\gamma'(t)|\,\dd t.
\]
\end{proposition}

\begin{proof}
逐坐标应用 \cref{theorem:4-2-3-2}，得到
$\gamma(t)-\gamma(s)=\int_s^t\gamma'(r)\,\dd r$；有限维范数等价性说明
$|\gamma'|\in L^1$。因此对任意分割，
\[
 \sum_j|\gamma(x_j)-\gamma(x_{j-1})|
 \le\sum_j\int_{x_{j-1}}^{x_j}|\gamma'|
 =\int_a^b|\gamma'|,
\]
故 $\ell(\gamma)\le\int|\gamma'|$。

反向令 $u(t)=\gamma'(t)/|\gamma'(t)|$ 于 $\gamma'(t)\ne0$，其余处令 $u(t)=0$。
记有限 Radon 测度 $\lambda=|\gamma'|\,\dd t$。先以有限值简单函数
$s=\sum_{\ell=1}^La_\ell\1_{E_\ell}$ 逼近 $u$，使 $|a_\ell|\le1$ 且
$\int|u-s|\,\dd\lambda<\varepsilon/2$。对两两不交的值层 $E_\ell$，正则性给出两两不交的紧集
$K_\ell\subset E_\ell$ 及两两不交的开集 $U_\ell\supset K_\ell$，并可使
\[
 \sum_{\ell=1}^L\lambda(E_\ell\setminus K_\ell)<\frac\varepsilon8,
 \qquad
 \sum_{\ell=1}^L\lambda(U_\ell\setminus K_\ell)<\frac\varepsilon8.
\]
在每个 $U_\ell$ 内用有限多个半开区间覆盖 $K_\ell$，细分全部端点，并将同一值层中的小区间合并为 $A_\ell$。
则 $A_\ell\subset U_\ell$两两不交、$K_\ell\subset A_\ell$，且
\[
 \sum_{\ell=1}^L\lambda(E_\ell\triangle A_\ell)
 \le\sum_{\ell=1}^L
 \bigl(\lambda(E_\ell\setminus K_\ell)+\lambda(U_\ell\setminus K_\ell)\bigr)
 <\frac\varepsilon4.
\]
因而区间阶梯向量 $\phi=\sum_\ell a_\ell\1_{A_\ell}$ 满足
\[
 \int|u-\phi|\,\dd\lambda
 \le\int|u-s|\,\dd\lambda
    +2\sum_\ell\lambda(E_\ell\triangle A_\ell)<\varepsilon.
\]
把全部区间端点写成一个共同分割后，这就是
\[
 \phi=\sum_jc_j\1_{(x_{j-1},x_j]},\qquad |c_j|\le1,\qquad
 \int|u-\phi|\,|\gamma'|<\varepsilon.
\]
于是
\[
 \int_a^b|\gamma'|-\varepsilon
 <\int_a^b\langle\phi,\gamma'\rangle
 =\sum_j\left\langle c_j,
   \gamma(x_j)-\gamma(x_{j-1})\right\rangle
 \le\sum_j|\gamma(x_j)-\gamma(x_{j-1})|
 \le\ell(\gamma).
\]
令 $\varepsilon\downarrow0$ 得到反向不等式。
\end{proof}

若 $\gamma$ 在每个分割段以常向量 $v_j$ 为速度，则公式右端为
$\sum_j|v_j|(x_j-x_{j-1})$，正好是折线边长之和；停驻区间的零速度也确实不给长度。

先构造原函数、再应用微积分基本定理的方法还给出一维换元。只写形式微分
$\dd s=\phi'(t)\dd t$ 会掩盖两个问题：复合函数是否仍绝对连续，以及链式法则在哪些点成立。

\begin{proposition}[绝对连续参数下的一维换元]\label{proposition:4-2-3-change-of-variables}
设 $\phi\in AC([a,b])$，$h$ 在包含 $\phi([a,b])$ 的区间上连续。则
\[
 \int_a^b h(\phi(t))\phi'(t)\,\dd t
 =\int_{\phi(a)}^{\phi(b)}h(s)\,\dd s,
\]
右侧按有向区间理解。特别地，若 $\phi$ 非减，这就是从 $t$--区间到
$[\phi(a),\phi(b)]$ 的通常换元公式。
\end{proposition}

\begin{proof}
在包含 $\phi([a,b])$ 的紧区间上定义
\[
 H(y)=\int_{\phi(a)}^y h(s)\,\dd s.
\]
普通微积分基本定理给 $H\in C^1$ 与 $H'=h$；在这个紧区间上，$H$ 还是
$\|h\|_\infty$-Lipschitz。若一族不交区间的总长度足够小，$\phi$ 的绝对连续性使
$\sum|\Delta\phi|$ 任意小，进而
\[
 \sum|\Delta(H\circ\phi)|
 \le\|h\|_\infty\sum|\Delta\phi|
\]
也任意小。因此 $H\circ\phi\in AC([a,b])$。在 $\phi$ 可微的几乎处处点，经典一元链式法则给
\[
 (H\circ\phi)'(t)=H'(\phi(t))\phi'(t)
 =h(\phi(t))\phi'(t).
\]
对 $H\circ\phi$ 应用 \cref{theorem:4-2-3-2}，便得
\[
 \int_a^b h(\phi(t))\phi'(t)\,\dd t
 =H(\phi(b))-H(\phi(a))
 =\int_{\phi(a)}^{\phi(b)}h(s)\,\dd s.
\]
\end{proof}

\begin{example}[平方参数换元]
在 $[0,1]$ 上取 $\phi(t)=t^2$ 与 $h(s)=(1+s)^{-1}$。命题给出
\[
 \int_0^1\frac{2t}{1+t^2}\,\dd t
 =\int_0^1\frac{\dd s}{1+s}=\log2.
\]
左边也可直接令 $u=1+t^2$ 验算；这次有限计算同时核对了导数因子与端点方向。
\end{example}

对一般 BV 函数，增量测度的 Lebesgue 分解还可以进一步把奇异质量中的原子逐个取出。

\begin{theorem}[BV 的三部分分解]\label{theorem:4-2-3-3}
设 $F$ 是右连续、$F(a)=0$ 的 BV 函数。存在唯一分解
\[
 F=F_{ac}+F_j+F_{sc},
\]
其中 $F_{ac}$ 绝对连续，$F_j$ 是绝对收敛跳跃级数的累计函数，$F_{sc}$ 连续且其增量测度
奇异于 Lebesgue 测度、没有原子。相应增量测度分解为
\[
 \nu_F=fm+\sum_{x\in(a,b]}\nu_F(\{x\})\delta_x+\nu_{sc}.
\]
\end{theorem}

\begin{proof}
对 $\nu_F$ 的正负部分分别应用 Lebesgue 分解定理
\cref{theorem:3-3-1-2}，再合并密度与奇异部分，得到
\[
 \nu_F=fm+\nu_s,\qquad f\in L^1([a,b]),\qquad \nu_s\perp m.
\]
有限测度 $|\nu_s|$ 的原子至多可数。更具体地，对每个 $n$，满足
$|\nu_s(\{x\})|\geq1/n$ 的原子只有有限个；
对 $n$ 取并即可。并且
\[
 \sum_{x\in(a,b]}|\nu_s(\{x\})|
 \le |\nu_s|((a,b])<\infty.
\]
所以
\[
 \nu_j=\sum_x\nu_s(\{x\})\delta_x
\]
在全变差范数下收敛，$\nu_{sc}:=\nu_s-\nu_j$ 是有限、奇异且无原子的有符号测度。

分别取规范累计函数
\[
 F_{ac}(x)=\int_a^xf(t)\,\dd t,
 \quad
 F_j(x)=\sum_{a<y\le x}\nu_s(\{y\}),
 \quad
 F_{sc}(x)=\nu_{sc}((a,x]).
\]
\Cref{theorem:4-2-3-1} 给出 $F_{ac}\in AC$；绝对收敛给 $F_j$ 右连续并具有所列跳跃；
$\nu_{sc}$ 无原子，
故其累计函数从左、从右都连续。三个增量测度之和等于 $\nu_F$，且三者在 $a$ 都为零，
\cref{theorem:4-2-1-4} 的唯一性给 $F=F_{ac}+F_j+F_{sc}$。

若另有一个这类分解，取增量测度后，Lebesgue 分解的唯一性先确定绝对连续项和奇异项；奇异测度的
每个原子质量又由单点取值唯一确定，故纯跳跃与无原子奇异项也分别相同。再用共同规范恢复函数，
得到唯一性。
\end{proof}

Cantor 函数就是 $F_{sc}$ 可以非零的最小模型。阶梯函数则只有 $F_j$。普通 a.e. 导数只返回
$F_{ac}'=f$；它既不返回跳跃，也不返回 Cantor 型增长。

为了刻画绝对连续性，还需要解决一个不易被一句直觉替代的问题：若连续 BV 函数把零测集仍送到
零测集，为什么它的\emph{变差测度}也不能藏在零测集上？函数本身的 $N$ 性质并不会形式地传给
Jordan 分量。答案是一维 Banach indicatrix 公式，它按高度计算曲线穿越次数。

\begin{lemma}[一维 Banach indicatrix 公式]\label{lemma:4-2-3-banach-indicatrix}
设 $F:[a,b]\to\R$ 连续且属于 $BV([a,b])$；形成增量测度时以
$F-F(a)$ 作规范化，并把 $(a,b]$ 上的测度 $|\nu_{F-F(a)}|$ 以
$\mu(\{a\})=0$ 延拓到 $[a,b]$，所得测度记为 $\mu$。减常数不改变下文的变差或穿越重数。
存在可测的穿越重数
$N(F,E,y)\in\{0,1,2,\ldots,\infty\}$，使每个 Borel 集 $E\subset[a,b]$ 满足
\begin{equation}\label{eq:4-2-banach-indicatrix}
 \mu(E)=\int_\R N(F,E,y)\,\dd y.
\end{equation}
而且 $N(F,E,y)=0$ 对 a.e. $y\notin F(E)$ 成立。对折线且 $y$ 不取顶点高度时，
$N(F,E,y)$ 就是参数落在 $E$ 中、曲线穿过水平线 $y$ 的次数；平坦段造成的无穷多个普通原像
只出现在至多可数个平坦高度上，不改变积分。
\end{lemma}

\begin{proof}
取含 $a,b$ 的可数稠密集 $D\subset[a,b]$，并记 $Z=F(D)$；$Z$ 可数，因而是零测集。
若 $s<t$ 属于 $D$，对顶点都在 $D$ 中的分割
$Q=\{s=t_0<\cdots<t_r=t\}$ 定义
\[
 n_Q(y)=\sum_{j=1}^r
 \1_{(\min\{F(t_{j-1}),F(t_j)\},\,
          \max\{F(t_{j-1}),F(t_j)\})}(y).
\]
每一项的积分就是相应纵向增量的长度，所以
\begin{equation}\label{eq:4-2-indicatrix-polygonal}
 \int_\R n_Q(y)\,\dd y
 =\sum_{j=1}^r|F(t_j)-F(t_{j-1})|.
\end{equation}
若在一段中插入顶点 $d\in D$，则对 $y\notin Z$，旧段跨过高度 $y$ 时两个新段至少有一个跨过
$y$；旧段不跨过而新顶点落在另一侧时则新增两次。因此细化使 $n_Q(y)$ 不减。

定义
\[
 N_{s,t}(y)=\sup_Q n_Q(y)\quad(y\notin Z),\qquad N_{s,t}(y)=0\quad(y\in Z),
\]
上确界取上述所有 $D$-顶点分割。这样的分割只有可数多个，所以 $N_{s,t}$ 可测，且定义与任何枚举
无关。枚举全部分割并逐次取共同细化，所得穿越数单调增加到这个上确界。连续性还说明，把任意有限
分割的顶点移到 $D$ 中可使变差和的改变量任意小。详细地，给定
$P=\{s=x_0<\cdots<x_m=t\}$ 与 $\eta>0$，由 $F$ 的一致连续性和 $D$ 的稠密性，可保持次序地取
\[
 d_0=s,\quad d_m=t,\quad d_i\in D,\qquad
 |F(d_i)-F(x_i)|<\frac{\eta}{2m}\quad(1\le i<m).
\]
于是
\[
 \begin{aligned}
 &\left|\sum_{i=1}^m|F(d_i)-F(d_{i-1})|
       -\sum_{i=1}^m|F(x_i)-F(x_{i-1})|\right|\\
 &\qquad\le
 \sum_{i=1}^m\bigl(|F(d_i)-F(x_i)|
                   +|F(d_{i-1})-F(x_{i-1})|\bigr)<\eta.
 \end{aligned}
\]
故 $D$-分割给出的变差上确界仍是 $V_s^t(F)$。单调收敛定理与
\eqref{eq:4-2-indicatrix-polygonal} 因而给出
\begin{equation}\label{eq:4-2-indicatrix-interval}
 \int_\R N_{s,t}(y)\,\dd y=V_s^t(F)=\mu((s,t]).
\end{equation}
最后一个等号来自变差测度等式在子区间上的应用；$F$ 连续使 $\nu_F$ 及 $\mu$ 都没有原子，
所以端点取舍无关。

若 $r<s<t$ 都在 $D$ 中，左右两个分割合并后给出
$N_{r,t}(y)\ge N_{r,s}(y)+N_{s,t}(y)$。反过来，在 $[r,t]$ 的任意分割中插入 $s$：一条旧线段
若跨过 $y$，分成的两条新线段中恰有一条跨过 $y$；若旧线段不跨过 $y$，插点至多新增两次穿越。
因此细化后的穿越数不小于原数，而细化后的数等于其左右两部分之和，至多为
$N_{r,s}(y)+N_{s,t}(y)$。对原分割取上确界，得到反向不等式。故对 $y\notin Z$ 有精确可加性
\begin{equation}\label{eq:4-2-indicatrix-additivity}
 N_{r,t}(y)=N_{r,s}(y)+N_{s,t}(y).
\end{equation}
由 \eqref{eq:4-2-indicatrix-interval}，
$k(y):=N_{a,b}(y)<\infty$ 对 a.e. $y$ 成立。删去 $Z$ 与这个零测例外集后固定 $y$，简记
$k=k(y)$。

下面把区间上的整数权构造成有限测度。写
$D=\{d_1,d_2,\ldots\}$。递归构造有限 $D$-分割：令 $P_m$ 细化 $P_{m-1}$，包含
$d_1,\ldots,d_m$，并另外加入有限多个 $D$ 中的点使网格小于 $1/m$；稠密性保证最后一步可以完成。
于是 $(P_m)$ 逐次细化、网格趋于零，并且每个 $D$ 中的点从某一层起永远是顶点。给 $P_m$ 的每个半开小区间
$I=(s,t]$ 赋整数权
\[
 w_y(I)=N_{s,t}(y).
\]
式 \eqref{eq:4-2-indicatrix-additivity} 说明父区间的权等于其所有子区间权之和，且每一层的总权
都是 $k$。在第一层放置 $k$ 个标号，并先按 $P_1$ 各小区间的整数权分配，使每个小区间恰携带
与其权重相同数目的标号；每次细化时，再把每个父区间携带的标号按各子区间的整数权分配下去。
由父权等于子权之和归纳可知，每一层的每个小区间始终恰携带与其权重相同数目的标号。
每个标号由此确定一列闭包递减、直径趋于零的小区间，其交为一个点 $x_i(y)$。

标号所在的每个小区间权至少为一。按 $N_{s,t}$ 的上确界定义，其中存在一个有限分割，至少有一段的
两个端值严格分处 $y$ 两侧；连续性给出该小区间闭包中的一点 $z_m$ 满足 $F(z_m)=y$。
随网格趋零，$z_m\to x_i(y)$，故 $F(x_i(y))=y$。特别地
$x_i(y)\notin D$，因为 $y\notin F(D)$。这也消除了半开区间端点的歧义。若若干标号收敛到同一点，
仍分别计入其重数，并定义
\[
 \kappa_y=\sum_{i=1}^k\delta_{x_i(y)}.
\]
当 $s,t\in D$ 同时成为 $P_m$ 的顶点后，分配规则表明，落在 $(s,t]$ 内的标号数恰为
$N_{s,t}(y)$；所以
\[
 \kappa_y((s,t])=N_{s,t}(y).
\]
对先前删去的高度置 $\kappa_y=0$。这样每个 $\kappa_y$ 从定义起就是有限 Borel 测度，而且
\[
 \supp\kappa_y\subset F^{-1}(\{y\}).
\]

在 $D$-端点半开区间上，$y\mapsto\kappa_y(A)$ 可测；有限不交并的 $\kappa_y$ 值是相应
区间值的有限和，故同一结论成立于这些区间组成的环。还要核对从这个环推广到全部
Borel 集时没有暗中要求原子位置 $x_i(y)$ 本身可测。令 $\mathcal C$ 为所有满足
$y\mapsto\kappa_y(E)$ 可测的 Borel 集 $E$。先有
$[a,b]\in\mathcal C$，因为 $\kappa_y([a,b])$ 等于可测函数 $N_{a,b}(y)$ 在保留高度集上的限制，
在删去的高度上为零。若 $E\in\mathcal C$，则
\[
 \kappa_y([a,b]\setminus E)=\kappa_y([a,b])-\kappa_y(E)
\]
可测；若 $(E_j)$ 两两不交且都属于 $\mathcal C$，则
$\kappa_y(\bigcup_jE_j)=\sum_j\kappa_y(E_j)$ 是非负可测函数。因此 $\mathcal C$ 是一个
$\lambda$-系统，并包含
\[
 \mathcal P=\{[a,b],\varnothing\}
 \cup\{(s,t]:s,t\in D,\ s<t\}.
\]
这是一个生成 $\Borel([a,b])$ 的 $\pi$-系统。
$\pi$--$\lambda$ 定理给出 $\mathcal C=\Borel([a,b])$。于是可定义
\[
 \lambda(E)=\int_\R\kappa_y(E)\,\dd y.
\]
若 $(E_j)$ 两两不交，则
$\kappa_y(\bigcup_jE_j)=\sum_j\kappa_y(E_j)$；对非负级数使用单调收敛定理，得到
$\lambda(\bigcup_jE_j)=\sum_j\lambda(E_j)$。此外，$\kappa_y([a,b])$ 与 $k(y)$ 在删去的
零测高度集之外相等，因而
$\lambda([a,b])=\int k(y)\dd y=V_a^b(F)<\infty$；所以 $\lambda$ 是有限 Borel 测度。
由 \eqref{eq:4-2-indicatrix-interval}，$\lambda$ 与 $\mu$ 在所有 $D$-端点半开区间上相等。
此外 $a\in D$ 且保留高度满足 $y\notin F(D)$，故所有 $x_i(y)$ 都不等于 $a$；删去高度又给零测度，
所以 $\lambda(\{a\})=0=\mu(\{a\})$。这些半开区间与 $\{a\}$ 生成
$\Borel([a,b])$，有限测度唯一性遂给 $\lambda=\mu$。定义
\[
 N(F,E,y)=\kappa_y(E),
\]
就得到 \eqref{eq:4-2-banach-indicatrix}。支撑包含关系还说明
$y\notin F(E)$ 时 $N(F,E,y)=0$。

若 $F$ 是折线且 $y$ 不等于任何顶点高度，每个非水平线段至多贡献一次，以上有限分割计数就是图像
穿过水平线的次数。平坦段只有在其平坦高度才造成无穷多个普通原像；有限折线只有有限个这种高度，
一般连续 BV 函数中相应例外也已被上述零测高度集排除。特别地，对开集 $U$，
\begin{equation}\label{eq:4-2-indicatrix-open}
 \int N(F,U,y)\,\dd y=\mu(U),
\end{equation}
并且
\begin{equation}\label{eq:4-2-indicatrix-domination}
 N(F,U,y)\le N(F,[a,b],y),
\end{equation}
左侧被积函数只可能在 $F(U)$ 上非零。
\end{proof}

\begin{lemma}[连续 BV 的变差--$N$ 引理]\label{lemma:4-2-3-variation-N}
设 $F$ 连续且有界变差。若 $F$ 具有 Luzin $N$ 性质，则
$|\nu_{F-F(a)}|\ll m$。
\end{lemma}

\begin{proof}
记 $\mu=|\nu_{F-F(a)}|$，并令
$N(y)=N(F,[a,b],y)$。由
\cref{lemma:4-2-3-banach-indicatrix}，$N\in L^1(\R)$。
先设 $E\subset[a,b]$ 为 Borel 集且 $m(E)=0$。由 $N$ 性质，$m(F(E))=0$。

先取紧集 $K\subset E$。给定 $\varepsilon>0$，积分的绝对连续性给出开集
$V\supset F(K)$，使 $\int_VN(y)\,\dd y<\varepsilon$。具体地，先取
$\delta>0$ 使 $m(A)<\delta$ 蕴含 $\int_A N<\varepsilon$；$F(K)$ 是紧零测集，故
Lebesgue 外正则性允许取开集 $V\supset F(K)$ 且 $m(V)<\delta$。集合
$U=F^{-1}(V)$ 是含 $K$ 的相对开集。由
\eqref{eq:4-2-indicatrix-open}、\eqref{eq:4-2-indicatrix-domination} 及局部穿越只落在 $F(U)$ 中，
\[
 \mu(K)\le\mu(U)
 =\int N(F,U,y)\,\dd y
 \le\int_{F(U)}N(y)\,\dd y
 \le\int_VN(y)\,\dd y<\varepsilon.
\]
故 $\mu(K)=0$。有限 Borel 测度在紧区间上内正则，于是
$\mu(E)=\sup\{\mu(K):K\subset E,\ K\text{ 紧}\}=0$。
若 $A$ 只是 Lebesgue 可测零集，取 Borel 零集 $E\supset A$；$\mu(E)=0$。有限 Borel 测度
$\mu$ 的 Lebesgue 完成化因而在 $A$ 上也取零；等价地，对绝对连续性所需的每个 Borel 零集已经有
$\mu(E)=0$。这证明 $\mu\ll m$。
\end{proof}

\begin{example}[重数不能删掉]\label{ex:4-2-3-indicatrix}
在 $[0,1]$ 上取帐篷函数
\[
 F(x)=1-|2x-1|.
\]
它的像集是 $[0,1]$，长度为 $1$；但它先以斜率 $2$ 上升、再以斜率 $-2$ 下降，故
$V_0^1(F)=2$。对几乎每个 $y\in(0,1)$，水平线与图像相交两次，所以
\[
 \int_0^1N(F,[0,1],y)\,\dd y=\int_0^12\,\dd y=2.
\]
只量像集而不数穿越次数会恰好少算一半。
\end{example}

\begin{proposition}[Banach--Zarecki 判据]\label{proposition:4-2-3-5}
函数 $F:[a,b]\to\R$ 绝对连续，当且仅当它连续、属于 $BV([a,b])$，并具有 Luzin $N$ 性质。
\end{proposition}

\begin{proof}
正向由 \cref{proposition:4-2-3-ac-basic} 得到。

反向设三个条件成立。减去 $F(a)$ 后使用
\cref{theorem:4-2-1-4} 得到增量测度 $\nu_F$。由
\cref{lemma:4-2-3-variation-N}，$|\nu_F|\ll m$；对正负部分应用 Radon--Nikodym 定理，
存在 $f\in L^1([a,b])$ 使 $\nu_F=fm$。于是对每个 $x$，
\[
 F(x)-F(a)=\nu_F((a,x])=\int_a^xf(t)\,\dd t.
\]
\Cref{theorem:4-2-3-1} 说明右侧绝对连续，所以 $F$ 绝对连续。
\end{proof}

判据中的两个附加条件都不可删。非恒定阶梯函数属于 BV，并且把任何零测集送到有限集或可数集，
因而有 $N$ 性质，但它不连续。Cantor 函数连续且 BV，却把零测 Cantor 集映到整个 $[0,1]$，
所以没有 $N$ 性质。

\begin{proposition}[积分部分公式]\label{proposition:4-2-3-4}
若 $F,G\in AC([a,b])$，则
\[
 F(b)G(b)-F(a)G(a)
 =\int_a^bF'(x)G(x)\,\dd x
  +\int_a^bF(x)G'(x)\,\dd x.
\]
\end{proposition}

\begin{proof}
绝对连续函数在紧区间上有界。对任意不交区间族，逐段使用精确恒等式
\[
 F(\beta)G(\beta)-F(\alpha)G(\alpha)
 =F(\beta)(G(\beta)-G(\alpha))
  +G(\alpha)(F(\beta)-F(\alpha)).
\]
因此，只要总长度充分小，所有乘积增量的绝对值之和至多
\[
 \|F\|_\infty\sum|\Delta G|
 +\|G\|_\infty\sum|\Delta F|,
\]
并可由 $F,G$ 的绝对连续性控制为任意小。具体地，给定 $\varepsilon>0$，分别取绝对连续性中对应于
\[
 \frac{\varepsilon}{2(1+\|G\|_\infty)}
 \quad\text{和}\quad
 \frac{\varepsilon}{2(1+\|F\|_\infty)}
\]
的 $\delta_F,\delta_G$。当区间族总长度小于
$\min\{\delta_F,\delta_G\}$ 时，
\[
 \sum|\Delta(FG)|
 \le \|F\|_\infty\frac{\varepsilon}{2(1+\|F\|_\infty)}
    +\|G\|_\infty\frac{\varepsilon}{2(1+\|G\|_\infty)}
 <\varepsilon.
\]
故 $FG\in AC([a,b])$。
在 $F,G$ 同时可微的满测集上，普通乘积法则给
$(FG)'=F'G+FG'$。右侧属于 $L^1$，因为 $F,G$ 有界。对 $FG$ 应用
\cref{theorem:4-2-3-2}，得到
\[
 F(b)G(b)-F(a)G(a)=\int_a^b(FG)'
 =\int_a^bF'G+\int_a^bFG'.
\]
\end{proof}

\begin{example}[Cantor 函数的失败]\label{ex:4-2-3-2}
Cantor 函数连续、非减且 $C'=0$ a.e.，但 $C(1)-C(0)=1$。若错误地对它应用
\cref{theorem:4-2-3-2}，右侧会给
$\int_0^1C'=0$。失败点可由 Banach--Zarecki 判据精确定位：零测 Cantor 集被 $C$ 映满
$[0,1]$，所以 $C$ 不具 $N$ 性质，也不绝对连续。
\end{example}

\begin{example}[跳跃函数的失败]\label{ex:4-2-3-3}
$H=\1_{[0,1]}$ 在 $[-1,1]$ 上属于 BV，且 $H'=0$ 于 $0$ 以外每点；但
$H(1)-H(-1)=1$。它在 $0$ 不连续，增量测度为 $\delta_0$，所以 Banach--Zarecki 判据的连续性假设
也不能删除。
\end{example}

\begin{figure}[htbp]
\centering
\begin{tikzpicture}[x=.85cm,y=1.0cm]
  \draw[->] (0,0)--(12.4,0) node[right]{$x$};
  \draw[->] (0,0)--(0,4.5) node[above]{$F$};
  \draw[BookAccent,thick] (0,.2).. controls (2,.5) and (3,1.45)..(4,1.7)
    .. controls (6,2.0) and (8,3.15)..(12,3.7);
  \draw[BookWarning,thick]
    (0,0)--(2.2,0)--(2.2,.55)--(5.2,.55)--(5.2,1.15)--(8.6,1.15)--(8.6,1.8)--(12,1.8);
  \draw[BookWarm,thick]
    (0,.1)--(1.35,.1)--(1.35,.45)--(2.7,.45)--(2.7,.8)--(4.05,.8)--(4.05,1.15)--
    (5.4,1.15)--(5.4,1.5)--(6.75,1.5)--(6.75,1.85)--(8.1,1.85)--(8.1,2.2)--
    (9.45,2.2)--(9.45,2.55)--(10.8,2.55)--(10.8,2.9)--(12,2.9);
  \node[book strong node] at (2.1,3.35){$F_{ac}$：$L^1$ 斜率};
  \node[book warning node] at (6.2,4.0){$F_j$：离散跳跃};
  \node[book warm node] at (10.1,3.4){$F_{sc}$：奇异连续};
\end{tikzpicture}
\caption{BV 增量的三个分量：绝对连续密度、原子跳跃与无原子的奇异连续增长。图中三条曲线为分量，
其逐点相加才是原函数。}
\label{fig:bv-three-parts}
\end{figure}

\begin{remark}[混合测度的累计函数]
在 $[0,1]$ 上取有限正测度
\[
 \nu=p\delta_{1/2}+q\mu_C+(1-p-q)m|_{[0,1]},
 \qquad p,q\ge0,\quad p+q\le1.
\]
其规范累计函数含一个大小为 $p$ 的跳跃、一个总增长为 $q$ 的 Cantor 分量，以及斜率为
$1-p-q$ 的绝对连续分量。a.e. 导数只返回最后一项。这也是概率分布函数中混合型分布的标准模型。
\end{remark}

稍后定义 $W^{1,1}(a,b)$ 后，\cref{theorem:4-2-3-2} 将推出每个一维 $W^{1,1}$ 等价类都有唯一的
绝对连续代表；
从导数重建函数仍需另给一个加法常数。

\begin{exercise}
从绝对连续定义直接证明 Lipschitz 函数绝对连续，并验证 $x\mapsto\sqrt x$ 给出反例说明逆命题失败。
\end{exercise}
\begin{exercise}
验证 Cantor 函数属于 $BV([0,1])$，并直接用定义~\ref{def:4-2-3-1} 证明它不绝对连续：对每一层
Cantor 基本区间取其内部稍作缩短的有限不交子区间，使总长度趋于零，同时使端点函数增量之和保持
任意接近 $1$。
\end{exercise}
\begin{exercise}
完整证明 \eqref{eq:4-2-integral-absolute-continuity}，并据此证明 $I_f$ 绝对连续。
\end{exercise}
\begin{exercise}
证明 Cantor 函数把 Cantor 集映到 $[0,1]$，从而不具有 Luzin $N$ 性质。
\end{exercise}
\begin{exercise}
先证明 $FG\in AC([a,b])$，再独立推导积分部分公式；指出若把 $F,G$ 换成具有共同跳跃的 BV 函数，
乘积法则中会出现哪一项。
\end{exercise}

绝对连续性把“增量测度有密度”和“函数可由导数积分恢复”变成同一件事。粗糙函数即使不绝对连续，
仍可通过平移平均获得光滑近似；下一节将把这一有限平均过程组织成 mollifier 与卷积理论。

\section{光滑逼近、分布与 Sobolev 空间}\label{sec:04-03}

\subsection{Mollifier、截断与光滑逼近}\label{subsec:4-3-1}
\index{mollifier}\index{光滑逼近}\index{近似恒等}

许多形式计算要求对积分中的函数求导，而被积函数可能没有经典导数。自然的处理方法是在小球内作
加权平均，但这需要分别证明三个结论：积分有定义，微分可以转移到核上，以及平均尺度趋于零时平均
在指定拓扑中收敛到原函数。若函数只定义在开集上，还必须保证靠近边界的平均窗口不越出定义域。

Friedrichs 光滑化把这种小尺度加权平均组织成可重复使用的方法：先在光滑对象上完成微分与积分
运算，再通过范数收敛返回原对象。其正确使用依赖对支撑、尺度与收敛方式的明确控制。

先从最简单的跳跃开始。在一维中，取支撑于
$[-\varepsilon,\varepsilon]$、积分为 $1$ 的非负偶核 $\rho_\varepsilon$，并把
$f=\1_{(0,1)}$ 看成 $\R$ 上的函数，则
\begin{equation}\label{eq:4-3-step-mollification}
 (f*\rho_\varepsilon)(x)
 =\int_{x-1}^{x}\rho_\varepsilon(y)\dd y.
\end{equation}
当 $\varepsilon<x<1-\varepsilon$ 时右端等于 $1$；在 $(0,\varepsilon)$ 与
$(1-\varepsilon,1)$ 中，它分别从 $1/2$ 过渡到 $1$、从 $1$ 过渡到 $1/2$。竖直跳跃被摊进
宽度至多 $2\varepsilon$ 的过渡层，但远离跳点的数值没有改变。

\begin{definition}[标准 mollifier]\label{def:4-3-1-1}
令
\[
 \eta(t)=
 \begin{cases}e^{-1/t},&t>0,\\0,&t\leq0,
 \end{cases}
 \qquad
 \rho(x)=c_n\eta(1-|x|^2),
\]
其中
\[
 c_n^{-1}=\int_{|x|<1}\exp\!\left(-\frac1{1-|x|^2}\right)\dd x.
\]
则 $\rho\in C_c^\infty(\R^n)$、$\rho\geq0$、$\supp\rho\subset\overline{B(0,1)}$ 且
$\int\rho=1$。对 $\varepsilon>0$ 定义
\[
 \rho_\varepsilon(x)=\varepsilon^{-n}\rho(x/\varepsilon).
\]
称 $(\rho_\varepsilon)$ 为标准 mollifier 族。
\end{definition}

这里的光滑性并非来自对分段公式的乐观判断。对每个 $k\geq0$，$\eta^{(k)}(t)$ 在 $t>0$
上等于 $e^{-1/t}$ 乘以 $t^{-1}$ 的一个多项式。确切地，存在多项式 $P_k$，使
\[
 \eta^{(k)}(t)=e^{-1/t}P_k(t^{-1})\qquad(t>0).
\]
这可归纳证明：$P_0=1$，若 $s=t^{-1}$，则
\[
 \frac{\dd}{\dd t}\bigl(e^{-s}P_k(s)\bigr)
 =e^{-s}s^2\bigl(P_k(s)-P_k'(s)\bigr),
\]
故 $P_{k+1}(s)=s^2(P_k(s)-P_k'(s))$ 仍为多项式。又对每个 $N\ge0$，令 $s=1/t$ 得
\[
 t^{-N}e^{-1/t}=s^Ne^{-s}\longrightarrow0\qquad(t\downarrow0).
\]
因而每个 $\eta^{(k)}(t)$ 在 $t\downarrow0$ 时趋于零。更具体地，若已定义
$\eta^{(k-1)}(0)=0$，则
\[
 \frac{\eta^{(k-1)}(t)-\eta^{(k-1)}(0)}t
 =e^{-1/t}t^{-1}P_{k-1}(t^{-1})\longrightarrow0;
\]
当 $t<0$ 时相应差商恒为零。因此从 $k=1$ 起归纳可知，$\eta$ 在零点的每阶双侧导数都存在且等于零。所以上述
$\rho$ 的确是光滑紧支撑函数。换元还给出
\[
 \int\rho_\varepsilon=1,\qquad
 \supp\rho_\varepsilon\subset\overline{B(0,\varepsilon)}.
\]

\begin{example}[阶跃函数平滑]\label{ex:4-3-1-1}
公式 \eqref{eq:4-3-step-mollification} 给出了完整计算。若 $x\notin\{0,1\}$，则对充分小的
$\varepsilon$，卷积值已经等于 $\1_{(0,1)}(x)$。在两个跳点，由 $\rho$ 的偶对称性，
卷积值趋于 $1/2$，而不是所选代表在该点的值。因而点态极限首先识别的是 Lebesgue 点，不能用来
恢复零测集上的任意赋值。
\end{example}

\begin{example}[Dirac 质量的平滑]\label{ex:4-3-1-2}
$\rho_\varepsilon$ 本身是总质量为一、支撑收缩到零的光滑密度。下一小节定义分布卷积后，将由定义
直接得到 $\delta_0*\rho_\varepsilon=\rho_\varepsilon$。此处只需观察右端的显式公式。
\end{example}

若 $f$ 只定义在开集内，卷积窗口必须完整留在域中。

\begin{definition}[卷积平滑]\label{def:4-3-1-2}
设 $\Omega\subset\R^n$ 开，$f\in L^1_{\mathrm{loc}}(\Omega)$。当
$\dist(x,\Omega^c)>\varepsilon$ 时，定义
\[
 f_\varepsilon(x)
 =\int_{B(0,\varepsilon)}f(x-y)\rho_\varepsilon(y)\dd y.
\]
若 $f$ 已定义在整个 $\R^n$，则也写作 $f*\rho_\varepsilon$。
\end{definition}

\begin{definition}[内部耗尽]\label{def:4-3-1-3}
对 $\varepsilon>0$，置
\[
 \Omega_\varepsilon
 =\{x\in\Omega:\dist(x,\Omega^c)>\varepsilon, |x|<\varepsilon^{-1}\}.
\]
这是开集，并且 $\overline{\Omega_\varepsilon}$ 是 $\Omega$ 中的紧集。若
$0<\varepsilon'<\varepsilon$，则 $\Omega_\varepsilon\subset\Omega_{\varepsilon'}$，且
$\bigcup_{k\geq1}\Omega_{1/k}=\Omega$。
\end{definition}

最后一句可由 Heine--Borel 定理验证：闭包既闭又有界，而距离条件在取闭包后变成
$\dist(x,\Omega^c)\geq\varepsilon$，所以闭包仍包含于 $\Omega$。因此
$\overline{\Omega_\varepsilon}\Subset\Omega$。

以后使用标准多重指标记号：若
$\alpha=(\alpha_1,\ldots,\alpha_n)\in\mathbb N_0^n$，则
$|\alpha|=\alpha_1+\cdots+\alpha_n$，并令
\[
 \partial^\alpha
 =\partial_1^{\alpha_1}\cdots\partial_n^{\alpha_n}.
\]
$\alpha=0$ 时把 $\partial^\alpha$ 读作恒等运算。

\begin{proposition}[卷积光滑性]\label{proposition:4-3-1-1}
若 $f\in L^1_{\mathrm{loc}}(\R^n)$、$\psi\in C_c^\infty(\R^n)$，则
$f*\psi\in C^\infty(\R^n)$，且对每个多重指标 $\alpha$，
\[
 \partial^\alpha(f*\psi)=f*(\partial^\alpha\psi).
 \tag{4.3.1}\label{eq:4-3-convolution-differentiate}
\]
若 $f$ 只定义在 $\Omega$，同一结论在所有卷积窗口紧含于 $\Omega$ 的区域成立。
\end{proposition}

\begin{proof}
把卷积写成
\[
 (f*\psi)(x)=\int f(y)\psi(x-y)\dd y.
\]
固定 $x_0$，取闭球 $K$ 使 $x$ 在 $x_0$ 的一个邻域中变化时，所有
$x-\supp\psi$ 都包含在 $K$ 中。于是积分只涉及 $f\1_K\in L^1$。以第一坐标为例，均值定理给出
\[
 \left|\frac{\psi(x+he_1-y)-\psi(x-y)}h\right|
 \leq\|\partial_1\psi\|_\infty
\]
（$h\ne0$），而差商逐点趋于 $\partial_1\psi(x-y)$。因此
$|f(y)|\|\partial_1\psi\|_\infty\1_K(y)$ 是可积支配函数，支配收敛定理给出
\[
 \partial_1(f*\psi)(x)=\int f(y)\partial_1\psi(x-y)\dd y.
\]
同一支配论证还说明右端关于 $x$ 连续。把 $\psi$ 依次换成它的各阶偏导，归纳即得
\eqref{eq:4-3-convolution-differentiate}。若论域是 $\Omega$，固定卷积窗口紧含于 $\Omega$ 的点
$x_0$，再取其一个足够小的邻域 $U$，使
\[
 \bigcup_{x\in U}(x-\supp\psi)\subset K\Subset\Omega.
\]
上面的所有差商便都由 $|f|\1_K$ 乘相应核导数的一致界支配，所以同一结论在该区域成立。
\end{proof}

\begin{figure}[htbp]
  \centering
  \begin{tikzpicture}[x=1.15cm,y=1.15cm]
    \draw[->,BookInk!75] (-.3,0)--(6.35,0) node[right] {$x$};
    \draw[->,BookInk!75] (0,-.15)--(0,2.15);
    \draw[BookWarning,semithick] (0,0)--(1,0)--(1,1.55)--(2.35,1.55)--(2.35,.55)--(3.4,.55)--(3.4,1.25)--(4.7,1.25)--(4.7,0)--(6,0);
    \draw[BookAccent,very thick,smooth] plot coordinates {(.45,.03) (.8,.15) (1.05,.75) (1.35,1.35) (1.8,1.48) (2.15,1.25) (2.45,.85) (2.8,.65) (3.2,.72) (3.5,1.0) (4.25,1.15) (4.65,.72) (5,.15) (5.4,.03)};
    \draw[BookWarm,semithick] (5.45,0)--(5.45,1.85);
    \draw[BookWarm,decorate,decoration={brace,amplitude=4pt}] (5.05,1.65)--(5.85,1.65)
      node[midway,above=5pt,font=\small] {卷积窗口};
    \node[font=\small,text=BookWarning] at (2.05,1.9) {粗糙函数};
    \node[font=\small,text=BookAccent] at (3.75,.35) {平滑后};
    \node[font=\small,text=BookWarm,anchor=west] at (5.55,1.15) {域的边界};
  \end{tikzpicture}
  \caption{核在内部滑过粗糙函数时产生光滑平均；若竖线代表域的边界，靠近它的窗口会越界。}
  \label{fig:4-3-1-mollifier-window}
\end{figure}

近似问题的核心是平移连续性。下面先证明所需的 $L^p$ 版本。

\begin{lemma}[$L^p$ 中的平移连续性]\label{lemma:4-3-1-translation}
若 $1\leq p<\infty$、$f\in L^p(\R^n)$，并记 $\tau_yf(x)=f(x-y)$，则
\[
 \|\tau_yf-f\|_p\longrightarrow0\qquad(y\to0).
\]
\end{lemma}

\begin{proof}
若 $f\in C_c(\R^n)$，取一个有限测度紧集 $K$，使 $|y|<1$ 时
$\supp(\tau_yf-f)\subset K$。紧支撑连续函数一致连续，故
\[
 \|\tau_yf-f\|_p^p
 \leq |K|\,\|\tau_yf-f\|_\infty^p\longrightarrow0.
\]
对一般 $f\in L^p$，由定理~\ref{theorem:3-3-3-1} 取 $g\in C_c(\R^n)$ 使
$\|f-g\|_p<\delta$。Lebesgue 测度平移不变，因而
\[
 \|\tau_yf-f\|_p
 \leq\|\tau_y(f-g)\|_p+\|\tau_yg-g\|_p+\|g-f\|_p
 <2\delta+\|\tau_yg-g\|_p.
\]
先令 $y\to0$，再令 $\delta\downarrow0$，结论成立。
\end{proof}

\begin{lemma}[积分平均的 Minkowski 估计]\label{lemma:4-3-1-integral-minkowski}
设 $1\leq p\leq\infty$，$F:\R^n\times\R^m\to\C$ 联合可测，并且
$y\mapsto\|F(\,\cdot,y)\|_{L^p(\R^n)}$ 可积。则
\[
 \left\|\int_{\R^m}F(\,\cdot,y)\dd y\right\|_{L^p(\R^n)}
 \leq\int_{\R^m}\|F(\,\cdot,y)\|_{L^p(\R^n)}\dd y.
 \tag{4.3.2a}\label{eq:4-3-integral-minkowski}
\]
左端的积分对 a.e. $x$ 绝对存在。
\end{lemma}

\begin{proof}
当 $p=1$ 时，逐点三角不等式与 Tonelli 定理给出
\[
 \int\left|\int F(x,y)\dd y\right|\dd x
 \leq\iint|F(x,y)|\dd y\dd x
 =\int\|F(\,\cdot,y)\|_1\dd y.
\]
当 $p=\infty$ 时，令
$a(y)=\|F(\,\cdot,y)\|_\infty$。可测集
\[
 N=\{(x,y):|F(x,y)|>a(y)\}
\]
的每个 $y$--截面都是零测集，故 Tonelli 与 Fubini 定理说明：对几乎每个 $x$，
$|F(x,y)|\leq a(y)$ 对几乎每个 $y$ 成立。因此原积分对这样的 $x$ 绝对存在，并且
\[
 \left|\int F(x,y)\dd y\right|
 \leq\int a(y)\dd y.
\]
取本性上确界即得所需估计。这里的 Fubini 步骤给出了一个与 $y$ 无关的 $x$--零集。

若 $1<p<\infty$，先把 $F$ 截断为有界且支撑于有限测度矩形的函数。以 $p'$ 表示共轭指数，
$L^p$ 对偶表示与 Fubini 定理给出
\begin{align*}
 \left\|\int F(\,\cdot,y)\dd y\right\|_p
 &=\sup_{\|g\|_{p'}\leq1}
   \left|\iint F(x,y)\overline{g(x)}\dd x\dd y\right|\\
 &\leq\int\sup_{\|g\|_{p'}\leq1}
       \left|\int F(x,y)\overline{g(x)}\dd x\right|\dd y
 \leq\int\|F(\,\cdot,y)\|_p\dd y.
\end{align*}
对一般 $F$，令
\[
 G_N(x,y)=|F(x,y)|\1_{B_{\R^n}(0,N)\times B_{\R^m}(0,N)}(x,y)
             \1_{\{|F(x,y)|\leq N\}}.
\]
则 $G_N\uparrow |F|$，而上述估计与
$\|G_N(\,\cdot,y)\|_p\leq\|F(\,\cdot,y)\|_p$ 给出
\[
 \left\|\int G_N(\,\cdot,y)\dd y\right\|_p
 \leq\int\|F(\,\cdot,y)\|_p\dd y.
\]
由单调收敛与 Fatou 引理，$x\mapsto\int|F(x,y)|\dd y$ 属于 $L^p$，并满足同一范数界。
所以原积分 a.e. 绝对存在；再用
$|\int F(x,y)\dd y|\leq\int|F(x,y)|\dd y$ 即得结论。
\end{proof}

\begin{theorem}[近似恒等]\label{theorem:4-3-1-2}
若 $1\leq p<\infty$ 且 $f\in L^p(\R^n)$，则
\[
 \|f*\rho_\varepsilon-f\|_p\longrightarrow0.
 \tag{4.3.2}\label{eq:4-3-mollifier-lp}
\]
若 $f\in C_0(\R^n)$，则收敛在一致范数中成立。若
$f\in L^p_{\mathrm{loc}}(\Omega)$，则对每个 $K\Subset\Omega$，
\[
 \|f*\rho_\varepsilon-f\|_{L^p(K)}\longrightarrow0
\]
（$\varepsilon$ 充分小时左端在 $K$ 上有定义）。
\end{theorem}

\begin{proof}
先设 $f\in L^p(\R^n)$。由 $\int\rho_\varepsilon=1$，逐点有
\[
 f*\rho_\varepsilon-f
 =\int\rho_\varepsilon(y)(\tau_yf-f)\dd y.
\]
引理~\ref{lemma:4-3-1-integral-minkowski} 给出
\[
 \|f*\rho_\varepsilon-f\|_p
 \leq\int\rho_\varepsilon(y)\|\tau_yf-f\|_p\dd y
 \leq\sup_{|y|\leq\varepsilon}\|\tau_yf-f\|_p.
\]
最后一项由引理~\ref{lemma:4-3-1-translation} 趋于零。

若 $f\in C_0$，则 $f$ 在 $\R^n$ 上一致连续。确实，可先在一个大闭球上用紧致性，再在球外利用
$f(x)\to0$。因此同一恒等式给出
\[
 \|f*\rho_\varepsilon-f\|_\infty
 \leq\sup_{|y|\leq\varepsilon}\|\tau_yf-f\|_\infty\longrightarrow0.
\]

最后取 $K\Subset\Omega$，再取 $\chi\in C_c(\Omega)$，使 $\chi=1$ 于 $K$ 的某个开邻域。
函数 $g=\chi f$ 的零延拓属于 $L^p(\R^n)$。当 $\varepsilon$ 小于该邻域宽度时，
$f*\rho_\varepsilon=g*\rho_\varepsilon$ 于 $K$，故局部结论由全局结论得到。
\end{proof}

$p=\infty$ 不能补进定理。令 $H=\1_{(0,\infty)}$；对每个 $\varepsilon>0$，在
$0<x<\varepsilon$ 的一个正测度区间上，$H*\rho_\varepsilon(x)$ 严格介于 $0$ 与 $1$，并且
当 $x\downarrow0$ 时趋于 $1/2$。所以
$\|H*\rho_\varepsilon-H\|_\infty\geq1/2$。因此跳跃函数排除了 $L^\infty$ 范数收敛。

\begin{example}[边界泄漏]\label{ex:4-3-1-3}
把 $(0,1)$ 上的常数函数 $1$ 零延拓为 $\widetilde f=\1_{(0,1)}$。公式
\eqref{eq:4-3-step-mollification} 表明，当 $0<x<\varepsilon$ 时
\[
 (\widetilde f*\rho_\varepsilon)(x)
 =\int_{-\varepsilon}^{x}\rho_\varepsilon(y)\dd y<1,
\]
并且 $x\downarrow0$ 时趋于 $1/2$。零延拓的卷积在每个内部紧集上仍逼近 $1$，却不保持边界附近的
常值。若问题要求零边界条件，零延拓可能正合适；若要求保持非零边界数据，就必须另作延拓或只在内部卷积。
\end{example}

\begin{proposition}[Lebesgue 点收敛]\label{proposition:4-3-1-4}
若 $f\in L^1_{\mathrm{loc}}(\R^n)$，则在 $f$ 的每个 Lebesgue 点 $x$，
\[
 f*\rho_\varepsilon(x)\longrightarrow f(x).
\]
\end{proposition}

\begin{proof}
由 $0\leq\rho\leq\|\rho\|_\infty$ 及支撑条件，
\begin{align*}
 |f*\rho_\varepsilon(x)-f(x)|
 &\leq \|\rho\|_\infty\varepsilon^{-n}
       \int_{B(0,\varepsilon)}|f(x-y)-f(x)|\dd y\\
 &=|B(0,1)|\|\rho\|_\infty
   \frac1{|B(x,\varepsilon)|}\int_{B(x,\varepsilon)}|f(z)-f(x)|\dd z.
\end{align*}
Lebesgue 点的定义使最后一项趋于零。
\end{proof}

这里的尺度上界不能删去。令 $a_k=2^{-k}$、$r_k=2^{-3k-4}$，并在 $\R$ 上置
\[
 f=\1_{\bigcup_{k\geq1}(-a_k-r_k,-a_k+r_k)},\qquad f(0)=0.
\]
若 $2^{-(m+1)}<r\leq2^{-m}$，则
\[
 \frac1{2r}\int_{-r}^{r}|f(t)|\dd t
 \leq \frac{\sum_{k\geq m-1}2r_k}{2^{-(m+1)}}
 =\frac{\sum_{k\ge m-1}2^{-3k-3}}{2^{-(m+1)}}
 =\frac{16}{7}\,2^{-2m}\longrightarrow0;
\]
所以零是 $f$ 的 Lebesgue 点。另一方面，取非负 $\kappa_k\in C_c^\infty$，使
$\int\kappa_k=1$ 且
$\supp\kappa_k\subset(a_k-r_k,a_k+r_k)\subset(-2a_k,2a_k)$。于是
\[
 (f*\kappa_k)(0)=\int f(-y)\kappa_k(y)\dd y=1.
\]
这些核非负、质量一且支撑收缩到零，却始终集中在稀疏尖峰上。Lebesgue 点只控制无权球平均；标准
mollifier 的 $\varepsilon^{-n}$ 上界则把加权平均控制在相应球平均之下，从而排除这种偏心集中。

下面同时处理边界距离、紧支撑与光滑性。

\begin{theorem}[光滑紧支撑函数的稠密性]\label{theorem:4-3-1-3}
对任意开集 $\Omega\subset\R^n$ 与 $1\leq p<\infty$，
$C_c^\infty(\Omega)$ 在 $L^p(\Omega)$ 中稠密。
\end{theorem}

\begin{proof}
固定 $f\in L^p(\Omega)$ 与 $\delta>0$。有限 Radon 测度
$\nu(E)=\int_E|f|^p\dd x$ 的内正则性给出紧集 $K\Subset\Omega$，使
$\int_{\Omega\setminus K}|f|^p<\delta^p$。由欧氏空间的局部 Urysohn 构造，取
$\chi\in C_c(\Omega)$，$0\leq\chi\leq1$ 且 $\chi=1$ 于 $K$。令 $g=\chi f$ 并零延拓到
$\R^n$，则
\[
 \|g-f\|_{L^p(\Omega)}^p
 =\int_{\Omega\setminus K}|1-\chi|^p|f|^p\dd x<\delta^p.
\]
$S=\supp g\subset\supp\chi$ 是 $\Omega$ 中的紧集。取
$0<\varepsilon<\tfrac12\dist(S,\Omega^c)$。命题~\ref{proposition:4-3-1-1} 说明
$g*\rho_\varepsilon$ 光滑，而
\[
 \supp(g*\rho_\varepsilon)
 \subset S+\overline{B(0,\varepsilon)}\Subset\Omega.
\]
再由定理~\ref{theorem:4-3-1-2} 选取足够小的 $\varepsilon$，使
$\|g*\rho_\varepsilon-g\|_p<\delta$。于是
$g*\rho_\varepsilon\in C_c^\infty(\Omega)$，且它到 $f$ 的 $L^p$ 距离小于 $2\delta$。
\end{proof}

上述定理讨论 $L^p$ 范数中的稠密性。Sobolev 范数中的情形有所不同：稍后将证明
$C^\infty(\Omega)\cap W^{1,p}(\Omega)$ 在 $W^{1,p}(\Omega)$ 中稠密，而
$C_c^\infty(\Omega)$ 的 Sobolev 范数闭包定义为 $W_0^{1,p}(\Omega)$，它一般真小于
$W^{1,p}(\Omega)$。卷积本身并不消除边界条件。

在 \S4.4 中，热核将提供另一类光滑化：固定正时间时卷积产生经典空间导数，而时间趋于零时近似
恒等恢复初值。

\begin{exercise}
证明对每个多重指标 $\alpha$，
$\|\partial^\alpha\rho_\varepsilon\|_1
=\varepsilon^{-|\alpha|}\|\partial^\alpha\rho\|_1$。再据此估计
$\|\partial^\alpha(f*\rho_\varepsilon)\|_p$，其中 $f\in L^p$。
\end{exercise}

\begin{exercise}
设 $f\in C_0(\R^n)$。不引用本节定理的 $C_0$ 部分，直接从“在大球内一致连续、在大球外一致小”
证明 $\|f*\rho_\varepsilon-f\|_\infty\to0$。
\end{exercise}

\begin{exercise}
从函数 $\eta(t)=e^{-1/t}\1_{(0,\infty)}(t)$ 出发，逐阶验证定义
\ref{def:4-3-1-1} 中的 $\rho$ 属于 $C_c^\infty(\R^n)$，并核对归一化常数确实使
$\int\rho=1$。
\end{exercise}

\begin{exercise}
设 $f\in L^1_{\mathrm{loc}}(\R^n)$、$\psi\in C_c^\infty(\R^n)$。为差商找出一个固定的
$L^1$ 支配函数，完整证明
$\partial^\alpha(f*\psi)=f*(\partial^\alpha\psi)$。
\end{exercise}

\begin{exercise}
沿“先截断到紧集，再选择小于边界距离的卷积尺度”的路线，独立证明
$C_c^\infty(\Omega)$ 在 $L^p(\Omega)$ 中稠密，$1\leq p<\infty$。
\end{exercise}

平滑近似让导数从粗糙函数转移到光滑核上。若进一步把核换成任意光滑紧支撑的“探针”，我们就能
把这次转移本身定义为导数；这正是分布语言的起点。

\subsection{分布、弱导数与测试函数对偶}\label{subsec:4-3-2}
\index{分布}\index{测试函数}\index{弱导数}

令 $H=\1_{(0,\infty)}$。原点以外，$H$ 的经典导数为零；若把导数简单地取为零，跳跃信息便会
被遗漏。另一方面，对每个 $\varphi\in C_c^\infty(\R)$，分部积分的形式计算给出
\[
 -\int_\R H(x)\varphi'(x)\dd x
 =-\int_0^\infty\varphi'(x)\dd x=\varphi(0).
 \tag{4.3.3}\label{eq:4-3-heaviside-motivation}
\]
右端只依赖测试函数在跳点的值。这提示我们：与其强迫导数逐点存在，不如把对象定义成它对所有光滑
紧支撑观测的响应。

Schwartz 分布理论把这种“由全部测试响应刻画对象”的想法系统化。以下采用逐紧集的半范数估计表述
连续性，这足以定义弱导数、卷积与分布收敛。

\begin{definition}[测试函数与分布]\label{def:4-3-2-1}
记 $\mathcal D(\Omega)=C_c^\infty(\Omega)$。一个\emph{分布}是线性泛函
$T:\mathcal D(\Omega)\to\C$，并且对每个紧集 $K\Subset\Omega$，存在常数
$C_K>0$ 与非负整数 $m_K$，使
\[
 |\langle T,\varphi\rangle|
 \leq C_K\sum_{|\alpha|\leq m_K}\|\partial^\alpha\varphi\|_\infty,
 \qquad \supp\varphi\subset K.
 \tag{4.3.4}\label{eq:4-3-distribution-continuity}
\]
全体分布记为 $\mathcal D'(\Omega)$。
\end{definition}

固定 $K$ 时，测试函数列 $\varphi_k$ 收敛到零，意指它们的支撑都包含于 $K$，且每个阶数
$\alpha$ 都满足 $\|\partial^\alpha\varphi_k\|_\infty\to0$。估计
\eqref{eq:4-3-distribution-continuity} 立即推出
$\langle T,\varphi_k\rangle\to0$。

\begin{definition}[正则分布]\label{def:4-3-2-2}
若 $f\in L^1_{\mathrm{loc}}(\Omega)$，定义
\[
 \langle T_f,\varphi\rangle=\int_\Omega f(x)\varphi(x)\dd x.
\]
称 $T_f$ 为 $f$ 诱导的正则分布。
\end{definition}

这确实满足定义：若 $\supp\varphi\subset K$，则
\[
 |\langle T_f,\varphi\rangle|
 \leq\left(\int_K|f|\right)\|\varphi\|_\infty.
\]
更重要的是，这个嵌入不丢失 a.e. 信息。

\begin{lemma}[测试函数识别局部可积函数]\label{lemma:4-3-2-regular-injective}
若 $u\in L^1_{\mathrm{loc}}(\Omega)$ 且
$\int u\varphi=0$ 对每个 $\varphi\in\mathcal D(\Omega)$ 成立，则 $u=0$ a.e.
因而 $T_f=T_g$ 当且仅当 $f=g$ a.e.
\end{lemma}

\begin{proof}
固定 $K\Subset\Omega$。对每个 $0<\varepsilon<\dist(K,\Omega^c)$ 及 $x\in K$，函数
$y\mapsto\rho_\varepsilon(x-y)$ 属于 $\mathcal D(\Omega)$，故
\[
 (u*\rho_\varepsilon)(x)
 =\int u(y)\rho_\varepsilon(x-y)\dd y=0.
\]
在稍大的紧邻域上截断 $u$ 后，定理~\ref{theorem:4-3-1-2} 的局部 $L^1$ 结论给出当
$\varepsilon\downarrow0$ 时 $u*\rho_\varepsilon\to u$ 于 $L^1(K)$。每个充分小的 $\varepsilon$
对应的左端都恒为零，故 $u=0$ a.e. 于 $K$。用
$\Omega_{1/k}$ 耗尽 $\Omega$ 即得全局结论。
\end{proof}

\begin{definition}[分布导数与弱导数]\label{def:4-3-2-3}
对 $T\in\mathcal D'(\Omega)$，定义
\[
 \langle\partial_jT,\varphi\rangle
 =-\langle T,\partial_j\varphi\rangle.
 \tag{4.3.5}\label{eq:4-3-distribution-derivative}
\]
若 $f,g\in L^1_{\mathrm{loc}}(\Omega)$ 且 $\partial_jT_f=T_g$，则称 $g$ 是 $f$ 的第 $j$ 个
\emph{弱导数}，仍记作 $\partial_jf$。
\end{definition}

负号来自经典分部积分：紧支撑测试函数使边界项消失。该定义还允许继续求任意阶分布导数。

\begin{proposition}[分布导数永远存在]\label{proposition:4-3-2-1}
每个分布都可以任意次微分；对多重指标 $\alpha,\beta$，
\[
 \partial^\alpha\partial^\beta T=\partial^{\alpha+\beta}T.
\]
特别地，分布的混合偏导与求导次序无关。
\end{proposition}

\begin{proof}
若 $T$ 在紧集 $K$ 上满足阶数 $m$ 的估计，则
\[
 |\langle\partial_jT,\varphi\rangle|
 =|\langle T,\partial_j\varphi\rangle|
 \leq C_K\sum_{|\alpha|\leq m}\|\partial^{\alpha+e_j}\varphi\|_\infty,
\]
所以 $\partial_jT$ 仍是分布。反复应用定义可作任意次微分。又因测试函数的经典偏导交换，
\[
 \langle\partial_i\partial_jT,\varphi\rangle
 =\langle T,\partial_j\partial_i\varphi\rangle
 =\langle T,\partial_i\partial_j\varphi\rangle
 =\langle\partial_j\partial_iT,\varphi\rangle.
\]
邻位交换给出任意求导次序的结论。
\end{proof}

\begin{example}[Heaviside 函数]\label{ex:4-3-2-1}
由 \eqref{eq:4-3-heaviside-motivation}，
\[
 \langle\partial T_H,\varphi\rangle=\varphi(0).
\]
定义 $\langle\delta_0,\varphi\rangle=\varphi(0)$，便得到
$\partial T_H=\delta_0$。点值泛函满足
$|\varphi(0)|\leq\|\varphi\|_\infty$，所以它的确是分布。它不是正则分布：若
$T_g=\delta_0$，则在每个不含零的开区间上测试可得 $g=0$ a.e.，因而 $g=0$ a.e. 于
$\R$；但取 $\varphi(0)=1$ 又会给 $\langle T_g,\varphi\rangle=1$，矛盾。因此 $H$ 没有
局部可积的弱导数。
\end{example}

\begin{example}[绝对值]\label{ex:4-3-2-2}
在 $(-\infty,0)$ 与 $(0,\infty)$ 分别积分，得到
\begin{align*}
 -\int_\R|x|\varphi'(x)\dd x
 &=-\int_{-\infty}^0\varphi(x)\dd x+
   \int_0^\infty\varphi(x)\dd x
 =\int_\R \sgn(x)\varphi(x)\dd x.
\end{align*}
因此 $(|x|)'=\sgn x$ 是弱导数。再次求分布导数，
\[
 -\int_\R\sgn(x)\varphi'(x)\dd x=2\varphi(0),
\]
故 $(|x|)''=2\delta_0$。
\end{example}

\begin{theorem}[弱导数唯一性]\label{theorem:4-3-2-2}
若 $g,h\in L^1_{\mathrm{loc}}(\Omega)$ 都是 $f$ 的第 $j$ 个弱导数，则 $g=h$ a.e.
更一般地，正则分布的代表在 a.e. 意义下唯一。
\end{theorem}

\begin{proof}
由定义，$T_g=\partial_jT_f=T_h$，故
$\int(g-h)\varphi=0$ 对每个测试函数成立。应用引理
\ref{lemma:4-3-2-regular-injective} 即得 $g=h$ a.e.
\end{proof}

分布也可以用测试核平滑，但必须先说明平移后的核何时仍留在论域内。对
$x\in\Omega_\varepsilon^\circ:=\{x:\dist(x,\Omega^c)>\varepsilon\}$，定义
\[
 (T*\rho_\varepsilon)(x)
 =\langle T,\rho_\varepsilon(x-\,\cdot)\rangle.
 \tag{4.3.6}\label{eq:4-3-distribution-convolution}
\]

\begin{proposition}[卷积与分布导数]\label{proposition:4-3-2-3}
函数 \eqref{eq:4-3-distribution-convolution} 在 $\Omega_\varepsilon^\circ$ 上属于
$C^\infty$，并且
\[
 \partial^\alpha(T*\rho_\varepsilon)
 = (\partial^\alpha T)*\rho_\varepsilon.
 \tag{4.3.7}\label{eq:4-3-distribution-convolution-derivative}
\]
\end{proposition}

\begin{proof}
固定 $x_0\in\Omega_\varepsilon^\circ$。当 $x$ 在 $x_0$ 的充分小邻域中变化时，所有
$y\mapsto\rho_\varepsilon(x-y)$ 的支撑都落在某个固定 $K\Subset\Omega$。对每个 $m$，经典
Taylor 公式说明差商
\[
 \frac{\rho_\varepsilon(x+he_j-\,\cdot)-
       \rho_\varepsilon(x-\,\cdot)}h
\]
连同其 $m$ 阶以内的各偏导，在 $K$ 上一致趋于
$\partial_j\rho_\varepsilon(x-\,\cdot)$。分布连续性估计遂允许在配对外取极限，得到
\[
 \partial_j(T*\rho_\varepsilon)(x)
 =\langle T,\partial_{x_j}\rho_\varepsilon(x-\,\cdot)\rangle.
\]
反复使用同一论证给出光滑性。另一方面，
\begin{align*}
 ((\partial_jT)*\rho_\varepsilon)(x)
 &=\langle\partial_jT,\rho_\varepsilon(x-\,\cdot)\rangle\\
 &=-\langle T,\partial_{y_j}\rho_\varepsilon(x-y)\rangle
 =\langle T,\partial_{x_j}\rho_\varepsilon(x-y)\rangle.
\end{align*}
两个负号在最后一步抵消。归纳得到 \eqref{eq:4-3-distribution-convolution-derivative}。
\end{proof}

由分布卷积的定义，例~\ref{ex:4-3-1-2} 中的恒等式现在可写为
\[
 (\delta_0*\rho_\varepsilon)(x)
 =\langle\delta_0,\rho_\varepsilon(x-\,\cdot)\rangle
 =\rho_\varepsilon(x).
\]

\begin{remark}[基本解的卷积形式]\label{ex:4-3-2-3}
若能找到分布 $\Phi$ 满足 $\Delta\Phi=\delta_0$，那么形式上
$u=\Phi*f$ 会满足 $\Delta u=f$。这条思路把微分方程的逆运算改写成卷积。要使该计算严格成立，
还需分别确定卷积的可定义性、基本解的具体形式与边界条件；这些问题将在 Fourier 与热核理论中处理。
\end{remark}

若 $T_k,T\in\mathcal D'(\Omega)$，称 $T_k$ 在\emph{分布意义下收敛}于 $T$，并写作
$T_k\to T$ in $\mathcal D'(\Omega)$，如果对每个 $\varphi\in\mathcal D(\Omega)$ 都有
\[
 \langle T_k,\varphi\rangle\longrightarrow
 \langle T,\varphi\rangle.
\]
这个定义只要求对每个测试函数的配对收敛，并不蕴含正则分布代表的逐点收敛或 $L^p$ 范数收敛。

\begin{theorem}[分布收敛与弱导数闭性]\label{theorem:4-3-2-4}
若 $f_k\to f$ 于 $L^1_{\mathrm{loc}}(\Omega)$，则对每个
$\varphi\in\mathcal D(\Omega)$，
\[
 \langle T_{f_k},\varphi\rangle\longrightarrow\langle T_f,\varphi\rangle.
\]
若对每个 $j=1,\ldots,n$ 另有 $g_{k,j}\to g_j$ 于 $L^1_{\mathrm{loc}}$，且
$\partial_jT_{f_k}=T_{g_{k,j}}$，则对每个 $j$ 都有 $\partial_jT_f=T_{g_j}$。

此外，设 $1<p<\infty$，$T_{f_k}\to T_f$ 于分布意义，并仍有
$\partial_jT_{f_k}=T_{g_{k,j}}$ 对每个 $j=1,\ldots,n$ 成立。若每个 $j$ 都满足
$\sup_k\|g_{k,j}\|_{L^p(\Omega)}<\infty$，则存在一个共同子列及
$g_j\in L^p(\Omega)$（$j=1,\ldots,n$），使对每个 $j$ 与每个
$\psi\in L^{p'}(\Omega)$，
\[
 \int g_{k_\ell,j}\psi\longrightarrow\int g_j\psi,
\]
并且 $\partial_jT_f=T_{g_j}$。
\end{theorem}

\begin{proof}
若 $K=\supp\varphi$，则
\[
 \left|\int(f_k-f)\varphi\right|
 \leq\|\varphi\|_\infty\|f_k-f\|_{L^1(K)}\longrightarrow0.
\]
同样，$\partial_j\varphi$ 有界且紧支撑，所以在
\[
 \int g_{k,j}\varphi=-\int f_k\partial_j\varphi
\]
中两边分别取极限，便得
$\int g_j\varphi=-\int f\partial_j\varphi$。

证明加强部分。$L^{p'}(\Omega)$ 可分，取可数稠密集
$\{\psi_m:m\geq1\}$。对每个 $j,m$，\Holder{} 不等式使数列
$\int g_{k,j}\psi_m$ 有界。由于 $j$ 只有有限多个，逐次抽取并对 $(j,m)$ 取对角子列，可令这些
数列同时收敛。以下固定 $j$。在
$\operatorname{span}\{\psi_m\}$ 上定义极限泛函 $L$；若
$M=\sup_k\|g_{k,j}\|_p$，则
\[
 |L(\psi)|\leq M\|\psi\|_{p'}.
\]
所以 $L$ 唯一延拓为 $L^{p'}$ 上的连续线性泛函。由
$L^{p'}$ 对偶表示定理~\ref{theorem:3-3-3-3}，存在 $h_j\in L^p$，使
\[
 L(\psi)=\int\psi\,\overline{h_j}.
\]
置 $g_j=\overline{h_j}$，便有 $L(\psi)=\int g_j\psi$ 且
$\|g_j\|_p=\|h_j\|_p\leq M$。给定一般
$\psi\in L^{p'}$，先以某个稠密族元素逼近 $\psi$，再用上面的统一界控制两端误差，得到
\[
 \int g_{k_\ell,j}\psi\to\int g_j\psi.
\]
最后取 $\psi=\varphi\in\mathcal D(\Omega)$，并在弱导数等式中令 $\ell\to\infty$；分布收敛处理
右端，刚得的标量积分收敛处理左端，于是 $\partial_jT_f=T_{g_j}$。因 $j$ 任意，结论对全部坐标
同时成立。
\end{proof}

\begin{example}[$p=1$ 端点的测度极限]\label{ex:4-3-2-4}
在 $(-1,1)$ 上定义
\[
 u_k(x)=
 \begin{cases}
 0,&x\leq0,\\
 kx,&0<x<1/k,\\
 1,&x\geq1/k.
 \end{cases}
\]
则 $\|u_k-H\|_1=1/(2k)\to0$，而
$u_k'=k\1_{(0,1/k)}$ 且 $\|u_k'\|_1=1$。对测试函数 $\varphi$，
\[
 \int k\1_{(0,1/k)}(x)\varphi(x)\dd x
 =\int_0^1\varphi(t/k)\dd t\longrightarrow\varphi(0).
\]
因此导数在分布意义下趋于 $\delta_0$，而不是某个 $L^1$ 函数。这里
$u_k\to H$ 只发生在 $L^1$ 中，绝不是在 $W^{1,1}$ 中；仅有 $L^1$ 范数有界不能代替上一
定理所需的标量积分子列紧性。导数质量集中时，$BV$ 正是容纳极限的更大空间。
\end{example}

\begin{figure}[htbp]
  \centering
  \begin{tikzpicture}[node distance=18mm and 20mm]
    \node[book warning node] (T) {粗糙对象 $T$};
    \node[book node,left=of T] (phi) {测试函数 $\varphi$};
    \node[book strong node,right=of T] (num) {数值 $\langle T,\varphi\rangle$};
    \node[book warm node,below=of phi] (dphi) {$\partial_j\varphi$};
    \node[book strong node,below=of T] (dT) {$\partial_jT$};
    \draw[book arrow] (phi)--(T);
    \draw[book arrow] (T)--(num);
    \draw[book arrow] (phi)--node[left,font=\small] {$\partial_j$}(dphi);
    \draw[book accent arrow] (T)--node[right,font=\small] {定义}(dT);
    \draw[book arrow] (dphi)--(dT);
    \node[font=\small,text=BookWarm,below=3mm of dT]
      {$\langle\partial_jT,\varphi\rangle=-\langle T,\partial_j\varphi\rangle$};
  \end{tikzpicture}
  \caption{分布只接受光滑测试；求导把微分箭头移到测试函数并增加负号。}
  \label{fig:4-3-2-distribution-black-box}
\end{figure}

分布与函数的乘法也有明确边界。若 $a\in C^\infty(\Omega)$，可以定义
$\langle aT,\varphi\rangle=\langle T,a\varphi\rangle$。但 $H\delta_0$ 没有由现有等价关系决定的
天然含义：Heaviside 的两个代表只在零点取值不同，诱导同一个正则分布；形式写成
$H(0)\delta_0$ 却会因这个不可观测点值而改变。准确的结论不是“分布从不能相乘”，而是一般分布之间
没有与 a.e. 等价及上述运算同时相容的规范乘法。

论域也必须写清。若把 $(0,1)$ 上的常数函数 $1$ 只看作该开区间内的正则分布，其分布导数为零；
若先零延拓
到 $\R$，所得 $\1_{(0,1)}$ 的分布导数却是 $\delta_0-\delta_1$。域内求导不会看见域外跳跃，
零延拓则会把边界跳跃变成边界测度。后面 $W_0^{1,p}$ 的闭包定义正是控制这种延拓行为的一种方式。

分布导数在弱解中把微分方程改写为对所有测试函数成立的积分恒等式。它还自然产生稠密定义的无界
微分算子；第~7 章建立 Fourier 变换后，分布微分将对应频率变量中的乘法。

\begin{exercise}
直接由定义证明 $\partial^\alpha\delta_0$ 满足
$\langle\partial^\alpha\delta_0,\varphi\rangle
=(-1)^{|\alpha|}\partial^\alpha\varphi(0)$，并验证它满足分布连续性估计。
\end{exercise}

\begin{exercise}
设 $f_k\to f$ 于 $L^1_{\mathrm{loc}}$。证明对每个固定 $\varepsilon>0$，
$f_k*\rho_\varepsilon\to f*\rho_\varepsilon$ 在卷积可定义的每个内部紧集上一致收敛。
\end{exercise}

\begin{exercise}
直接计算 $|x|$ 的一阶与二阶分布导数，并逐项核对第二阶结果中的系数与符号。
\end{exercise}

\begin{exercise}
只用定义~\ref{def:4-3-2-3} 与测试函数经典混合偏导的交换性，证明任意分布的混合偏导交换。
\end{exercise}

\begin{exercise}
设 $f_k\to f$ 于 $L^1_{\mathrm{loc}}(\Omega)$。对固定测试函数给出完整估计，证明
$T_{f_k}\to T_f$ 于分布意义。
\end{exercise}

\begin{exercise}
从 $-\int_0^\infty\varphi'$ 出发重新证明 $H'=\delta_0$，并说明为何这个导数不能由
$L^1_{\mathrm{loc}}$ 函数表示。
\end{exercise}

若一个函数的分布导数仍由 $L^p$ 函数表示，便可以同时测量函数和导数的大小。下一小节研究由此得到的
一阶 Sobolev 空间；分布导数的闭性将直接给出其完备性。

\subsection{一阶 Sobolev 空间、\Poincare{} 与紧性}\label{subsec:4-3-3}
\index{Sobolev 空间}\index{Poincare@\Poincare{} 不等式}\index{Rellich--Kondrachov 定理}

经典可微性把要求放在每一点；能量估计通常只能给出
$u,\partial_1u,\ldots,\partial_n u$ 的积分控制。绝对值函数与阶跃函数说明这不是同一种粗糙：前者
只有一个折角，分布导数仍是有界函数；后者有跳跃，导数已变成 Dirac 测度。

Sobolev 空间把“函数与弱导数同时受积分控制”组织成完备空间，\Poincare{} 不等式负责消除加法常数，
Rellich--Kondrachov 定理则在次临界范围把能量界转成紧性。三者解决的是不同问题；把它们并列在
同一小节，正是为了看清完备、范数控制与紧性之间不能互相替代。

\begin{example}[绝对值函数]\label{ex:4-3-3-1}
由例~\ref{ex:4-3-2-2}，$u(x)=|x|$ 在 $(-1,1)$ 上的弱导数为 $\sgn x$。因此对每个
$1\leq p<\infty$，$u$ 与 $u'$ 都属于 $L^p(-1,1)$，尽管 $u$ 在零点没有经典导数。
\end{example}

\begin{example}[跳跃不在 $W^{1,p}$]\label{ex:4-3-3-2}
Heaviside 函数在 $(-1,1)$ 上的分布导数是 $\delta_0$。引理
\ref{lemma:4-3-2-regular-injective} 说明它不可能同时是某个 $L^p$ 函数诱导的分布。因此跳跃函数
不属于下面定义的任何 $W^{1,p}(-1,1)$。这给出 $BV$ 允许跳跃而一阶 Sobolev 空间排除跳跃的
基本区别。
\end{example}

本小节除专门指出的端点外，始终取 $1\leq p<\infty$。

\begin{definition}[Sobolev 空间]\label{def:4-3-3-1}
设 $\Omega\subset\R^n$ 开。$W^{1,p}(\Omega)$ 由所有满足下列条件的 a.e. 等价类组成：
$u\in L^p(\Omega)$，且每个分布偏导 $\partial_ju$ 都由某个 $L^p(\Omega)$ 函数表示。赋予范数
\[
 \|u\|_{W^{1,p}(\Omega)}
 =\|u\|_{L^p(\Omega)}+
  \sum_{j=1}^n\|\partial_ju\|_{L^p(\Omega)}.
 \tag{4.3.8}\label{eq:4-3-sobolev-norm}
\]
并记 $\nabla u=(\partial_1u,\ldots,\partial_n u)$。
\end{definition}

\begin{definition}[$W_0^{1,p}$]\label{def:4-3-3-2}
$W_0^{1,p}(\Omega)$ 是 $C_c^\infty(\Omega)$ 在范数
\eqref{eq:4-3-sobolev-norm} 下的闭包。
\end{definition}

这个定义不需要预先建立边界迹，并以内部光滑逼近表达齐次边界条件。

\begin{definition}[平均零规范化]\label{def:4-3-3-3}
若 $0<|\Omega|<\infty$ 且 $u\in L^1(\Omega)$，记
\[
 u_\Omega=\frac1{|\Omega|}\int_\Omega u(x)\dd x.
\]
\end{definition}

常数的梯度为零，所以任何只用 $\nabla u$ 控制 $u$ 的不等式都必须先固定常数自由度；减去
$u_\Omega$ 是最自然的一种规范化。

\begin{theorem}[$W^{1,p}$ 完备]\label{theorem:4-3-3-1}
对每个开集 $\Omega\subset\R^n$ 与 $1\leq p<\infty$，
$W^{1,p}(\Omega)$ 是 Banach 空间。
\end{theorem}

\begin{proof}
设 $(u_k)$ 在 $W^{1,p}$ 中 Cauchy。$L^p$ 的完备性给出
$u,g_1,\ldots,g_n\in L^p(\Omega)$，使
\[
 u_k\to u,\qquad \partial_ju_k\to g_j
 \quad\text{于 }L^p(\Omega).
\]
每个紧集有限测度，\Holder{} 不等式说明这些收敛也发生在 $L^1_{\mathrm{loc}}$ 中。对
$\varphi\in C_c^\infty(\Omega)$，
\[
 \int g_j\varphi
 =\lim_{k\to\infty}\int(\partial_ju_k)\varphi
 =-\lim_{k\to\infty}\int u_k\partial_j\varphi
 =-\int u\partial_j\varphi.
\]
故 $g_j=\partial_ju$ 是弱导数，$u\in W^{1,p}$。各分量的 $L^p$ 收敛再给
$\|u_k-u\|_{W^{1,p}}\to0$。
\end{proof}

一维中，弱导数会恢复一个真正的绝对连续代表。端点值属于这个代表，而不属于原来任意选择的
$L^1$ 代表。

这里还需固定复值情形的含义。称 $F:[a,b]\to\C$ 绝对连续，是指 $\Re F$ 与 $\Im F$
都绝对连续；等价地，可在定义~\ref{def:4-2-3-1} 中把绝对值理解为复数模。后一种表述推出
前一种，因为 $|\Delta\Re F|,|\Delta\Im F|\le|\Delta F|$；反过来用
$|z|\le|\Re z|+|\Im z|$，分别把允许误差取为 $\varepsilon/2$ 即得。因而
\cref{theorem:4-2-3-1,theorem:4-2-3-2} 对复值函数逐分量成立。

\begin{proposition}[一维 Sobolev 的绝对连续代表]\label{proposition:4-3-3-ac-representative}
若 $u\in W^{1,1}(a,b)$，则存在唯一的绝对连续函数 $\widetilde u:[a,b]\to\C$，使
$\widetilde u=u$ a.e.，并且
\[
 \widetilde u(x)=c+\int_a^xu'(t)\dd t.
 \tag{4.3.9}\label{eq:4-3-ac-representative}
\]
\end{proposition}

\begin{proof}
令 $F(x)=\int_a^xu'(t)\dd t$。对实部与虚部分别应用 Lebesgue 基本定理
\ref{theorem:4-2-3-1}，$F$ 绝对连续且 $F'=u'$ a.e.。于是
$v=u-F$ 的分布导数为零。

固定闭区间 $I\Subset(a,b)$。对充分小的 $\varepsilon$，命题
\ref{proposition:4-3-2-3} 给
$(v*\rho_\varepsilon)'=v'*\rho_\varepsilon=0$ 于 $I$，所以
$v*\rho_\varepsilon$ 在 $I$ 上是常数。又由局部近似恒等，
$v*\rho_\varepsilon\to v$ 于 $L^1(I)$。常数函数空间在 $L^1(I)$ 中闭，故 $v$ a.e. 等于
$I$ 上的一个常数。两个相交内区间上的常数必须一致；以可数个相交内区间耗尽 $(a,b)$，得到
$v=c$ a.e.。于是 $c+F$ 是所求绝对连续代表。

若两个连续函数 a.e. 相等而在某点不同，则其差的绝对值在该点附近有正下界，因而不同的点集有
正测度，矛盾。故代表唯一。
\end{proof}

\begin{example}[一维代表]\label{ex:4-3-3-3}
若 $u\in W^{1,1}(a,b)$，公式 \eqref{eq:4-3-ac-representative} 给
\[
 \widetilde u(y)-\widetilde u(x)=\int_x^yu'(t)\dd t.
\]
特别地，$\widetilde u(a)=c$ 是绝对连续代表的端点值。随意改变原 $L^1$ 代表在 $a$ 点的数值并不
改变 Sobolev 元素，也不改变这个由内部数据连续延拓得到的 $c$。
\end{example}

在高维中，一维结论仍沿几乎每条坐标线成立。下面记录之后反射延拓所需的精确形式。

\begin{lemma}[坐标线上的绝对连续性]\label{lemma:4-3-3-acl}
设 $Q=I_1\times\cdots\times I_n\Subset\Omega$，$u\in W^{1,p}(\Omega)$。固定 $j$。
对几乎每个横截坐标 $\widehat x_j\in\prod_{i\ne j}I_i$，截面
$t\mapsto u(x_1,\ldots,t,\ldots,x_n)$ 有绝对连续代表，其一维导数 a.e. 等于
$\partial_ju$ 的相应截面。
\end{lemma}

\begin{proof}
先固定 $u$ 与各个 $\partial_ju$ 的 Lebesgue 可测代表；由前述完备化代表结论，也可把它们取成
Borel 代表。于是下面出现在积空间上的函数都具有真正的乘积可测代表，Fubini 截面不是对
a.e. 等价类作未定义的逐点取值。
取 $\alpha\in C_c^\infty(\prod_{i\ne j}I_i)$ 与
$\beta\in C_c^\infty(I_j)$。弱导数定义与 Fubini 定理给出
\[
 \int\alpha(\widehat x_j)
 \left[\int u(\widehat x_j,t)\beta'(t)\dd t
       +\int\partial_ju(\widehat x_j,t)\beta(t)\dd t\right]
 \dd\widehat x_j=0.
\]
固定 $\beta$。由于上式对每个 $\alpha$ 成立，局部可积函数识别引理说明方括号对几乎每个
$\widehat x_j$ 为零。为取得与 $\beta$ 无关的零集，取紧区间
$J_m\Subset I_j$ 递增耗尽 $I_j$；对每个 $m$，在所有支撑包含于 $J_m$ 的光滑函数中取一个关于
$C^1(J_m)$ 范数的可数稠密子集，再对 $m$ 取并，得到可数族 $\mathcal B$。同时由 Fubini 定理，
删去一个零集后，$u(\widehat x_j,\cdot)$ 与
$\partial_ju(\widehat x_j,\cdot)$ 都属于 $L^1(I_j)$。再删去 $\mathcal B$ 中各测试函数对应零集的
可数并。

现在固定保留下来的 $\widehat x_j$ 与任意 $\beta\in C_c^\infty(I_j)$。取某个 $J_m$，使其内部包含
$\supp\beta$，并选 $\beta_\ell\in\mathcal B$，使
$\beta_\ell\to\beta$、$\beta_\ell'\to\beta'$ 在 $J_m$ 上一致。于是
\begin{align*}
 &\left|\int u(\beta_\ell'-\beta')\right|
 +\left|\int\partial_ju(\beta_\ell-\beta)\right|\\
 &\qquad\leq
 \|u(\widehat x_j,\cdot)\|_{L^1(I_j)}
       \|\beta_\ell'-\beta'\|_\infty
 +\|\partial_ju(\widehat x_j,\cdot)\|_{L^1(I_j)}
       \|\beta_\ell-\beta\|_\infty
 \longrightarrow0.
\end{align*}
因此截面恒等式对所有 $\beta$ 同时成立，该截面的分布导数就是
$\partial_ju$ 的相应截面。命题
\ref{proposition:4-3-3-ac-representative} 给出所述绝对连续代表。
\end{proof}

下一步要在开域内同时平滑函数与全部弱导数。卷积尺度不能统一选择：域可能无界，边界距离也可能沿
不同位置趋于零。局部有限单位分解正是为每一块保留独立尺度。

\begin{lemma}[光滑单位分解的局部构造]\label{lemma:4-3-3-smooth-partition}
每个开集 $\Omega\subset\R^n$ 有一族局部有限的
$\theta_k\in C_c^\infty(\Omega)$，满足 $0\leq\theta_k\leq1$ 与
$\sum_{k\geq1}\theta_k=1$。还可同时取局部有限的相对紧开集 $O_k$，使
$\supp\theta_k\subset O_k\Subset\Omega$。若预先给定一个可数开覆盖 $(V_j)$，还可使每个 $k$ 都有某个
$j(k)$ 满足 $\supp\theta_k\Subset V_{j(k)}$；同一个 $V_j$ 可以对应多个指标 $k$。
\end{lemma}

\begin{proof}
取 $U_k=\Omega_{1/k}$（$k\geq1$）。由定义~\ref{def:4-3-1-3}，
$\overline U_k\subset U_{k+1}$ 且 $\bigcup_kU_k=\Omega$；并约定
$U_j=\varnothing$ 当 $j\leq0$。壳层
$A_k=\overline U_k\setminus U_{k-2}$ 是紧集，并可由有限多个闭球覆盖，而这些球的二倍闭球仍包含在
$U_{k+1}\setminus\overline U_{k-3}$ 内；若给定了覆盖 $(V_j)$，还可逐球要求其二倍闭球包含在
某个 $V_j$ 内。对其中一个闭球 $K=\overline{B(c,r)}$，其二倍闭球已包含在上述指定开集内。令
$G=B(c,3r/2)$，并取 $0<\delta<r/2$，置
\[
 \phi_K=\1_G*\rho_\delta.
\]
于是 $0\leq\phi_K\leq1$。若 $x\in K$ 且 $y\in\supp\rho_\delta$，则
$|x-y-c|\leq r+|y|<3r/2$，所以
\[
 \phi_K(x)=\int\rho_\delta(y)\dd y=1.
\]
另一方面
\[
 \supp\phi_K\subset\overline G+\overline{B(0,\delta)}
 \subset B(c,2r),
\]
故其支撑紧含于预先指定的开集；卷积求导公式给
$\phi_K\in C_c^\infty$。
保留每个球对应的函数并把所有壳层中的函数逐一排列成 $(\phi_\ell)$；这样在给定覆盖时，每个
$\phi_\ell$ 的支撑仍紧含于选定的某个 $V_{j(\ell)}$。每个壳层只有有限项，而第 $k$ 个壳层所生
函数的支撑连同某个稍大的相对紧开邻域都位于
$U_{k+1}\setminus\overline U_{k-3}$，所以这些邻域组成局部有限族；把它们依次记为 $O_\ell$。
具体地，若某个点的一个邻域包含于 $U_m$，则该邻域不再与 $k\geq m+3$ 的上述壳层邻域相交，
而较小的 $k$ 只有有限多个。另一方面，每个点按其首次落入某个 $U_k$ 的指标属于相应 $A_k$，
故闭球确实覆盖整个 $\Omega$。因此每个 $x\in\Omega$ 至少落在一个相应闭球上，对该指标有
$\phi_\ell(x)=1$，从而局部有限和
\[
 S(x)=\sum_\ell\phi_\ell(x)\geq1.
\]
于是 $S$ 是严格正的光滑函数。置 $\theta_\ell=\phi_\ell/S$，则
$0\leq\theta_\ell\leq1$ 且
$\sum_\ell\theta_\ell=S/S=1$。由于每一点附近只有有限项，求和、除法与求导均为通常的有限运算。
\end{proof}

\begin{proposition}[Meyers--Serrin 光滑稠密]\label{proposition:4-3-3-1}
对任意开集 $\Omega\subset\R^n$ 与 $1\leq p<\infty$，
$C^\infty(\Omega)\cap W^{1,p}(\Omega)$ 在 $W^{1,p}(\Omega)$ 中稠密。
\end{proposition}

\begin{proof}
取引理~\ref{lemma:4-3-3-smooth-partition} 的单位分解及局部有限开邻域 $(O_k)$，并记
$K_k=\supp\theta_k\Subset O_k$。先验证乘积法则。
若 $v_k=\theta_ku$，则对测试函数 $\varphi$，
\begin{align*}
 \int v_k\partial_j\varphi
 &=\int u\partial_j(\theta_k\varphi)-\int u(\partial_j\theta_k)\varphi\\
 &=-\int\bigl(\theta_k\partial_ju+u\partial_j\theta_k\bigr)\varphi.
\end{align*}
故
\[
 \partial_jv_k=\theta_k\partial_ju+u\partial_j\theta_k,\qquad
 v_k\in W^{1,p}(\Omega),\quad \supp v_k\subset K_k.
 \tag{4.3.10}\label{eq:4-3-sobolev-product}
\]
把 $v_k$ 零延拓到 $\R^n$。为核对零延拓没有产生边界测度，取
$\chi_k\in C_c^\infty(\Omega)$ 在 $K_k$ 的一个邻域上等于一。对
$\Phi\in C_c^\infty(\R^n)$，弱导数等式应用于 $\chi_k\Phi$，并利用
$v_k$ 及其弱导数都支撑于 $K_k$，得到
\[
 \int_{\R^n}\widetilde v_k\,\partial_j\Phi
 =-\int_{\R^n}\widetilde{\partial_jv_k}\,\Phi.
\]
故零延拓仍属于 $W^{1,p}(\R^n)$，其弱导数是原弱导数的零延拓。由命题
\ref{proposition:4-3-2-3} 或直接测试计算，
\[
 \partial_j(v_k*\rho_\delta)=(\partial_jv_k)*\rho_\delta.
\]
定理~\ref{theorem:4-3-1-2} 同时应用于 $v_k$ 与各个 $\partial_jv_k$，给出
$v_k*\rho_\delta\to v_k$ 于 $W^{1,p}(\R^n)$。

给定 $\varepsilon>0$，逐块选择 $\delta_k>0$，使
\[
 \delta_k<\tfrac12\dist(K_k,\Omega^c),\qquad
 \delta_k<\dist(K_k,O_k^c),\qquad
 \|v_k*\rho_{\delta_k}-v_k\|_{W^{1,p}(\R^n)}
 <2^{-k}\varepsilon,
\]
于是膨胀支撑 $K_k+\overline{B(0,\delta_k)}$ 包含于 $O_k$，因而仍局部有限。记
\[
 e_k=v_k*\rho_{\delta_k}-v_k.
\]
上式给出
\[
 \sum_{k=1}^\infty\|e_k\|_{W^{1,p}(\Omega)}
 \leq\sum_{k=1}^\infty\|e_k\|_{W^{1,p}(\R^n)}
 <\varepsilon.
\]
$W^{1,p}(\Omega)$ 的完备性因而说明 $\sum_ke_k$ 在该空间中收敛。另一方面，支撑的局部有限性说明
这个空间极限作为分布恰等于逐点局部有限和。于是
\[
 w=\sum_{k=1}^\infty v_k*\rho_{\delta_k}
\]
在每个点附近只是有限和，故 $w\in C^\infty(\Omega)$。又因
$u=\sum_kv_k$ a.e.，所以在分布意义以及 $W^{1,p}$ 中
$w-u=\sum_ke_k$。因此
\[
 \|w-u\|_{W^{1,p}(\Omega)}
 \leq\sum_{k=1}^\infty
 \|v_k*\rho_{\delta_k}-v_k\|_{W^{1,p}(\R^n)}<\varepsilon.
\]
这证明稠密性。
\end{proof}

命题中的近似函数不要求紧支撑。若要求近似函数具有紧支撑，其闭包恰是定义
\ref{def:4-3-3-2} 中的 $W_0^{1,p}$；例如有界区间上的常数一不能由
$C_c^\infty$ 在 $W^{1,p}$ 中逼近，因为稍后的 \Poincare{} 不等式会迫使这种逼近的函数范数也随
导数范数趋于零。

\begin{example}[长盒上的尺度依赖]\label{ex:4-3-3-long-box}
令 $\Omega_L=(0,L)\times(0,1)^{n-1}$，$u(x)=x_1-L/2$。对 $x_1$ 积分可知
$u_{\Omega_L}=0$，而 $|\nabla u|=1$。直接缩放 $x_1=Lt$ 得
\[
 \|u\|_p^p=L^{p+1}\int_0^1|t-1/2|^p\dd t,\qquad
 \|\nabla u\|_p^p=L.
\]
两者之比等于只依赖 $p$ 的常数乘以 $L$。因此 \Poincare{} 常数必须依赖于区域的尺度，不能仅依赖
$n,p$。
\end{example}

先在区间上看常数自由度如何被删掉。由绝对连续代表，若 $I=(a,b)$、$\ell=b-a$，则
\begin{align*}
 |u(x)-u_I|
 &\leq\frac1\ell\int_I|u(x)-u(y)|\dd y
 \leq\int_I|u'(t)|\dd t.
\end{align*}
因此
\begin{equation}\label{eq:4-3-interval-poincare}
 \|u-u_I\|_{L^p(I)}
 \leq \ell^{1/p}\|u'\|_{L^1(I)}
 \leq \ell\|u'\|_{L^p(I)}.
\end{equation}
高维凸域仍可沿线段积分，但若直接换元，会出现看似不可积的 $t^{-n}$。把哪些线段经过同一点重新
计数后，奇性恰好合并成一个可积势核。

\begin{lemma}[凸域中的势核估计]\label{lemma:4-3-3-convex-potential}
设 $\Omega\subset\R^n$ 是有界凸开集，$D=\diam\Omega$。若
$u\in C^1(\Omega)$ 且 $u,\nabla u\in L^p(\Omega)$，则对 a.e. $x\in\Omega$，
\[
 |u(x)-u_\Omega|
 \leq \frac{D^n}{n|\Omega|}
 \int_\Omega\frac{|\nabla u(z)|}{|x-z|^{n-1}}\dd z.
 \tag{4.3.11}\label{eq:4-3-convex-potential}
\]
当 $n=1$ 时把核 $|x-z|^{1-n}$ 读作 $1$。
\end{lemma}

\begin{proof}
由凸性，连接 $x,y\in\Omega$ 的线段留在 $\Omega$，故
\[
 |u(x)-u(y)|
 \leq |x-y|\int_0^1|\nabla u(x+t(y-x))|\dd t.
\]
对 $y$ 平均，并固定 $x,t$ 作换元 $z=x+t(y-x)$。此时
$\dd y=t^{-n}\dd z$、$|x-y|=t^{-1}|x-z|$，且
$x+t(\Omega-x)\subset\Omega$。反过来，若某个 $z$ 出现在该缩小域中，则
$|x-z|=t|x-y|\leq tD$，所以 $t\geq|x-z|/D$。Tonelli 定理于是给出
\begin{align*}
 |u(x)-u_\Omega|
 &\leq\frac1{|\Omega|}\int_0^1
   \int_{x+t(\Omega-x)}|x-z|t^{-n-1}|\nabla u(z)|\dd z\dd t\\
 &\leq\frac1{|\Omega|}\int_\Omega |x-z||\nabla u(z)|
   \int_{|x-z|/D}^1t^{-n-1}\dd t\dd z\\
 &=\frac1{n|\Omega|}\int_\Omega |x-z||\nabla u(z)|
   \left[\left(\frac{D}{|x-z|}\right)^n-1\right]\dd z\\
 &\leq\frac{D^n}{n|\Omega|}
   \int_\Omega\frac{|\nabla u(z)|}{|x-z|^{n-1}}\dd z.
\end{align*}
对角线是零测集，不影响积分。
\end{proof}

\begin{theorem}[\Poincare{} 不等式]\label{theorem:4-3-3-2}
设 $\Omega\subset\R^n$ 是有界凸域，$1\leq p<\infty$。存在常数
$C=C(n,p,\Omega)$，使每个 $u\in W^{1,p}(\Omega)$ 满足
\[
 \|u-u_\Omega\|_{L^p(\Omega)}
 \leq C\|\nabla u\|_{L^p(\Omega)}.
 \tag{4.3.12}\label{eq:4-3-poincare-mean}
\]
此外，存在 $C_0=C_0(n,p,\Omega)$，使每个 $u\in W_0^{1,p}(\Omega)$ 满足
\[
 \|u\|_{L^p(\Omega)}\leq C_0\|\nabla u\|_{L^p(\Omega)}.
 \tag{4.3.13}\label{eq:4-3-poincare-zero}
\]
\end{theorem}

\begin{proof}
先设 $u\in C^\infty(\Omega)\cap W^{1,p}(\Omega)$。把
$g=|\nabla u|\1_\Omega$ 零延拓到 $\R^n$，并令
$K(w)=|w|^{1-n}\1_{B(0,D)}(w)$。引理
\ref{lemma:4-3-3-convex-potential} 的右端不超过常数乘
$\int K(w)g(x-w)\dd w$。而
\[
 \|K\|_1
 =|S^{n-1}|\int_0^D r^{1-n}r^{n-1}\dd r
 =|S^{n-1}|D<\infty.
\]
对非负简单函数先用有限和的 Minkowski 不等式，再以单调收敛逼近，可得
\begin{align*}
 \left\|\int K(w)g(\,\cdot-w)\dd w\right\|_{L^p(\R^n)}
 &\leq\int K(w)\|g(\,\cdot-w)\|_p\dd w\\
 &=\|K\|_1\|g\|_p.
\end{align*}
于是 \eqref{eq:4-3-poincare-mean} 对光滑 $u$ 成立。命题
\ref{proposition:4-3-3-1} 给出 $u_k\to u$ 于 $W^{1,p}$；\Holder{} 不等式还给
$(u_k)_\Omega\to u_\Omega$，故令 $k\to\infty$ 即得一般情形。

再证 \eqref{eq:4-3-poincare-zero}。零边界版本不能由平均零版本直接删除 $u_\Omega$ 得到。取有界开盒
$Q$，使 $\overline\Omega\subset Q$；$Q$ 是凸域且 $A=Q\setminus\Omega$ 有正测度。若
$u\in C_c^\infty(\Omega)$，把它零延拓为 $\widetilde u\in C_c^\infty(Q)$。已证的平均版本给出
\[
 \|\widetilde u-(\widetilde u)_Q\|_{L^p(Q)}
 \leq C_Q\|\nabla\widetilde u\|_{L^p(Q)}.
\]
由于 $\widetilde u=0$ 于 $A$，
\[
 |(\widetilde u)_Q|\,|A|^{1/p}
 =\|\widetilde u-(\widetilde u)_Q\|_{L^p(A)}
 \leq C_Q\|\nabla\widetilde u\|_p.
\]
故
\[
 \|u\|_{L^p(\Omega)}
 \leq C_Q\left(1+\frac{|Q|^{1/p}}{|A|^{1/p}}\right)
       \|\nabla u\|_{L^p(\Omega)}.
\]
最后按定义以 $C_c^\infty(\Omega)$ 中的序列逼近任意 $u\in W_0^{1,p}$。零延拓后的函数和弱导数
在 $L^p(Q)$ 中分别收敛，因此极限正是 $u$ 的零延拓；令序列趋于极限便得结论。
\end{proof}

连通性不可省略。若 $\Omega$ 是两个不交开球之并，取函数在一个分量上为零、另一个分量上为一，
则弱梯度为零而 $u-u_\Omega$ 不为零。对一般有界连通 Lipschitz 域，\Poincare{} 仍成立，但需要
边界坐标图或链式覆盖；上面的核心定理只声称凸域版本。

\begin{figure}[htbp]
  \centering
  \begin{tikzpicture}[scale=1.0]
    \draw[book region] plot[smooth cycle,tension=.8] coordinates
      {(0,0) (1.1,-.65) (3.3,-.5) (4.7,.3) (4.1,1.8) (2.2,2.15) (.45,1.45)};
    \coordinate (x) at (.7,.45);
    \coordinate (y) at (4.0,1.45);
    \fill[BookInk] (x) circle (1.7pt) node[below left,font=\small] {$x$};
    \fill[BookInk] (y) circle (1.7pt) node[above right,font=\small] {$y$};
    \draw[BookAccent,very thick] (x)--(y);
    \fill[BookWarm] ($(x)!.58!(y)$) circle (1.7pt)
      node[below right,font=\small] {$z=x+t(y-x)$};
    \draw[BookWarm,decorate,decoration={brace,amplitude=4pt}]
      ($(x)+(0,-.25)$)--($(y)+(0,-.25)$)
      node[midway,below=5pt,font=\small] {$u(y)-u(x)=\int\nabla u\cdot\dd z$};
    \node[font=\small,text=BookAccent] at (2.45,2.5) {对所有 $y$ 平均，常数方向消失};
  \end{tikzpicture}
  \caption{凸性保证连接 $x,y$ 的整条线段留在域内；沿线积分并对 $y$ 平均产生可积势核。}
  \label{fig:4-3-3-poincare-segments}
\end{figure}

沿线段积分还有一个更直接的定量结果，它既控制 mollifier 误差，也将成为紧性准则的输入。

\begin{proposition}[平移估计]\label{proposition:4-3-3-3}
暂记 $\tau_hu(x)=u(x-h)$。若 $u\in W^{1,p}(\R^n)$、$1\leq p<\infty$，则
\[
 \|\tau_hu-u\|_{L^p(\R^n)}
 \leq |h|\,\|\,|\nabla u|\,\|_{L^p(\R^n)}.
 \tag{4.3.14}\label{eq:4-3-translation-estimate}
\]
\end{proposition}

\begin{proof}
若 $u\in C^\infty(\R^n)\cap W^{1,p}$，则
\[
 u(x-h)-u(x)=-\int_0^1h\cdot\nabla u(x-th)\dd t.
\]
积分型 Minkowski 不等式、Cauchy--Schwarz 与平移不变性给出
\[
 \|\tau_hu-u\|_p
 \leq |h|\int_0^1\|\nabla u(\,\cdot-th)\|_p\dd t
 =|h|\|\nabla u\|_p.
\]
对一般 $u$，用命题~\ref{proposition:4-3-3-1} 在 $\R^n$ 上取
$u_k\to u$ 于 $W^{1,p}$。平移是 $L^p$ 等距，故左端与右端都可令 $k\to\infty$。
\end{proof}

若 $p=1$ 且 $u\geq0$、$\int u=1$，那么 $u(x)\dd x$ 与
$u(x-h)\dd x$ 是两个概率测度。此时直接把总变差距离定义为密度差的一半，
\[
 d_{\mathrm{TV}}
 =\frac12\|\tau_hu-u\|_1
 \leq\frac{|h|}{2}\|\nabla u\|_1.
\]
若 $X$ 具有密度 $u$，则这两个测度分别是 $\mathcal L(X+h)$ 与 $\mathcal L(X)$。因此上式也给出
随机向量平移前后分布的总变差估计。

下面两项选读结果说明一阶能量还能提供更强的可积性或连续性。为建立它们，先给出 Lipschitz 边界上
所需的延拓机制。称有界开集 $\Omega$ 为 Lipschitz 域，如果每个
边界点附近经过刚性坐标变换后，都存在长方体 $U'\times(-a,a)$ 与 Lipschitz 函数
$\phi:U'\to\R$，使
\[
 \Omega\cap(U'\times(-a,a))
 =\{(x',x_n):x_n>\phi(x')\}
\]
（缩小长方体后理解）。边界紧致性允许只取有限张图。

连续性结论还需要固定所用的函数空间。若 $K\subset\R^n$ 紧、$0<\alpha\le1$，定义
\begin{equation*}
 [u]_{C^{0,\alpha}(K)}
 :=\sup_{\substack{x,y\in K\\x\ne y}}
   \frac{|u(x)-u(y)|}{|x-y|^\alpha},
 \qquad
 \|u\|_{C^{0,\alpha}(K)}
 :=\|u\|_\infty+[u]_{C^{0,\alpha}(K)}.
\end{equation*}
记 $C^{0,\alpha}(K)$ 为上述范数有限的连续函数空间，并约定
$C^{0,0}(K)=C(K)$、$\|u\|_{C^{0,0}}=\|u\|_\infty$。例如在 $[0,1]$ 上，
$u(x)=x^\alpha$ 满足
$|x^\alpha-y^\alpha|\le|x-y|^\alpha$，故属于 $C^{0,\alpha}$；若
$\beta>\alpha$，则
$|u(x)-u(0)|/|x|^\beta=x^{\alpha-\beta}\to\infty$，所以一般不能把指数提高。

\begin{lemma}[一阶 Sobolev 的 Lipschitz 延拓]\label{lemma:4-3-3-lipschitz-extension}
若 $\Omega$ 是有界 Lipschitz 域且 $1\leq p<\infty$，则存在线性算子
\[
 E:W^{1,p}(\Omega)\longrightarrow W^{1,p}(\R^n)
\]
与固定球 $B(0,R)$，使 $Eu=u$ a.e. 于 $\Omega$、
$\supp Eu\subset\overline{B(0,R)}$，并且
\[
 \|Eu\|_{W^{1,p}(\R^n)}\leq C_{\Omega,p}\|u\|_{W^{1,p}(\Omega)}.
 \tag{4.3.15}\label{eq:4-3-extension-bound}
\]
\end{lemma}

\begin{proof}
先处理支撑紧含于一张边界图柱的一块函数；取两个嵌套图柱并乘一个在内柱上等于一的光滑截断，便可
设这一块在横向边界与柱的上下边界附近为零，只有图面 $s=0$ 可以与其支撑相交。最后的有限单位分解
会把一般函数化成这样的块。令
\[
 F(x',s)=(x',s+\phi(x')).
\]
对每个固定 $x'$，这是竖直方向的平移，所以 Fubini 定理给
$\int |v\circ F|^p=\int|v|^p$（在对应长方体中）。Lipschitz 函数 $\phi$ 限制在每条坐标线上
仍为 Lipschitz，因而由一维微积分基本定理几乎处处可微，且相应导数的绝对值不超过 Lipschitz
常数 $L$。
对每个 $i<n$，Fubini 定理把坐标线上的零测例外集合并成 $U'$ 中的零测集；因此
$\partial_i\phi$ 几乎处处存在且 $|\partial_i\phi|\leq L$。若 $v$ 光滑，则对固定其余坐标，
映射
\[
 r\longmapsto v(x'+re_i,s+\phi(x'+re_i))
\]
是一个 $C^1$ 函数与一维绝对连续函数的复合，一维链式法则给出
\[
 \partial_i(v\circ F)=(\partial_iv)\circ F
   +(\partial_nv)\circ F\,\partial_i\phi\quad(i<n),
 \qquad
 \partial_n(v\circ F)=(\partial_nv)\circ F.
 \tag{4.3.16}\label{eq:4-3-flatten-chain}
\]
Fubini 定理把这些逐线恒等式变成弱导数恒等式。再利用固定 $x'$ 时
$s\mapsto s+\phi(x')$ 的 Jacobian 为一，得到
\[
 \|v\circ F\|_{W^{1,p}}\leq C(n,L)\|v\|_{W^{1,p}}.
\]
对一般 $v\in W^{1,p}$，先由 Meyers--Serrin 命题
\ref{proposition:4-3-3-1} 在图内以光滑函数逼近；上式说明复合列在平坦半长方体中 Cauchy，
而竖直变量替换还给出
$\|(v_m-v)\circ F\|_p=\|v_m-v\|_p$，所以该极限确是 $v\circ F$ 的 a.e. 等价类。
式~\eqref{eq:4-3-flatten-chain} 随后由弱导数闭性成立。

在平坦坐标中，把 $w(x',s)$ 从 $s>0$ 偶反射为
\[
 \mathcal Rw(x',s)=w(x',|s|).
\]
先看切向导数。把测试积分分成 $s>0$ 与 $s<0$ 两部分，并在后一部分作代换 $s\mapsto-s$；
随后对 $x_i$ 使用 $w$ 的弱导数定义，得到
\[
 \partial_i\mathcal Rw(x',s)=
 \begin{cases}\partial_iw(x',s),&s>0,\\
                \partial_iw(x',-s),&s<0,
 \end{cases}
\qquad i<n.
\]
再把引理~\ref{lemma:4-3-3-acl} 用在竖直直线上。一维绝对连续代表的偶延拓仍绝对连续，故
$\partial_n$ 在 $s<0$ 一侧等于 $-(\partial_nw)(x',-s)$，且界面上没有 Dirac 项。逐项换元给出
\[
 \|\mathcal Rw\|_{L^p}^p=2\|w\|_{L^p}^p,\qquad
 \|\partial_j\mathcal Rw\|_{L^p}^p
 =2\|\partial_jw\|_{L^p}^p\quad(1\leq j\leq n),
\]
于是 $\|\mathcal Rw\|_{W^{1,p}}=2^{1/p}\|w\|_{W^{1,p}}$。再以
\[
 F^{-1}(x',x_n)=(x',x_n-\phi(x'))
\]
送回原坐标。对 $F^{-1}$ 重复 \eqref{eq:4-3-flatten-chain} 的逐坐标计算（把
$\phi$ 换成 $-\phi$）得到反向估计
$\|z\circ F^{-1}\|_{W^{1,p}}\leq C(n,L)\|z\|_{W^{1,p}}$，故得到单张图上的有界反射延拓。

最后用有限张边界图 $U_1,\ldots,U_N$ 与一个相对紧的内部开集 $U_0\Subset\Omega$
覆盖 $\overline\Omega$。这里需要的是环境空间中的有限单位分解，而不是支撑紧含于 $\Omega$ 的
单位分解：先缩小为仍覆盖 $\overline\Omega$ 的开集 $V_i$，使
$\overline V_i\subset U_i$；再取 $\chi_i\in C_c^\infty(U_i)$，满足
$0\leq\chi_i\leq1$ 且 $\chi_i=1$ 于 $\overline V_i$。函数
$S=\sum_{i=0}^N\chi_i$ 在 $\overline\Omega$ 的某个开邻域 $W$ 上严格为正。取开集
$W_1$ 与函数 $\zeta\in C_c^\infty(W)$，使
\[
 \overline\Omega\subset W_1,\qquad \overline W_1\Subset W,
 \qquad \zeta=1\quad\hbox{于 }W_1.
\]
在 $W$ 上置 $\theta_i=\zeta\chi_i/S$，并在 $W$ 外置零。因为
$\supp\zeta\Subset W$，这些零延拓都是全空间的光滑紧支撑函数；在 $\Omega\subset W_1$ 上有
\[
 \sum_{i=0}^N\theta_i=\zeta\frac{\sum_i\chi_i}{S}=1.
\]
这个截断步骤也保证没有在 $\{S=0\}$ 的边界留下未定义的比值。

内部块 $\theta_0u$ 的支撑紧含于 $U_0\Subset\Omega$，故可零延拓。对 $i\geq1$，
$\supp\theta_i\Subset U_i$；缩小图柱后，上述平坦化与反射在包含
$\supp\theta_i\cap\overline\Omega$ 的邻域中给出 $\theta_i u$ 的局部延拓。再取
$\zeta_i\in C_c^\infty(U_i)$，使它在该支撑的一个邻域上等于一，并把局部延拓乘以
$\zeta_i$ 后零延拓到全空间。这样得到线性算子 $E_i$，满足
\[
 E_i(\theta_i u)=\theta_i u\quad\hbox{a.e. 于 }\Omega,
 \qquad
 \|E_i(\theta_i u)\|_{W^{1,p}(\R^n)}
 \leq C_i\|u\|_{W^{1,p}(\Omega)}.
\]
这里的估计依次使用乘积法则~\eqref{eq:4-3-sobolev-product}、平坦化的双向估计、反射估计和
再次乘截断的乘积法则。定义
\[
 Eu=\widetilde{\theta_0u}+\sum_{i=1}^NE_i(\theta_i u),
\]
其中第一项表示零延拓。有限三角不等式给出
\eqref{eq:4-3-extension-bound}；在 $\Omega$ 内各块之和为 $u$。所有截断支撑的有限并包含在某个
固定球中，故 $Eu$ 具有与 $u$ 无关的固定紧支撑。
\end{proof}

延拓把区域上的估计化为 $\R^n$ 上的估计。在整空间中，关键端点估计只需一维微积分、Fubini 与
广义 \Holder{} 不等式。

\begin{lemma}[欧氏 Sobolev 与 Morrey 核心估计]\label{lemma:4-3-3-euclidean-sobolev}
设 $u\in C_c^\infty(\R^n)$。
\begin{enumerate}
  \item 若 $n\geq2$，则
  \[
   \|u\|_{L^{n/(n-1)}}\leq C_n\|\nabla u\|_{L^1}.
   \tag{4.3.17}\label{eq:4-3-sobolev-l1}
  \]
  \item 若 $1<p<n$、$p^*=np/(n-p)$，则
  \[
   \|u\|_{L^{p^*}}\leq C_{n,p}\|\nabla u\|_{L^p}.
   \tag{4.3.18}\label{eq:4-3-sobolev-p}
  \]
  \item 若 $p>n$、$\alpha=1-n/p$，则
  \[
   |u(x)-u(y)|\leq C_{n,p}|x-y|^\alpha\|\nabla u\|_{L^p}
   \quad(x,y\in\R^n).
   \tag{4.3.19}\label{eq:4-3-morrey}
  \]
\end{enumerate}
\end{lemma}

\begin{proof}
对第一项，沿第 $i$ 条坐标轴积分，得到
\[
 |u(x)|\leq g_i(\widehat x_i),\qquad
 g_i(\widehat x_i)=\int_\R|\partial_i u(x_1,\ldots,t,\ldots,x_n)|\dd t.
\]
所以
\[
 |u(x)|^{n/(n-1)}
 \leq\prod_{i=1}^n g_i(\widehat x_i)^{1/(n-1)}.
\]
反复应用 Fubini 与普通 \Holder{} 给出
\[
 \int_{\R^n}\prod_{i=1}^ng_i(\widehat x_i)^{1/(n-1)}\dd x
 \leq\prod_{i=1}^n
 \left(\int_{\R^{n-1}}g_i\right)^{1/(n-1)}.
 \tag{4.3.20}\label{eq:4-3-coordinate-holder}
\]
具体归纳如下。$n=2$ 时两个因子分别只依赖另一个坐标，Fubini 直接给等号。设结论已对
$n-1$ 个变量成立。先固定 $x'=(x_1,\ldots,x_{n-1})$，对 $x_n$ 积分；因 $g_n$ 与 $x_n$
无关，对其余 $n-1$ 个因子使用指数 $n-1$ 的 \Holder{}，得到
\[
 g_n(x')^{1/(n-1)}
 \prod_{i=1}^{n-1}G_i(x'_{\widehat i})^{1/(n-1)},
 \qquad
 G_i(x'_{\widehat i})=\int_\R g_i(x'_{\widehat i},x_n)\dd x_n.
\]
这里 $G_i$ 是 $n-2$ 个变量的函数，不依赖 $x_i$。再在 $\R^{n-1}$ 上用 \Holder{}，指数分别取
$n-1$ 与 $(n-1)/(n-2)$，上式的积分不超过
\[
 \left(\int g_n\right)^{1/(n-1)}
 \left(\int_{\R^{n-1}}
       \prod_{i=1}^{n-1}G_i^{1/(n-2)}\right)^{(n-2)/(n-1)}.
\]
归纳假设把第二个因子控制为
$\prod_{i=1}^{n-1}(\int G_i)^{1/(n-1)}$；Fubini 又给
$\int G_i=\int g_i$。这就证明 \eqref{eq:4-3-coordinate-holder}。由于
$\int g_i=\|\partial_i u\|_1$，取相应次方并用几何平均不超过算术平均，得到
\eqref{eq:4-3-sobolev-l1}。

设 $1<p<n$，令 $\gamma=p(n-1)/(n-p)>1$。把第一项应用于
$v_\delta=(|u|^2+\delta^2)^{\gamma/2}-\delta^\gamma$。对实值或复值 $u$，逐坐标链式法则与
$|\partial_i|u||\le |\partial_i u|$（在 $u\ne0$ 处直接计算，在 $u=0$ 处取连续正则化）给出
\[
 |\nabla v_\delta|
 \le \gamma |u|(|u|^2+\delta^2)^{\gamma/2-1}|\nabla u|.
\]
当 $1<\gamma\le2$ 时，右侧不超过
$\gamma|u|^{\gamma-1}|\nabla u|$；当 $\gamma>2$ 且 $0<\delta\le1$ 时，它在 $u$ 的固定紧支撑上
被 $C_u|\nabla u|$ 支配。故令 $\delta\downarrow0$ 时，左侧用 Fatou 引理，右侧用支配收敛定理，
再用 \Holder{} 不等式，得到
\begin{align*}
 \|u\|_{\gamma n/(n-1)}^\gamma
 &\leq C_n\gamma\int|u|^{\gamma-1}|\nabla u|\\
 &\leq C_n\gamma\|\nabla u\|_p
       \|u\|_{(\gamma-1)p'}^{\gamma-1}.
\end{align*}
直接计算
$\gamma n/(n-1)=(\gamma-1)p'=p^*$。若 $u\not\equiv0$，约去
$\|u\|_{p^*}^{\gamma-1}$；若 $u=0$ 结论成立。于是得到
\eqref{eq:4-3-sobolev-p}。

最后设 $p>n$。先把定理~\ref{theorem:4-3-3-2} 的 $p=1$ 情形从单位球缩放到
$B=B(x,r)$：对 $v(z)=u(x+rz)$ 应用单位球不等式，再作变量替换，得到
\[
 \frac1{|B|}\int_B|u-u_B|
 \leq C_nr\frac1{|B|}\int_B|\nabla u|
 \leq C_{n,p}r^{1-n/p}\|\nabla u\|_{L^p(B)}.
 \tag{4.3.21}\label{eq:4-3-ball-average}
\]
记 $B_k=B(x,2^{-k}r)$。由于 $|B_k|/|B_{k+1}|=2^n$，
\eqref{eq:4-3-ball-average} 给出
\[
 \begin{aligned}
 |u_{B_{k+1}}-u_{B_k}|
 &\le \frac1{|B_{k+1}|}\int_{B_{k+1}}|u-u_{B_k}|\\
 &\le 2^n\frac1{|B_k|}\int_{B_k}|u-u_{B_k}|
 \le C_{n,p}(2^{-k}r)^\alpha\|\nabla u\|_{L^p(B_k)}.
 \end{aligned}
\]
对 $k\geq0$ 求和；几何级数收敛，而光滑性给
$u_{B_k}\to u(x)$，故
\[
 |u(x)-u_{B(x,r)}|\leq C r^\alpha\|\nabla u\|_p.
\]
若 $r=|x-y|>0$，令
$B_x=B(x,2r)$、$B_y=B(y,2r)$ 与 $B_*=B(x,4r)$；于是
$B_x\cup B_y\subset B_*$，且 $|B_*|/|B_x|=|B_*|/|B_y|=2^n$。上式分别控制
$|u(x)-u_{B_x}|$ 与 $|u(y)-u_{B_y}|$。此外，对 $C=B_x$ 或 $B_y$，
\[
 |u_C-u_{B_*}|
 \le\frac1{|C|}\int_C|u-u_{B_*}|
 \le2^n\frac1{|B_*|}\int_{B_*}|u-u_{B_*}|
 \le C_{n,p}r^\alpha\|\nabla u\|_p.
\]
把这四个差用三角不等式相加，得到
$|u(x)-u(y)|\leq C_{n,p}r^\alpha\|\nabla u\|_p$，即
\eqref{eq:4-3-morrey}；$x=y$ 时结论成立。
\end{proof}

\begin{proposition}[Sobolev 连续嵌入（选读）]\label{proposition:4-3-3-4}
设 $\Omega$ 为有界 Lipschitz 域，$1\leq p<\infty$。
\begin{enumerate}
  \item 若 $p<n$，令 $p^*=np/(n-p)$，则
  $W^{1,p}(\Omega)\hookrightarrow L^q(\Omega)$ 对每个 $1\leq q\leq p^*$ 连续。
  \item 若 $p=n\geq2$，则 $W^{1,n}(\Omega)\hookrightarrow L^q(\Omega)$ 对每个有限
  $q$ 连续，但一般不嵌入 $L^\infty$。一维的 $W^{1,1}(a,b)$ 则嵌入 $L^\infty$。
  \item 若 $p>n$，则每个 Sobolev 元素有唯一的
  $C^{0,1-n/p}(\overline\Omega)$ 代表，并且相应嵌入连续。
\end{enumerate}
\end{proposition}

\begin{proof}
令 $E$ 为引理~\ref{lemma:4-3-3-lipschitz-extension} 的延拓算子。若 $1<p<n$，以
Meyers--Serrin 逼近 $Eu$，再用 \eqref{eq:4-3-sobolev-p} 及
\eqref{eq:4-3-extension-bound} 取极限，得到
$\|u\|_{p^*}\leq C\|u\|_{W^{1,p}}$。若 $p=1<n$，因 $Eu$ 有固定紧支撑，用 mollifier 对
$Eu$ 及其弱导数同时近似，得到 $C_c^\infty(\R^n)$ 中的 $W^{1,1}$ 逼近列。在每一项上应用
\eqref{eq:4-3-sobolev-l1}，再由 $W^{1,1}$ 收敛和 Fatou 引理得到
$\|u\|_{n/(n-1)}\le C\|u\|_{W^{1,1}}$。对 $1\leq q<p^*$，有限测度 \Holder{} 不等式给
$\|u\|_q\leq|\Omega|^{1/q-1/p^*}\|u\|_{p^*}$。

若 $p=n\geq2$ 且给定有限 $q$，取 $1\leq p_0<n$ 使 $q\leq p_0^*$。由于 $Eu$ 支撑在固定
有限测度球中，\Holder{} 不等式给
$\|Eu\|_{W^{1,p_0}}\leq C\|Eu\|_{W^{1,n}}$；应用刚证的 $p_0$ 情形即可。一维中，
绝对连续代表满足
\[
 |\widetilde u(x)|
 \leq |u_I|+\int_I|u'|
 \leq |I|^{-1}\|u\|_1+\|u'\|_1,
\]
故嵌入 $L^\infty$。

若 $p>n$，因 $Eu$ 有紧支撑，以 mollifier 卷积可取
$v_k\in C_c^\infty(\R^n)$ 使 $v_k\to Eu$ 于 $W^{1,p}$。对 $v_k$ 应用
\eqref{eq:4-3-morrey}，右端由 $\|u\|_{W^{1,p}}$ 一致控制。取包含所有支撑的固定球 $B$。对
$w=v_k-v_\ell$ 与每个 $x\in B$，\eqref{eq:4-3-morrey} 给出
\[
 \begin{aligned}
 |w(x)|
 &\leq |w_B|+\frac1{|B|}\int_B|w(x)-w(y)|\dd y\\
 &\leq |B|^{-1/p}\|w\|_{L^p(B)}
   +C_{n,p}(\diam B)^\alpha\|\nabla w\|_{L^p(\R^n)}.
 \end{aligned}
\]
故 $(v_k)$ 在一致范数中 Cauchy，设其一致极限为 $v$；$L^p$ 极限的唯一性说明
$v=Eu$ a.e.。再令 $k\to\infty$ 于
\[
 |v_k(x)-v_k(y)|\leq C_{n,p}|x-y|^\alpha\|\nabla v_k\|_p
\]
并使用 \eqref{eq:4-3-extension-bound}，得到
$[v]_{C^{0,\alpha}}\leq C\|u\|_{W^{1,p}(\Omega)}$。同时对 $w=v_k$ 的上一估计给出
$\|v\|_\infty\leq C_{\Omega,p}\|u\|_{W^{1,p}(\Omega)}$。限制到
$\overline\Omega$ 即得完整的 $C^{0,1-n/p}$ 范数估计。若两个连续代表在 $\Omega$ 中 a.e. 相等而
某个内点处之差非零，连续性会给出一个差仍非零的开球，与 a.e. 相等矛盾；故二者先在 $\Omega$
处处相等，再由 $\Omega$ 在 $\overline\Omega$ 中稠密而在边界上也相等。因此代表唯一。
\end{proof}

$p=n$ 的 $L^\infty$ 端点确实错误。设 $n\geq2$，在 $B(0,e^{-2})$ 上令
\[
 u(x)=\log\log\frac e{|x|},
\]
并在外侧作光滑截断。它在零点附近无界，而
\[
 |\nabla u(x)|=\frac1{|x|\log(e/|x|)},
\qquad
 \int_{|x|<e^{-2}}|\nabla u|^n\dd x
 =|S^{n-1}|\int_0^{e^{-2}}\frac{\dd r}{r\log^n(e/r)}<\infty.
\]
$u$ 本身也属于 $L^n$：令 $s=\log(e/r)$，则 $r=e^{1-s}$，从而
\[
 \int_{|x|<e^{-2}}|u(x)|^n\dd x
 =|S^{n-1}|e^n\int_3^\infty(\log s)^n e^{-ns}\dd s<\infty.
\]
为验证原点没有产生额外分布项，令
$u_M=\min\{u,M\}$，其中外侧截断预先取成非负径向光滑函数。对充分大的 $M$，集合
$\{u\geq M\}$ 是以原点为心的小球；$u_M$ 在该球内恒为 $M$，在球外分片 $C^1$，并在界面两侧
取同一个值。对测试函数分别在内球与外环作经典分部积分，界面上的两个法向边界项因迹相同而抵消，
故
\[
 \nabla u_M=\1_{\{u<M\}}\nabla u\qquad\text{a.e.}
\]
确为其弱梯度。又由支配收敛定理，
$u_M\to u$、$\nabla u_M\to\nabla u$ 于 $L^n$。弱导数闭性遂说明
$u\in W^{1,n}$，但它没有有界代表。

连续嵌入只给统一界；从有界列中抽出强收敛子列还需要排除振荡与质量逃逸。所需的整空间紧性准则
可以用小立方体平均直接证明。

\begin{lemma}[Kolmogorov--Riesz 的本章特例]\label{lemma:4-3-3-kolmogorov-riesz}
设 $1\leq p<\infty$，$\mathcal F\subset L^p(\R^n)$ 满足
\begin{align*}
 &\sup_{f\in\mathcal F}\|f\|_p<\infty,\tag{4.3.22}\label{eq:4-3-kr-bound}\\
 &\lim_{h\to0}\sup_{f\in\mathcal F}\|\tau_hf-f\|_p=0,\tag{4.3.23}\label{eq:4-3-kr-translation}\\
 &\lim_{R\to\infty}\sup_{f\in\mathcal F}
   \|f\1_{\{|x|>R\}}\|_p=0.\tag{4.3.24}\label{eq:4-3-kr-tail}
\end{align*}
则 $\mathcal F$ 在 $L^p(\R^n)$ 中相对紧。
\end{lemma}

\begin{proof}
我们证明 $\mathcal F$ 全有界；$L^p$ 完备性随后给出相对紧性。给定 $\varepsilon>0$，先由
\eqref{eq:4-3-kr-tail} 选 $R$，使所有函数在立方体
$Q=[-R,R]^n$ 外的 $L^p$ 范数小于 $\varepsilon$。再选充分大的整数 $M$，令
$\ell=2R/M$ 足够小，并把 $Q$ 精确分成 $M^n$ 个边长为 $\ell$ 的半开小立方体
$Q_1,\ldots,Q_N$。定义（在 $Q$ 外取零）
\[
 P_\ell f=\sum_{j=1}^N f_{Q_j}\1_{Q_j},\qquad
 f_{Q_j}=\frac1{|Q_j|}\int_{Q_j}f.
\]
Jensen 不等式给
\begin{align*}
 \|f-P_\ell f\|_{L^p(Q)}^p
 &\leq\sum_j\frac1{|Q_j|}
   \int_{Q_j}\int_{Q_j}|f(x)-f(y)|^p\dd y\dd x\\
 &=\ell^{-n}\int_{[-\ell,\ell]^n}
   \sum_j\int_{Q_j\cap(Q_j+h)}
   |f(x)-f(x-h)|^p\dd x\dd h\\
 &\leq\ell^{-n}\int_{[-\ell,\ell]^n}
   \|\tau_hf-f\|_{L^p(\R^n)}^p\dd h\\
 &\leq2^n\sup_{|h|\leq\sqrt n\ell}
   \|\tau_hf-f\|_{L^p(\R^n)}^p.
 \tag{4.3.25}\label{eq:4-3-cube-average-approx}
\end{align*}
第二行使用 $h=x-y$；对固定 $h$，各集合 $Q_j\cap(Q_j+h)$ 两两不交，因而第三行可把积分域扩大到
$\R^n$。由 \eqref{eq:4-3-kr-translation} 取 $M$ 足够大，使最后一行小于
$\varepsilon^p$。

$P_\ell\mathcal F$ 位于由 $\1_{Q_1},\ldots,\1_{Q_N}$ 张成的有限维空间。再次用 Jensen
不等式，
\[
 \|P_\ell f\|_p^p
 =\sum_{j=1}^N |Q_j|\,|f_{Q_j}|^p
 \leq\sum_{j=1}^N\int_{Q_j}|f|^p
 \leq\|f\|_p^p,
\]
所以它由 \eqref{eq:4-3-kr-bound} 一致有界；有限维有界集全有界。因此存在有限个 $L^p$ 球覆盖
$P_\ell\mathcal F$。每个 $f$ 到 $P_\ell f$ 的距离至多为内部误差加尾误差，小于
$2\varepsilon$。放大有限个球半径便覆盖 $\mathcal F$，所以 $\mathcal F$ 全有界。
\end{proof}

\begin{theorem}[Rellich--Kondrachov（选读）]\label{theorem:4-3-3-4}
设 $\Omega\subset\R^n$ 为有界 Lipschitz 域。
\begin{enumerate}
  \item 若 $1\leq p<n$，则嵌入
  $W^{1,p}(\Omega)\hookrightarrow L^q(\Omega)$ 对每个
  $1\leq q<p^*=np/(n-p)$ 都是紧的；临界 $q=p^*$ 一般不紧。
  \item 若 $p=n$，则对每个有限 $q$，嵌入
  $W^{1,n}(\Omega)\hookrightarrow L^q(\Omega)$ 是紧的。
  \item 若 $p>n$，则对每个 $0\leq\beta<1-n/p$，嵌入
  $W^{1,p}(\Omega)\hookrightarrow C^{0,\beta}(\overline\Omega)$ 是紧的；
  $\beta=0$ 表示 $C^0$。因此它也紧嵌入每个 $L^q$，$1\leq q\leq\infty$。
\end{enumerate}
\end{theorem}

\begin{proof}
设 $(u_k)$ 在 $W^{1,p}(\Omega)$ 中有界，并令 $v_k=Eu_k$，其中 $E$ 来自引理
\ref{lemma:4-3-3-lipschitz-extension}。于是 $(v_k)$ 在 $W^{1,p}(\R^n)$ 中有界，并且支撑都包含在
同一个球内。命题~\ref{proposition:4-3-3-3} 给出统一平移模
\[
 \sup_k\|\tau_hv_k-v_k\|_p
 \leq |h|\sup_k\|\nabla v_k\|_p\longrightarrow0.
\]
固定支撑使尾条件 \eqref{eq:4-3-kr-tail} 自动成立。引理
\ref{lemma:4-3-3-kolmogorov-riesz} 因而给出一个在 $L^p(\R^n)$ 中收敛的子列，限制回
$\Omega$ 后仍在 $L^p(\Omega)$ 中收敛。

若 $p<n$，命题~\ref{proposition:4-3-3-4} 给出 $(u_k)$ 在 $L^{p^*}$ 中一致有界。对
$p<q<p^*$，取 $0<\theta<1$ 使
$1/q=\theta/p+(1-\theta)/p^*$。\Holder{} 不等式给
\[
 \|u_k-u_\ell\|_q
 \leq\|u_k-u_\ell\|_p^\theta
       \|u_k-u_\ell\|_{p^*}^{1-\theta},
\]
所以该子列在 $L^q$ 中 Cauchy。若 $q\leq p$，有限测度 \Holder{} 不等式直接把 $L^p$ 收敛降到
$L^q$。这证明所有 $q<p^*$ 的紧性。

若 $p=n\geq2$ 且给定有限 $q$，选 $p_0<n$ 使 $q<p_0^*$。有界域上的 \Holder{} 不等式说明
$W^{1,n}$ 有界列也是 $W^{1,p_0}$ 有界列，应用第一部分即可。若 $n=p=1$，一维估计说明该列在
$L^\infty$ 中一致有界；上文已经从 Kolmogorov--Riesz 得到 $L^1$ 收敛子列。对每个有限
$q>1$，
\[
 \|u_k-u_\ell\|_q^q
 \leq\|u_k-u_\ell\|_\infty^{q-1}\|u_k-u_\ell\|_1,
\]
故同一子列在 $L^q$ 中收敛；$q=1$ 已经得到。

最后设 $p>n$，$\alpha=1-n/p$。命题~\ref{proposition:4-3-3-4} 说明选定的连续代表一致有界且
具有统一的 $\alpha$-\Holder{} 模。$\overline\Omega$ 紧，\ArzelaAscoli{} 定理给出一致收敛子列。
若 $0<\beta<\alpha$，对任意连续函数 $w$ 有插值估计
\[
 [w]_{C^{0,\beta}}
 \leq 2^{1-\beta/\alpha}
      \|w\|_\infty^{1-\beta/\alpha}
      [w]_{C^{0,\alpha}}^{\beta/\alpha}.
 \tag{4.3.26}\label{eq:4-3-holder-interpolation}
\]
事实上，对 $r=|x-y|$，分别使用
$|w(x)-w(y)|\leq2\|w\|_\infty$ 与
$|w(x)-w(y)|\leq[w]_\alpha r^\alpha$，再取两上界的加权几何平均即可。把
\eqref{eq:4-3-holder-interpolation} 用于子列两项之差，一致收敛与统一 $C^{0,\alpha}$ 界给出
$C^{0,\beta}$ 收敛。$\beta=0$ 已由 \ArzelaAscoli{} 得到。最后
$\|w\|_{L^q}\leq|\Omega|^{1/q}\|w\|_\infty$ 给出所有 $L^q$ 结论，包括 $q=\infty$。
\end{proof}

临界指数为何不紧，可以直接缩放验算。取 $x_0\in\Omega$ 和非零
$\eta\in C_c^\infty(B(0,1))$，并令
\[
 u_k(x)=k^{(n-p)/p}\eta(k(x-x_0))
\]
（$k$ 充分大时 $\supp u_k\Subset\Omega$）。换元给出
\[
 \|u_k\|_{p^*}=\|\eta\|_{p^*},\qquad
 \|\nabla u_k\|_p=\|\nabla\eta\|_p,\qquad
 \|u_k\|_p=k^{-1}\|\eta\|_p.
\]
所以 $(u_k)$ 在 $W^{1,p}$ 中有界，并且对 $x\ne x_0$ 最终恒为零。若某子列在
$L^{p^*}$ 中强收敛，则可再抽取 a.e. 收敛子列，其极限只能为零；强收敛却会迫使
$\|u_k\|_{p^*}\to0$，与上式矛盾。

尾条件同样不是多余的。在 $\R^n$ 上取非零
$\eta\in C_c^\infty$ 并令 $w_k(x)=\eta(x-ke_1)$。其 $W^{1,p}$ 范数与平移模完全相同，但支撑
彼此向无穷远移动。取间距足够大的子列可使支撑两两不交，于是
$\|w_k-w_\ell\|_p^p=2\|\eta\|_p^p$；它没有 $L^p$ Cauchy 子列。有界域延拓后的固定支撑
正是用来排除这种整体逃逸。

Sobolev 嵌入与 Rellich--Kondrachov 紧性都依赖于区域边界、指数范围和尾部控制；定理中的这些条件
不可省略。在 PDE 中，$W_0^{1,p}$ 表达齐次边界条件，\Poincare{} 不等式把梯度估计提升为函数范数
估计，而 Rellich 紧性在次临界范围内把能量有界列转化为强收敛列。平移估计
\eqref{eq:4-3-translation-estimate} 还将用于下一节的热核正则化。

\begin{exercise}
证明若 $u\in W^{1,p}(\Omega)$ 且 $a\in C^1(\Omega)$、$a$ 与 $\nabla a$ 有界，则
$au\in W^{1,p}(\Omega)$，并求其弱导数。说明有界性假设在 $a$ 具有紧支撑时如何局部化。
\end{exercise}

\begin{exercise}
在区间 $(a,b)$ 上，用绝对连续代表证明
$\|u-u_{(a,b)}\|_p\leq(b-a)\|u'\|_p$。构造一列函数说明常数必须随区间长度按一次尺度变化。
\end{exercise}

\begin{exercise}
完成 \eqref{eq:4-3-coordinate-holder} 的归纳证明，并据此从
\eqref{eq:4-3-sobolev-l1} 逐行推出 \eqref{eq:4-3-sobolev-p}。
\end{exercise}

\begin{exercise}
在引理~\ref{lemma:4-3-3-kolmogorov-riesz} 的证明中，详细写出从双积分到
\eqref{eq:4-3-cube-average-approx} 的 $h=x-y$ 换元，并说明为何选取
$\ell=2R/M$ 可以避免边缘立方体。
\end{exercise}

\begin{exercise}
从 $L^p$ 的完备性和弱导数闭性出发，独立证明 $W^{1,p}(\Omega)$ 完备；证明中须写出测试函数配对的
取极限步骤。
\end{exercise}

\begin{exercise}
验证 $u(x)=|x|$ 属于 $W^{1,p}(-1,1)$，$1\leq p<\infty$，并说明它为什么没有处处经典导数。
\end{exercise}

\begin{exercise}
先对 $C_c^\infty(\R^n)$ 沿线段积分，再用 Meyers--Serrin 逼近，完整证明平移估计
\eqref{eq:4-3-translation-estimate}。
\end{exercise}

卷积在这里首先是平滑工具；平移估计又揭示它是由平移平均出来的算子。下一节把卷积、近似恒等、
热核以及强连续半群放进同一套运算结构中。

\section{卷积、近似恒等与半群}\label{sec:04-04}

\subsection{平移连续、卷积与 Young 不等式}\label{subsec:4-4-1}
\index{平移连续}\index{卷积}\index{Young 不等式}

两个函数相乘时比较的是同一点的取值。卷积换了一个问题：把其中一个函数平移以后，两者在所有位置
上的重叠总量是多少？平移量改变，重叠量便形成一个新函数。这个操作既可看成移动平均，也可看成
积空间中的质量经加法映射汇合；后一种看法解释了为什么卷积天然带有结合律。

卷积同时出现于 Fourier 分析中的平移不变运算与测度的质量叠加。下面先从函数积分与测度推前定义
卷积；它与独立随机变量之和及 Fourier 乘法的关系将在相应章节说明。

先从一个完全显式的例子开始。
\begin{example}[两个区间的重叠长度]\label{ex:4-4-1-1}
令 $f=g=\1_{[0,1]}$。固定 $x$ 时，乘积
$f(x-y)g(y)$ 恰在 $y\in[0,1]\cap[x-1,x]$ 上等于一，因此
\begin{equation}\label{eq:indicator-convolution}
 (f*g)(x)=m\bigl([0,1]\cap[x-1,x]\bigr)
 =\begin{cases}
 x,&0\le x\le1,\\
 2-x,&1<x\le2,\\
 0,&\text{其他}.
 \end{cases}
\end{equation}
矩形的两条竖边经过卷积变成三角形的两条斜边。该计算同时显示总质量保持为一，而卷积的峰值可由
两个因子的适当范数控制。
\end{example}

\begin{definition}[平移]\label{def:4-4-1-1}
对 $h\in\R^n$，定义
\[
  \tau_hf(x)=f(x-h).
\]
Lebesgue 测度的平移不变性给出 $\|\tau_hf\|_p=\|f\|_p$，其中
$1\le p\le\infty$。
\end{definition}

\begin{definition}[函数卷积]\label{def:4-4-1-2}
若右端绝对收敛，定义
\[
 (f*g)(x)=\int_{\R^n}f(x-y)g(y)\dd y.
\]
若两个函数只给出 a.e. 等价类，卷积首先也是 a.e. 意义的对象；除非另有连续代表或额外正则性，
不把上式理解为每一点都收敛的公式。
\end{definition}

以后对 Lebesgue 可测函数使用 Tonelli 或 Fubini 时，先固定与其 a.e. 相等的 Borel 代表。
于是 $(x,y)\mapsto f(x-y)g(y)$ 是
$\Borel(\R^n)\otimes\Borel(\R^n)$-可测函数；所得卷积的 a.e. 等价类与代表选择无关。
这一步避免把两个完成化等价类的形式乘积直接当成已经给定的乘积可测函数。

\begin{definition}[测度卷积]\label{def:4-4-1-3}
设 $\mu,\nu$ 是 $\R^n$ 上有限 Borel 测度，$S(x,y)=x+y$。定义
\[
 \mu*\nu=S_*(\mu\otimes\nu).
\]
也就是说，对 Borel 集 $E$，
$ (\mu*\nu)(E)=\iint\1_E(x+y)\dd\mu(x)\dd\nu(y)$。
\end{definition}

\begin{example}[Dirac 测度与平移]\label{ex:4-4-1-2}
Dirac 测度把定义中的“平移”原样保留下来：
\[
 \delta_0*\mu=\mu,\qquad \delta_h*(f\,m)=(\tau_hf)\,m.
\]
第一式由 $0+x=x$ 得到；第二式对 Borel 集 $E$ 计算
$\int\1_E(h+y)f(y)\dd y=\int_Ef(x-h)\dd x$。所以 $\delta_0$ 是有限测度卷积的单位元，
却不是 $L^1$ 函数。因此 $L^1$ 函数的卷积代数没有单位元。
\end{example}

平移等距并不自动给出 $h\to0$ 时的连续性。若只在一个固定函数上观察，连续性成立；若同时在
$L^p$ 单位球上观察，它会失败。

\begin{theorem}[$L^p$ 平移连续]\label{theorem:4-4-1-1}
若 $1\le p<\infty$ 且 $f\in L^p(\R^n)$，则
\[
 \|\tau_hf-f\|_p\longrightarrow0\qquad(h\to0).
\]
\end{theorem}

\begin{proof}
先设 $f\in C_c(\R^n)$，并取紧集 $K$ 包含
$\supp f+B(0,1)$。当 $|h|<1$ 时，$\tau_hf-f$ 支撑于 $K$。函数 $f$ 在 $K$ 的一个紧邻域上
一致连续，故
\[
 \sup_x|f(x-h)-f(x)|\longrightarrow0.
\]
于是
\[
 \|\tau_hf-f\|_p^p
 \le m(K)\sup_x|f(x-h)-f(x)|^p\longrightarrow0.
\]

对一般 $f\in L^p$，给定 $\varepsilon>0$，由定理
\ref{theorem:3-3-3-1} 取 $g\in C_c(\R^n)$ 使 $\|f-g\|_p<\varepsilon$。平移等距性给出
\[
 \|\tau_hf-f\|_p
 \le \|\tau_h(f-g)\|_p+\|\tau_hg-g\|_p+\|g-f\|_p
 <2\varepsilon+\|\tau_hg-g\|_p.
\]
令 $h\to0$，再令 $\varepsilon\downarrow0$，即得结论。
\end{proof}

这个结论在 $p=\infty$ 处断裂。Heaviside 函数
$H=\1_{(0,\infty)}$ 满足 $\|\tau_hH-H\|_\infty=1$ 对每个 $h\ne0$ 成立。即使
$p<\infty$，定理说的也只是对每个固定 $f$ 的强收敛，而不是算子范数收敛。给定 $h\ne0$，取半径
小于 $|h|/3$ 的球 $E$，则 $E$ 与 $E+h$ 不交，从而
\[
 \|\tau_h\1_E-\1_E\|_p=2^{1/p}\|\1_E\|_p.
\]
因此 $\|\tau_h-I\|_{L^p\to L^p}\ge2^{1/p}$。逐向量连续和单位球上一致连续是两种不同要求。

卷积估计需要积分型 Minkowski 不等式；下面从标量积分与 $L^p$ 对偶性证明其所需形式。

\begin{lemma}[积分型 Minkowski 不等式]\label{lemma:4-4-1-minkowski-integral}
设 $1\le p\le\infty$，$F:\R^n\times\R^n\to\C$ 联合可测，并且
$y\mapsto\|F(\cdot,y)\|_p$ 可积。则 $\int F(\cdot,y)\dd y$ 对 a.e. $x$ 绝对存在，属于
$L^p$，并满足
\[
 \left\|\int F(\cdot,y)\dd y\right\|_p
 \le\int\|F(\cdot,y)\|_p\dd y.
\]
同一结论对非负 $F$ 的扩展实值积分成立。
\end{lemma}

\begin{proof}
对 $p=1$，Tonelli 定理给出
\[
 \int\left|\int F(x,y)\dd y\right|\dd x
 \le\iint|F(x,y)|\dd y\dd x
 =\int\|F(\cdot,y)\|_1\dd y.
\]
对 $p=\infty$，令 $a(y)=\|F(\cdot,y)\|_\infty$。可测集
$N=\{(x,y):|F(x,y)|>a(y)\}$ 的每个 $y$-截面都是零测集，故 Tonelli 定理给出
$(m\otimes m)(N)=0$。由 Fubini 定理，对 a.e. $x$ 有 $|F(x,y)|\le a(y)$ 对 a.e. $y$ 成立，
于是
\[
 \left|\int F(x,y)\dd y\right|\le\int a(y)\dd y.
\]
取本性上确界便得结论。

设 $1<p<\infty$，共轭指数为 $q$。先假设 $F$ 有界并支撑在一个有限测度矩形中。由
定理~\ref{theorem:3-3-3-3} 的 $L^p$ 对偶表示和 Fubini 定理，
\begin{align*}
 \left\|\int F(\cdot,y)\dd y\right\|_p
 &=\sup_{\|\varphi\|_q\le1}
 \left|\int\!\int F(x,y)\overline{\varphi(x)}\dd y\dd x\right|\\
 &\le\sup_{\|\varphi\|_q\le1}
 \int\|F(\cdot,y)\|_p\|\varphi\|_q\dd y
 \le\int\|F(\cdot,y)\|_p\dd y.
\end{align*}
一般情形令
$F_N=F\1_{B(0,N)\times B(0,N)}\1_{\{|F|\le N\}}$。把上述估计用于
$|F_N|$，并以 Fatou 引理和单调收敛定理令 $N\to\infty$，便得到
$x\mapsto\int|F(x,y)|\dd y$ 属于 $L^p$ 及所需估计。复值积分随后由绝对存在性定义。
\end{proof}

\begin{theorem}[Young 不等式：基础型]\label{theorem:4-4-1-2}
若 $f\in L^1(\R^n)$、$g\in L^p(\R^n)$，其中 $1\le p\le\infty$，则
$f*g$ 对 a.e. $x$ 绝对存在，属于 $L^p$，并且
\[
 \|f*g\|_p\le\|f\|_1\|g\|_p.
\]
\end{theorem}

\begin{proof}
换元后可写
\[
 |f*g(x)|\le\int|f(y)|\,|g(x-y)|\dd y.
\]
对 $F(x,y)=|f(y)|\,|g(x-y)|$ 应用引理
\ref{lemma:4-4-1-minkowski-integral}，再用平移等距性，得到
\[
 \|f*g\|_p
 \le\int|f(y)|\,\|\tau_yg\|_p\dd y
 =\|f\|_1\|g\|_p.
\]
同一引理同时保证右侧定义卷积的绝对积分在 a.e. $x$ 有限。
\end{proof}

\begin{proposition}[卷积代数]\label{proposition:4-4-1-3}
$L^1(\R^n)$ 在卷积下是交换 Banach 代数：
\[
 f*g=g*f,\qquad (f*g)*h=f*(g*h),\qquad
 \|f*g\|_1\le\|f\|_1\|g\|_1.
\]
此外，$\tau_h(f*g)=(\tau_hf)*g=f*(\tau_hg)$，只要相应卷积有意义。
\end{proposition}

\begin{proof}
范数估计是定理~\ref{theorem:4-4-1-2} 的 $p=1$ 情形。交换律由 $z=x-y$ 换元：
\[
 \int f(x-y)g(y)\dd y=\int g(x-z)f(z)\dd z.
\]
对结合律，Tonelli 定理给出
\[
 \iiint |f(x-y-z)g(y)h(z)|\dd y\dd z\dd x
 =\|f\|_1\|g\|_1\|h\|_1<\infty.
\]
因而可用 Fubini 交换 $y,z$ 的积分并作线性换元，得到
$((f*g)*h)(x)=(f*(g*h))(x)$ 对 a.e. $x$ 成立。平移协变性直接由定义中的变量替换得到。
$L^1$ 已由定理~\ref{theorem:3-3-2-3} 证明完备，故这是 Banach 代数。
\end{proof}

有限测度的卷积也交换并结合。这里不需要写三重积分的形式符号：对有界非负 Borel 函数 $\varphi$，
推前积分公式给出
\[
 \int\varphi\dd((\mu*\nu)*\sigma)
 =\iiint\varphi(x+y+z)\dd\mu(x)\dd\nu(y)\dd\sigma(z),
\]
Tonelli 定理说明另一个括号给出同一数值。若 $\mu=fm$、$\nu=gm$，则对非负 Borel
$\varphi$ 再用 Tonelli 和换元可得
\[
 \int\varphi\dd(\mu*\nu)
 =\int\varphi(z)(f*g)(z)\dd z,
\]
所以测度卷积的密度正是函数卷积。

\begin{proposition}[Young 不等式：一般指数]\label{proposition:4-4-1-4}
设 $1\le p,q,r\le\infty$，且
\[
 \frac1p+\frac1q\ge1,\qquad
 1+\frac1r=\frac1p+\frac1q.
\]
若 $f\in L^p$、$g\in L^q$，则 $f*g\in L^r$ 且
$\|f*g\|_r\le\|f\|_p\|g\|_q$。
\end{proposition}

\begin{proof}
若 $r=\infty$，则 $q=p'$，逐点 \Holder{} 不等式给出
$|f*g(x)|\le\|f\|_p\|g\|_q$。若 $p=1$ 或 $q=1$，结论由基础型 Young 不等式及交换律得到。

以下设 $p,q,r<\infty$；由指数关系 $r\ge p,q$。先对非负、有界且紧支撑的 $f,g$ 计算。对固定
$x$，把被积函数分解为
\[
 |f(y)g(x-y)|
 =\bigl(|f(y)|^p|g(x-y)|^q\bigr)^{1/r}
 |f(y)|^{1-p/r}|g(x-y)|^{1-q/r}.
\]
三个 \Holder{} 指数的倒数分别为
\[
 \frac1r,\qquad \frac{1-p/r}{p},\qquad
 \frac{1-q/r}{q};
\]
它们之和等于一；若某个倒数为零，就把相应的恒等于一的因子删去。因此广义 \Holder{} 不等式给出
\[
 |f*g(x)|^r
 \le \|f\|_p^{r-p}\|g\|_q^{r-q}
 \int |f(y)|^p|g(x-y)|^q\dd y.
\]
对 $x$ 积分并用 Tonelli 定理及平移不变性，得到
\[
 \|f*g\|_r^r
 \le \|f\|_p^{r-p}\|g\|_q^{r-q}
 \|f\|_p^p\|g\|_q^q
 =\|f\|_p^r\|g\|_q^r.
\]
取 $r$ 次方根得到所需估计。

对一般 $f\in L^p$、$g\in L^q$，令
\[
 f_N=|f|\1_{B(0,N)}\1_{\{|f|\leq N\}},\qquad
 g_N=|g|\1_{B(0,N)}\1_{\{|g|\leq N\}}.
\]
则 $f_N\uparrow|f|$、$g_N\uparrow|g|$ a.e.，而已证估计给出
\[
 \|f_N*g_N\|_r^r
 \leq\|f_N\|_p^r\|g_N\|_q^r
 \leq\|f\|_p^r\|g\|_q^r.
\]
对每个 $x$，被积函数 $f_N(y)g_N(x-y)$ 随 $N$ 单调增加；单调收敛定理因而给出
\[
 (f_N*g_N)(x)\uparrow
 h(x):=\int_{\R^n}|f(y)|\,|g(x-y)|\dd y\in[0,\infty].
\]
再对 $x$ 使用单调收敛定理，得到
\[
 \|h\|_r^r=\lim_{N\to\infty}\|f_N*g_N\|_r^r
 \leq\|f\|_p^r\|g\|_q^r<\infty.
\]
所以 $h(x)<\infty$ 对 a.e. $x$ 成立；这些点上原卷积绝对存在且
$|f*g(x)|\leq h(x)$。于是
$\|f*g\|_r\leq\|h\|_r\leq\|f\|_p\|g\|_q$，完成一般情形。
\end{proof}

\begin{example}[Gaussian 测度在卷积下封闭]\label{ex:4-4-1-3}
一维中令
\[
 \gamma_{m,\sigma^2}(x)=\frac1{\sqrt{2\pi\sigma^2}}
 \exp\!\left[-\frac{(x-m)^2}{2\sigma^2}\right],\qquad\sigma>0.
\]
在卷积积分中完成平方：若 $v=\sigma_1^2+\sigma_2^2$，则
\begin{align*}
 &\frac{(x-y-m_1)^2}{2\sigma_1^2}
 +\frac{(y-m_2)^2}{2\sigma_2^2}\\
 &\quad=
 \frac{(x-m_1-m_2)^2}{2v}
 +\frac{v}{2\sigma_1^2\sigma_2^2}
 \left(y-m_2-\frac{\sigma_2^2}{v}(x-m_1-m_2)\right)^2.
\end{align*}
余下的 Gaussian 积分等于
$\sqrt{2\pi\sigma_1^2\sigma_2^2/v}$，故
\[
 \gamma_{m_1,\sigma_1^2}*\gamma_{m_2,\sigma_2^2}
 =\gamma_{m_1+m_2,\,\sigma_1^2+\sigma_2^2}.
\]
多维正定协方差的结论也可逐项核对如下。令
\[
 \gamma_{m,\Sigma}(x)
 =(2\pi)^{-n/2}(\det\Sigma)^{-1/2}
 \exp\!\left[-\frac12(x-m)^T\Sigma^{-1}(x-m)\right],
\]
并置 $z=x-m_1-m_2$、$w=y-m_2$ 及
$K=\Sigma_1^{-1}+\Sigma_2^{-1}$。则
\[
 \begin{aligned}
 &(z-w)^T\Sigma_1^{-1}(z-w)+w^T\Sigma_2^{-1}w\\
 &\quad=z^T(\Sigma_1+\Sigma_2)^{-1}z
 +(w-K^{-1}\Sigma_1^{-1}z)^TK
   (w-K^{-1}\Sigma_1^{-1}z).
 \end{aligned}
\]
因此
\[
 \begin{aligned}
 (\gamma_{m_1,\Sigma_1}*\gamma_{m_2,\Sigma_2})(x)
 &=\frac{e^{-z^T(\Sigma_1+\Sigma_2)^{-1}z/2}}
 {(2\pi)^n(\det\Sigma_1\det\Sigma_2)^{1/2}}
 \int_{\R^n}e^{-(w-K^{-1}\Sigma_1^{-1}z)^TK
 (w-K^{-1}\Sigma_1^{-1}z)/2}\,\dd w\\
 &=\frac{e^{-z^T(\Sigma_1+\Sigma_2)^{-1}z/2}}
 {(2\pi)^{n/2}\det(\Sigma_1+\Sigma_2)^{1/2}}
 =\gamma_{m_1+m_2,\,\Sigma_1+\Sigma_2}(x),
 \end{aligned}
\]
其中用了
\[
 \int_{\R^n}e^{-v^TKv/2}\,\dd v
 =(2\pi)^{n/2}(\det K)^{-1/2},
 \qquad
 \det K\det\Sigma_1\det\Sigma_2=\det(\Sigma_1+\Sigma_2).
\]
第一个积分恒等式由正交对角化
$K=U^T\operatorname{diag}(\lambda_1,\ldots,\lambda_n)U$、正交换元以及 $n$ 个一维
Gaussian 积分相乘得到；第二个恒等式来自
$K=\Sigma_1^{-1}(\Sigma_1+\Sigma_2)\Sigma_2^{-1}$ 后对行列式取积。
若协方差退化，密度公式已经不存在。
这时在 $\R^n$ 上取标准 Gaussian 测度 $\gamma_n$，并选择线性映射
$A_i:\R^n\to\R^n$ 使 $A_iA_i^T=\Sigma_i$，把相应测度写成
$(m_i+A_i\,\cdot)_*\gamma_n$。两份测度的卷积是积测度
$\gamma_n\otimes\gamma_n$ 经
$(u,v)\mapsto m_1+m_2+A_1u+A_2v$ 推前，仍给出协方差
$A_1A_1^T+A_2A_2^T$ 的 Gaussian 测度。推前表述同时涵盖退化情形，无需假设相应测度具有
Lebesgue 密度。
\end{example}

\begin{figure}[htbp]
\centering
\begin{tikzpicture}[x=1cm,y=1cm]
  \draw[->,BookInk!70] (-0.2,0)--(7.4,0) node[right] {$y$};
  \draw[BookAccent,thick] plot[smooth] coordinates
    {(0.2,0.05) (0.8,0.25) (1.4,1.0) (2.0,1.55) (2.6,0.95) (3.2,0.2) (3.7,0.04)};
  \draw[BookWarm,thick,dashed] plot[smooth] coordinates
    {(2.8,0.05) (3.4,0.32) (4.0,1.15) (4.6,1.45) (5.2,0.8) (5.8,0.18) (6.4,0.04)};
  \fill[BookAccent!12] (3.0,0)--plot[smooth] coordinates
    {(3.0,0.12) (3.3,0.24) (3.6,0.55) (3.9,0.98) (4.2,1.30) (4.5,1.42)}--(4.5,0)--cycle;
  \node[BookAccent] at (1.6,1.9) {$f(y)$};
  \node[BookWarm] at (5.4,1.75) {$g(x-y)$};
  \node[book strong node,text width=3.3cm] at (3.6,-1.0)
    {平移量 $x$ 改变重叠区域；\\重叠的加权面积形成 $(f*g)(x)$};
  \begin{scope}[xshift=8.3cm]
    \node[book node] (prod) at (0,1.25) {$\mu\otimes\nu$};
    \node[book strong node] (sum) at (3.15,1.25) {$\mu*\nu$};
    \draw[book accent arrow] (prod)--node[above,book label]
      {$S(x,y)=x+y$}(sum);
    \draw[BookInk!45,thin] (-.9,-.25) grid[xstep=.38,ystep=.28] (.9,.65);
    \foreach \x/\y in {-.65/.05,-.25/.4,.15/.1,.55/.48,.7/.18,-.45/.55}
      \fill[BookWarm] (\x,\y) circle (1.5pt);
    \draw[BookAccent,thick] (2.35,-.1).. controls (2.75,.45) and
      (3.45,.62)..(3.9,.12);
    \node[book label,align=center] at (1.55,-.85)
      {积空间中的质量沿\\等和方向汇入一维};
  \end{scope}
\end{tikzpicture}
\caption{左：卷积把相对位移下的全部重叠积分为一个数值。右：测度卷积是
$\mu\otimes\nu$ 经加法映射的推前。图形说明两种来源，但不能替代 Fubini、Tonelli 与 Young
不等式对积分存在性和换序的证明。}
\label{fig:4-4-1-convolution-overlap}
\end{figure}

\begin{remark}[与概率卷积的关系]
总质量一的有限测度在卷积下仍有总质量一。第 5 章定义随机向量和独立性后，$\mu*\nu$ 将解释为
两个独立随机向量之和的分布。这里的测度定义已经给出了这一解释所需的乘积测度与加法推前结构。
\end{remark}

\begin{exercise}
\begin{enumerate}
  \item 直接从区间交集重新推导式~\eqref{eq:indicator-convolution}，并验证其积分等于一。
  \item 证明若有限测度 $\mu,\nu$ 分别有密度 $f,g$，则 $\mu*\nu$ 有密度 $f*g$。
  \item 证明 $L^1$ 卷积的结合律，并指出绝对可积性在哪一步使用。
  \item 构造 $h_k\to0$ 与 $f_k\in L^p$、$\|f_k\|_p=1$，使
        $\|\tau_{h_k}f_k-f_k\|_p$ 不趋于零。
\end{enumerate}
\end{exercise}

卷积核若总质量为一并逐渐集中到原点，移动平均应当回到原函数。下一小节把这一现象从标准
mollifier 中抽离出来，并比较热核与 Poisson 核的不同时间尺度。


\subsection{近似恒等、Poisson 核与热核}\label{subsec:4-4-2}
\index{近似恒等}\index{热核}\index{Poisson 核}

“核越来越尖”并不是数学条件：尖峰可能具有错误的总质量，可能在远处留下不消失的尾部，也可能因
正负振荡而放大误差。近似恒等的定义相应要求保存常数、统一控制总变差，并使总变差集中到原点。

Poisson 核最初用于调和边值问题，热核则描述扩散方程；二者也是 Fourier 乘子与 Markov 演化的
基本模型。本小节用实分析证明质量集中、边界恢复与热核半群；Poisson 半群的 Fourier 证明将在
第~7 章给出。

\begin{definition}[近似恒等]\label{def:4-4-2-1}
一族 $\phi_\varepsilon\in L^1(\R^n)$（$\varepsilon>0$）称为近似恒等，如果
\[
 \int_{\R^n}\phi_\varepsilon(x)\dd x=1,
 \qquad \sup_{\varepsilon>0}\|\phi_\varepsilon\|_1<\infty,
\]
并且对每个 $\delta>0$，
\[
 \int_{|x|>\delta}|\phi_\varepsilon(x)|\dd x\longrightarrow0
 \qquad(\varepsilon\downarrow0).
\]
这里不要求 $\phi_\varepsilon$ 非负。
\end{definition}

最常见的构造是从一个 $\phi\in L^1$、$\int\phi=1$ 出发，令
$\phi_\varepsilon(x)=\varepsilon^{-n}\phi(x/\varepsilon)$。变量替换立即验证前两项，而
\[
 \int_{|x|>\delta}|\phi_\varepsilon(x)|\dd x
 =\int_{|z|>\delta/\varepsilon}|\phi(z)|\dd z\longrightarrow0
\]
验证了质量集中。

\begin{example}[正性不是定义的一部分]\label{ex:4-4-2-1}
取 $\eta\in C_c^\infty(\R^n)$，$\eta\ge0$ 且 $\int\eta=1$，并令
$\phi=\eta+a\partial_1\eta$，其中 $a\ne0$。由于紧支撑函数导数的积分为零，
$\int\phi=1$；相应伸缩族 $\phi_\varepsilon$ 是近似恒等，尽管适当选取 $a$ 后 $\phi$ 会变号。

统一 $L^1$ 控制却不能删掉。为看清这一点，可在圆周 $\mathbb T=\R/(2\pi\Z)$ 上比较
\[
 D_N(\theta)=\sum_{k=-N}^Ne^{ik\theta}
 =\frac{\sin((N+\tfrac12)\theta)}{\sin(\theta/2)}
\]
与 \Fejer{} 核
\[
 F_N(\theta)=\frac1{N+1}\sum_{j=0}^ND_j(\theta)
 =\frac1{N+1}
 \left(\frac{\sin((N+1)\theta/2)}{\sin(\theta/2)}\right)^2.
\]
二者相对于归一化测度 $(2\pi)^{-1}\dd\theta$ 都有积分一，但
$\|D_N\|_{L^1(\mathbb T)}$ 随 $N$ 呈对数增长。这里把这个常被一句带过的估计写出。
在 $|\theta|\le\pi$ 上，
\[
 |D_N(\theta)|\le
 \min\{2N+1,\pi/|\theta|\},
\]
所以分开 $|\theta|\le(N+1)^{-1}$ 与其余部分积分，得到
$\|D_N\|_1\le C\log(N+2)$。反向令 $M=N+\tfrac12$，并在
\[
 I_j=\left[\frac{(j+1/6)\pi}{M},
             \frac{(j+5/6)\pi}{M}\right],
 \qquad 1\le j\le N-1,
\]
上积分。此处 $|\sin(M\theta)|\ge1/2$，而
$\sin(\theta/2)\le\theta/2$，故
\[
 \int_{I_j}|D_N(\theta)|\dd\theta
 \ge\int_{I_j}\frac{\dd\theta}{\theta}
 =\log\frac{j+5/6}{j+1/6}\ge\frac{c}{j+1}.
\]
求和给 $\|D_N\|_1\ge c\log(N+1)$。另一方面，$F_N\ge0$，而逐项积分
$\int_{\mathbb T}D_j\,\dd\theta/(2\pi)=1$，所以 $\|F_N\|_1=1$。
这只是圆周上的旁证，不把周期核冒充为定义
\ref{def:4-4-2-1} 中的 $\R^n$ 核；它准确指出，缺少统一总变差控制时，质量一并不足以抑制振荡。
\end{example}

\begin{definition}[热核]\label{def:4-4-2-2}
对 $t>0$，定义
\[
 p_t(x)=(4\pi t)^{-n/2}\exp\!\left(-\frac{|x|^2}{4t}\right),
 \qquad H_tf=p_t*f.
\]
\end{definition}

\begin{definition}[Poisson 核]\label{def:4-4-2-3}
令
\[
c_n=\left(\int_{\R^n}(1+|z|^2)^{-(n+1)/2}\dd z\right)^{-1},
\]
并定义上半空间的 Poisson 核与 Poisson 算子为
\[
 q_t(x)=c_n\frac{t}{(t^2+|x|^2)^{(n+1)/2}},
 \qquad Q_tf=q_t*f,\qquad t>0.
\]
\end{definition}

\begin{proposition}[Poisson 核的归一化]
若 Euler Gamma 函数定义为
\[
 \Gamma(a)=\int_0^\infty s^{a-1}e^{-s}\,\dd s,
 \qquad a>0,
\]
则
\[
 c_n=\frac{\Gamma((n+1)/2)}{\pi^{(n+1)/2}},
 \qquad \int_{\R^n}q_t(x)\,\dd x=1.
\]
\end{proposition}

\begin{proof}
先把所用常数公式写明。对 $a,b>0$ 定义
\[
 B(a,b):=\int_0^\infty
   \frac{v^{a-1}}{(1+v)^{a+b}}\dd v.
\]
Tonelli 定理以及换元
$s=\rho\theta$、$t=\rho(1-\theta)$（Jacobian 为 $\rho$）给出
\begin{align*}
 \Gamma(a)\Gamma(b)
 &=\int_0^\infty\!\int_0^\infty
   s^{a-1}t^{b-1}e^{-(s+t)}\dd s\dd t\\
 &=\left(\int_0^\infty\rho^{a+b-1}e^{-\rho}\dd\rho\right)
   \left(\int_0^1\theta^{a-1}(1-\theta)^{b-1}\dd\theta\right)\\
 &=\Gamma(a+b)B(a,b),
\end{align*}
其中最后一步在第二个积分中用了 $v=\theta/(1-\theta)$。此外，一维 Gaussian 积分和
$s=r^2$ 给出
\[
 \Gamma(1/2)=2\int_0^\infty e^{-r^2}\dd r=\sqrt\pi.
\]
再将 $\int_{\R^n}e^{-|x|^2}\dd x=\pi^{n/2}$ 写成极坐标积分，得到
\[
 \pi^{n/2}=|\mathbb S^{n-1}|\int_0^\infty e^{-r^2}r^{n-1}\dd r
 =\frac{|\mathbb S^{n-1}|}{2}\Gamma(n/2),
\]
即 $|\mathbb S^{n-1}|=2\pi^{n/2}/\Gamma(n/2)$。现在作代换 $u=r^2$，有
\begin{align*}
 \int_{\R^n}(1+|z|^2)^{-(n+1)/2}\dd z
 &=|\mathbb S^{n-1}|\int_0^\infty
   \frac{r^{n-1}}{(1+r^2)^{(n+1)/2}}\,\dd r\\
 &=\frac{|\mathbb S^{n-1}|}{2}B\!\left(\frac n2,\frac12\right)
 =\frac{\pi^{(n+1)/2}}{\Gamma((n+1)/2)}.
\end{align*}
取倒数得到 $c_n$ 的公式；再令 $x=tz$，便有
\[
 \int_{\R^n}q_t(x)\,\dd x
 =c_n\int_{\R^n}(1+|z|^2)^{-(n+1)/2}\dd z=1.
\]
\end{proof}

符号 $H_t$ 与 $Q_t$ 分别表示热算子和 Poisson 算子；两族核的尺度与方程不同。

\begin{example}[热核的质量与扩散尺度]\label{ex:4-4-2-2}
一维 Gaussian 积分
\[
 \int_\R e^{-x^2/(4t)}\dd x=\sqrt{4\pi t}
\]
及乘积结构给出 $\int_{\R^n}p_t=1$。奇偶性给出
$\int x_jp_t(x)\dd x=0$。再对
$\int_\R e^{-a x^2}\dd x=\sqrt{\pi/a}$ 关于 $a$ 求导，并取 $a=(4t)^{-1}$，得到
\[
 \int_\R x^2(4\pi t)^{-1/2}e^{-x^2/(4t)}\dd x=2t.
\]
因此
\[
 \int_{\R^n}x_j^2p_t(x)\dd x=2t,
 \qquad \int_{\R^n}|x|^2p_t(x)\dd x=2nt.
\]
由点态不等式
$\1_{\{|x|>\delta\}}\le |x|^2/\delta^2$ 直接积分，得到
$\int_{|x|>\delta}p_t(x)\dd x\le2nt/\delta^2$；这表明 $t\downarrow0$ 时质量集中于原点，
二阶矩则表明典型扩散距离具有 $\sqrt t$ 的尺度。
\end{example}

\begin{theorem}[近似恒等的范数收敛]\label{theorem:4-4-2-1}
设 $(\phi_\varepsilon)$ 是近似恒等。
若 $1\le p<\infty$ 且 $f\in L^p(\R^n)$，则
\[
 \|\phi_\varepsilon*f-f\|_p\longrightarrow0.
\]
若 $f\in C_0(\R^n)$，则上述收敛在一致范数中成立。
\end{theorem}

\begin{proof}
记 $C=\sup_\varepsilon\|\phi_\varepsilon\|_1$。由积分型 Minkowski 不等式，
\begin{equation}\label{eq:approx-id-translation}
 \|\phi_\varepsilon*f-f\|_p
 \le\int_{\R^n}|\phi_\varepsilon(y)|
       \|\tau_yf-f\|_p\dd y.
\end{equation}
给定 $\eta>0$，定理~\ref{theorem:4-4-1-1} 给出 $\delta>0$，使
$|y|<\delta$ 时 $\|\tau_yf-f\|_p<\eta/(2C)$。又恒有
$\|\tau_yf-f\|_p\le2\|f\|_p$，故式~\eqref{eq:approx-id-translation} 的右端至多为
\[
 \frac\eta{2C}\int_{|y|<\delta}|\phi_\varepsilon(y)|\dd y
 +2\|f\|_p\int_{|y|\ge\delta}|\phi_\varepsilon(y)|\dd y.
\]
第一项不超过 $\eta/2$，第二项在 $\varepsilon\downarrow0$ 时趋于零。

若 $f\in C_0$，则 $f$ 在 $\R^n$ 上一致连续。确实，先在一个足够大的紧球内使用一致连续性，
球外则用 $f(x)\to0$。把上面的 $L^p$ 平移模量换成
$\|\tau_yf-f\|_\infty$，同一分割证明一致收敛。
\end{proof}

$L^p$ 收敛只使用总变差集中，并不保证每个指定点都收敛。点态恢复还要求核不把质量集中在越来越
细的偏心环层中；径向递减上界排除了这种偏心集中。

\begin{theorem}[Lebesgue 点处的恢复]\label{theorem:4-4-2-2}
设 $1\le p\le\infty$、$f\in L^p(\R^n)$，并令
$\phi_\varepsilon(x)=\varepsilon^{-n}\phi(x/\varepsilon)$，其中 $\int\phi=1$。
假设存在径向递减函数 $\Psi\in L^1(\R^n)$ 使 $|\phi|\le\Psi$。
则在 $f$ 的每个 Lebesgue 点 $x$，
\[
 (\phi_\varepsilon*f)(x)\longrightarrow f(x).
\]
此外，在右侧有意义的每一点，
\[
 \sup_{\varepsilon>0}|\phi_\varepsilon*f(x)|
 \le C_n\|\Psi\|_1Mf(x),
\]
其中 $M$ 是定义~\ref{def:4-1-1-1} 的中心极大函数，$C_n$ 只依赖维数。
\end{theorem}

\begin{proof}
把 $\Psi(z)$ 的径向递减代表写成 $\psi(|z|)$，并令
$A_k=\{z:2^{k-1}<|z|\le2^k\}$，$k\in\Z$。在 $A_k$ 上
$\Psi(z)\le\psi(2^{k-1})$，从而对任意局部可积 $g\ge0$，
\begin{align}
 \int\varepsilon^{-n}\Psi(y/\varepsilon)g(x-y)\dd y
 &\le \sum_{k\in\Z}\varepsilon^{-n}\psi(2^{k-1})
       \int_{|y|\le2^k\varepsilon}g(x-y)\dd y \notag\\
 &\le \sum_{k\in\Z}a_k
   \frac1{|B(x,2^k\varepsilon)|}
   \int_{B(x,2^k\varepsilon)}g(z)\dd z,                 \label{eq:radial-majorant}
\end{align}
其中 $a_k=|B(0,1)|2^{kn}\psi(2^{k-1})$。环层
$\{2^{k-2}<|z|\le2^{k-1}\}$ 两两不交，而且在该环层上
$\Psi(z)\ge\psi(2^{k-1})$。记 $\omega_n=|B(0,1)|$，则逐项有
\begin{align*}
 \int_{2^{k-2}<|z|\leq2^{k-1}}\Psi(z)\dd z
 &\geq \omega_n\bigl(2^{n(k-1)}-2^{n(k-2)}\bigr)
          \psi(2^{k-1})\\
 &=\frac{2^n-1}{2^{2n}}a_k.
\end{align*}
对 $k\in\Z$ 求和并使用这些环层的不交性，故
\begin{equation}\label{eq:radial-coefficients}
 \sum_{k\in\Z}a_k
 \leq\frac{2^{2n}}{2^n-1}\|\Psi\|_1.
\end{equation}
将 $g=|f|$ 代入式~\eqref{eq:radial-majorant}，便得到所述极大函数控制。

现在固定一个 Lebesgue 点 $x$，并取
$g_x(z)=|f(z)-f(x)|$。Lebesgue 点条件恰好说
\[
 \frac1{|B(x,r)|}\int_{B(x,r)}g_x(z)\dd z\longrightarrow0
 \qquad(r\downarrow0).
\]
取 $r_0>0$，使 $0<r<r_0$ 时这些平均不超过一。若 $1\leq p<\infty$ 且 $r\geq r_0$，
\Holder{} 不等式给出
\[
 \frac1{|B(x,r)|}\int_{B(x,r)}g_x
 \leq |B(x,r)|^{-1/p}\|f\|_p+|f(x)|
 \leq |B(0,r_0)|^{-1/p}\|f\|_p+|f(x)|.
\]
若 $p=\infty$，相应平均至多为 $\|f\|_\infty+|f(x)|$。所以这些球平均对全部 $r>0$ 确有
有限共同上界。
式~\eqref{eq:radial-majorant} 与式~\eqref{eq:radial-coefficients} 给出
\[
 |\phi_\varepsilon*f(x)-f(x)|
 \le\sum_{k\in\Z}a_k
 \frac1{|B(x,2^k\varepsilon)|}
 \int_{B(x,2^k\varepsilon)}g_x(z)\dd z.
\]
对每个固定 $k$，右侧相应球平均随 $\varepsilon\downarrow0$ 趋于零；又由
$\sum_ka_k<\infty$ 和刚得到的共同上界可在计数测度上使用支配收敛定理。因此右侧趋于零。
\end{proof}

热核与 Poisson 核都是非负、径向递减、质量一的伸缩族：
$p_t(x)=t^{-n/2}p_1(x/\sqrt t)$，而 $q_t(x)=t^{-n}q_1(x/t)$。
所以前两定理同时给出它们的 $L^p$、$C_0$ 与 Lebesgue 点边界恢复。两个尺度
$\sqrt t$ 与 $t$ 的差别将在它们满足的方程中再次出现。

\begin{proposition}[热核半群]\label{proposition:4-4-2-3}
对 $s,t>0$，
\[
 p_s*p_t=p_{s+t},\qquad H_sH_t=H_{s+t}.
\]
因而 $H_t$ 在每个 $L^p(\R^n)$（$1\le p\le\infty$）上为正收缩；当
$1\le p<\infty$ 时，$H_tf\to f$ 于 $L^p$。
\end{proposition}

\begin{proof}
完成平方得到
\[
 \frac{|x-y|^2}{4s}+\frac{|y|^2}{4t}
 =\frac{|x|^2}{4(s+t)}
 +\frac{s+t}{4st}\left|y-\frac{t}{s+t}x\right|^2.
\]
因此
\begin{align*}
 (p_s*p_t)(x)
 &=(4\pi s)^{-n/2}(4\pi t)^{-n/2}
 e^{-|x|^2/(4(s+t))}
 \int_{\R^n}e^{-(s+t)|y-tx/(s+t)|^2/(4st)}\dd y\\
 &=(4\pi(s+t))^{-n/2}e^{-|x|^2/(4(s+t))}=p_{s+t}(x).
\end{align*}
结合律给出算子半群恒等。正性来自 $p_t\ge0$，而基础型 Young 不等式与
$\|p_t\|_1=1$ 给出收缩性。最后，$(p_t)_{t>0}$ 在 $t\downarrow0$ 时是近似恒等，故
定理~\ref{theorem:4-4-2-1} 给出强连续性。
\end{proof}

\begin{remark}[Poisson 半群的 Fourier 表示]
恒等式 $q_s*q_t=q_{s+t}$ 成立，并推出 $Q_sQ_t=Q_{s+t}$。按照第~7 章采用的 Fourier 约定
$e^{-2\pi i x\cdot\xi}$，可由
$\widehat q_t(\xi)=e^{-2\pi t|\xi|}$ 立即验证。该 Fourier 公式将在第~7 章证明。

同一约定下，第~7 章还将核对
\[
 \widehat p_t(\xi)=e^{-4\pi^2t|\xi|^2},
 \qquad
 \widehat q_t(\xi)=e^{-2\pi t|\xi|}.
\]
若把 Fourier 核改成 $e^{-ix\cdot\xi}$，相应乘子则写成
$e^{-t|\xi|^2}$ 与 $e^{-t|\xi|}$。热核半群已经由上面的完成平方计算直接得到。
\end{remark}

\begin{proposition}[热方程与 Poisson 方程]\label{proposition:4-4-2-4}
若 $f\in C_c^\infty(\R^n)$，则
\[
 u(x,t)=H_tf(x),\qquad v(x,t)=Q_tf(x)
\]
分别满足
\[
 \partial_tu=\Delta_xu,\qquad
 (\Delta_x+\partial_t^2)v=0\qquad(t>0).
\]
对任意 $f\in L^p$、$1\le p\le\infty$，这些等式在每个固定的正时间带
$t\in[a,b]\subset(0,\infty)$ 上仍以经典空间导数和 $L^p$ 值时间导数成立。
在 $t=0$ 的初值或边界值恢复则分别采用定理~\ref{theorem:4-4-2-1} 与
定理~\ref{theorem:4-4-2-2}，不把正时间微分与边界极限混作同一步。
\end{proposition}

\begin{proof}
直接求导给出
\[
 \partial_tp_t(x)=
 \left(-\frac n{2t}+\frac{|x|^2}{4t^2}\right)p_t(x)
 =\Delta_xp_t(x).
\]
对 Poisson 核令 $R=t^2+|x|^2$、$a=(n+1)/2$。略去共同常数 $c_n$，逐项计算为
\begin{align*}
 \Delta_x(tR^{-a})
 &=-2antR^{-a-1}+4a(a+1)t|x|^2R^{-a-2},\\
 \partial_t^2(tR^{-a})
 &=-6atR^{-a-1}+4a(a+1)t^3R^{-a-2}.
\end{align*}
两式相加并利用 $|x|^2+t^2=R$，得到
\begin{align*}
 &(\Delta_x+\partial_t^2)(tR^{-a})\\
 &\quad=\bigl[-2a(n+3)+4a(a+1)\bigr]tR^{-a-1}=0,
\end{align*}
其中最后一步使用 $2(a+1)=n+3$。

对 $f\in C_c^\infty$，上述核及其在正时间带上的任意固定阶导数都有可积上界，故支配收敛定理
允许在卷积积分号下微分，并把核等式传给 $u,v$。

再考虑粗糙的 $f$。记 $k_t$ 为 $p_t$ 或 $q_t$。在任意紧正时间带
$[a,b]\subset(0,\infty)$ 上先证明所需的核估计。对热核，逐次求导给出
\[
 \partial_t^j\partial_x^\alpha p_t(x)
 =t^{-n/2-j-|\alpha|/2}
   P_{j,\alpha}\!\left(\frac{x}{\sqrt t}\right)
   \exp\!\left(-\frac{|x|^2}{4t}\right).
\]
其中 $P_{j,\alpha}$ 是一个多项式。对 Poisson 核，由
$q_t(x)=t^{-n}q_1(x/t)$ 归纳得到
\[
 \partial_t^j\partial_x^\alpha q_t(x)
 =t^{-n-j-|\alpha|}Q_{j,\alpha}(x/t),
 \qquad
 |Q_{j,\alpha}(z)|\le
 C_{j,\alpha}(1+|z|)^{-n-1-|\alpha|}.
\]
这里后一估计从
$q_1(z)=c_n(1+|z|^2)^{-(n+1)/2}$ 开始归纳：空间求导使右端的衰减幂至少增加一，
而一次时间求导在缩放变量中只把剖面 $Q$ 换成
$-(n+j+|\alpha|)Q-z\cdot\nabla Q$，不减弱已有的衰减幂。

第一个缩放公式经变量代换 $x=\sqrt t\,z$，第二个缩放公式经变量代换 $x=tz$，表明
方程中出现的每个 $\partial_t^j\partial_x^\alpha k_t$ 都属于 $L^1\cap C_0$，并且
其 $L^1$ 与一致范数在 $t\in[a,b]$ 上一致有界。它们还在这两个范数中关于 $t$ 连续。
事实上，若 $t_\nu\to t\in[a,b]$，则相应导数逐点收敛；热核族有形如
$C(1+|x|)^N e^{-|x|^2/(4b)}$ 的共同可积上界，Poisson 核族由
$C(1+|x|/b)^{-n-1-|\alpha|}$ 控制，故 $L^1$ 连续性来自支配收敛定理。
在任意固定紧集上的收敛由关于 $(t,x)$ 的联合连续性而一致；上述两个上界又使紧集外的
一致范数关于 $t\in[a,b]$ 一致趋于零，因而得到 $C_0$ 范数连续性。特别地，这些导数在
$L^1$ 与 $C_0$ 中都具有空间平移连续性。

先固定 $t$。对空间一阶差商，一维微积分基本定理给出
\[
 \frac{k_t(\,\cdot+he_j)-k_t}{h}-\partial_jk_t
 =\frac1h\int_0^h
   \bigl[\partial_jk_t(\,\cdot+se_j)-\partial_jk_t\bigr]\,\dd s.
\]
$L^1$ 平移连续性使右端的 $L^1$ 范数趋于零；把 $k_t$ 依次换成它的空间导数，便得到二阶空间
差商的同一结论。时间方向同样由
\[
 \frac{k_{t+h}-k_t}{h}-\partial_tk_t
 =\frac1h\int_0^h
 \bigl(\partial_tk_{t+s}-\partial_tk_t\bigr)\,\dd s
\]
得到，二阶时间导数亦然。把上述积分恒等式分别取 $L^1$ 范数与 $C_0$ 范数，并使用相应的平移
连续性或时间连续性，便知所有所需核差商同时在 $L^1$ 与 $L^\infty$ 中收敛到对应核导数。
基础型 Young 不等式遂给出所有这些差商卷积在 $L^p$ 中的极限；例如
\[
 \frac{(p_t*f)(\,\cdot+he_j)-p_t*f}{h}
 \longrightarrow(\partial_jp_t)*f
 \quad\hbox{于 }L^p,
\]
对 $p=\infty$ 也成立。因此
\[
 \partial_t(k_t*f)=(\partial_tk_t)*f,\qquad
 \partial_t^2(k_t*f)=(\partial_t^2k_t)*f,\qquad
 \partial_i\partial_j(k_t*f)=(\partial_i\partial_jk_t)*f
\]
作为 $L^p$ 值恒等式成立。另一方面，上述核差商还在 $L^{p'}$ 中收敛：当
$1<p<\infty$ 时由 $L^1\cap L^\infty$ 插值得到，两个端点分别直接使用 $L^\infty$ 与
$L^1$。\Holder{} 不等式因而允许对每个固定 $x$ 在积分中取差商极限，给出所述经典空间导数。
最后把 $\partial_tp_t=\Delta p_t$ 与
$\partial_t^2q_t=-\Delta q_t$ 同 $f$ 卷积，便得到两个方程。这个论证也覆盖不可分的
$L^\infty$ 端点，而没有借助光滑函数稠密性。
\end{proof}

\begin{example}[一维 Poisson 延拓的可计算边界值]\label{ex:4-4-2-3}
在 $n=1$ 时 $c_1=1/\pi$。取 $f=\1_{[-1,1]}$，则
\begin{align*}
 Q_tf(x)
 &=\frac1\pi\int_{-1}^1\frac{t}{t^2+(x-y)^2}\dd y\\
 &=\frac1\pi\left[
   \arctan\frac{1-x}{t}+\arctan\frac{1+x}{t}\right].
\end{align*}
命题~\ref{proposition:4-4-2-4} 表明它在上半平面调和。令 $t\downarrow0$ 可直接读出极限：
在 $|x|<1$ 为 $1$，在 $|x|>1$ 为 $0$，在 $x=\pm1$ 为 $1/2$。端点不是所选代表的
Lebesgue 点，而核给出的半和也没有与 Lebesgue 点定理冲突。
\end{example}

\begin{figure}[htbp]
\centering
\begin{tikzpicture}[x=1cm,y=1cm]
  \draw[->,BookInk!70] (-3.5,0)--(3.7,0) node[right] {$x$};
  \draw[->,BookInk!70] (0,-0.15)--(0,3.4);
  \draw[BookAccent,very thick,domain=-3.2:3.2,samples=100]
    plot (\x,{2.9*exp(-2.2*\x*\x)});
  \draw[BookWarm,thick,domain=-3.2:3.2,samples=100]
    plot (\x,{1.45*exp(-0.55*\x*\x)});
  \draw[BookInk!55,thick,dashed,domain=-3.2:3.2,samples=100]
    plot (\x,{0.9*exp(-0.22*\x*\x)});
  \node[BookAccent] at (0.75,2.85) {$t_1$};
  \node[BookWarm] at (1.45,1.45) {$t_2$};
  \node[BookInk!70] at (2.35,0.86) {$t_3$};
  \node[anchor=west,text width=4.0cm] at (4.2,2.1)
    {$0<t_1<t_2<t_3$。热核变矮、变宽，面积始终为一；反向令 $t\downarrow0$，
     核变窄而读取更局部的平均。};
\end{tikzpicture}
\caption{热核的扩散与集中。图中高度和宽度只作尺度示意，质量守恒由 Gaussian 积分而非图形证明。}
\label{fig:4-4-2-kernel-scales}
\end{figure}

\begin{remark}[Brownian 时间尺度]
测度 $p_t\dd x$ 的二阶矩矩阵为 $2tI_n$。第~6 章定义协方差为
$tI_n$ 的标准 Brownian 运动后，其转移核将是 $p_{t/2}$，生成元为
$\tfrac12\Delta$；$p_t$ 对应的则是空间尺度放大 $\sqrt2$ 后的过程，生成元为 $\Delta$。
这一时间归一化将在 Brownian 转移半群中再次出现。
\end{remark}

\begin{exercise}
\begin{enumerate}
  \item 从定义三条件独立证明定理~\ref{theorem:4-4-2-1}，并指出何处必须排除 $p=\infty$。
  \item 用极坐标和 Beta 积分验证定义~\ref{def:4-4-2-3} 中 $c_n$ 的 Gamma 函数公式。
  \item 不使用 Fourier 变换，按完成平方的计算证明 $p_s*p_t=p_{s+t}$。
  \item 直接验证 $q_t$ 调和，并用例~\ref{ex:4-4-2-3} 检查边界极限。
\end{enumerate}
\end{exercise}

核的卷积律把一段时间的演化拆成两段时间的复合。下一小节把这个结构从热核中抽象出来；
与此同时，短时间差商只对部分函数收敛，这会迫使我们认真记录一个算子的定义域。


\subsection{强连续半群、生成元与能量估计}\label{subsec:4-4-3}
\index{强连续半群}\index{生成元}\index{能量估计}

等式 $H_sH_t=H_{s+t}$ 说明：先扩散 $t$ 单位时间，再扩散 $s$ 单位时间，
与一次扩散 $s+t$ 单位时间相同。但这个代数等式并不保证短时差商
\[
  \frac{T_tf-f}{t}
\]
对每个 $f$ 都收敛。下面三个例子分别考察平移半群、热半群和 $L^\infty$ 端点；随后将定义
强连续性、生成元及其定义域，并严格识别前两个例子中的生成元。

Hille 与 Yosida 的生成元理论刻画产生强连续半群的无界算子。本小节从已知半群出发，证明局部一致
有界、生成元稠密且闭，并导出正实参数上的 Laplace 预解式。逆向生成定理需要更完整的预解式条件。

\begin{example}[平移半群]\label{ex:4-4-3-1}
在 $L^p(\R)$ 上，$1\le p<\infty$，令
\[
  T_tf(x)=f(x+t),\qquad t\ge0.
\]
则 $T_0=I$、$T_sT_t=T_{s+t}$，并且 $\|T_tf\|_p=\|f\|_p$。由定理
\ref{theorem:4-4-1-1}，$\|T_tf-f\|_p\to0$。若 $f$ 有弱导数 $f'\in L^p$，
一维绝对连续代表给出
\[
  \frac{T_tf-f}{t}=\frac1t\int_0^t T_sf'\,\dd s.
\]
这个公式暗示短时速度是 $f'$，但“差商收敛恰好等价于
$f\in W^{1,p}(\R)$”还需要反向论证。
\end{example}

\begin{example}[热半群]\label{ex:4-4-3-2}
令 $T_t=H_t$，并在 $t=0$ 时定义 $H_0=I$。上一小节已经证明
$H_sH_t=H_{s+t}$；对 $1\le p<\infty$，近似恒等定理给出
$H_tf\to f$ 于 $L^p$。核的时间微分恒等式
$\partial_tp_t=\Delta p_t$ 表明光滑函数的短时速度应为 $\Delta f$。
下面将在 $C_c^\infty(\R^n)$ 上证明这一点，而不对完整生成元域作出断言。
\end{example}

\begin{example}[$L^\infty$ 端点的失败]\label{ex:4-4-3-3}
仍取 $T_tf(x)=f(x+t)$，但把空间改为 $L^\infty(\R)$。对
$f=\1_{(0,\infty)}$ 与每个 $t>0$，$T_tf-f$ 在 $(-t,0]$ 上等于一，在其他点
几乎处处为零。因而
\[
  \|T_tf-f\|_\infty=1.
\]
平移公式没有变，但范数变了，时间连续性便消失了。
\end{example}

\begin{definition}[$C_0$ 半群]\label{def:4-4-3-1}
设 $B$ 是 Banach 空间。一族有界线性算子
$(T_t)_{t\ge0}\subset\mathcal L(B)$ 称为 $B$ 上的一个\emph{$C_0$ 半群}，如果
\[
 T_0=I,\qquad T_{s+t}=T_sT_t\quad(s,t\ge0),
\]
且对每个 $f\in B$ 都有
\[
 \|T_tf-f\|_B\longrightarrow0\qquad(t\downarrow0).
\]
最后一条称为强连续性。它是逐向量的范数连续，并不要求
$\|T_t-I\|_{\mathcal L(B)}\to0$。
\end{definition}

定义只说每条轨道在零时刻连续。为了对轨道积分，还需要一个共同的算子范数界。
这个界不是定义的形式后果；它的证明正是 Baire 范畴定理在这里的用武之地。

\begin{proposition}[紧时间区间上的一致有界]\label{proposition:4-4-3-0}
若 $(T_t)_{t\ge0}$ 是 $B$ 上的 $C_0$ 半群，则对每个 $T<\infty$，
\[
  M_T:=\sup_{0\le t\le T}\|T_t\|_{\mathcal L(B)}<\infty.
\]
\end{proposition}

\begin{proof}
先固定 $f\in B$。由 $T_rf\to f$，可取 $\delta>0$ 使
$\|T_rf\|\le\|f\|+1$ 对 $0\le r\le\delta$ 成立。若
$t\in[0,T]$，写成 $t=k\delta+r$，其中
$0\le r<\delta$ 且 $0\le k\le\lceil T/\delta\rceil$，则半群律给出
\[
  \|T_tf\|=\|T_{k\delta}T_rf\|
  \le \|T_{k\delta}\|(\|f\|+1).
\]
令 $N=\lceil T/\delta\rceil$。右边只涉及 $T_{k\delta}$（$0\leq k\leq N$）这有限多个已知有界
算子，因而
\[
 \sup_{0\leq t\leq T}\|T_tf\|
 \leq(\|f\|+1)\max_{0\leq k\leq N}\|T_{k\delta}\|<\infty.
\]

对每个 $t\in[0,T]$ 定义连续标量函数
$F_t:B\to\R$，$F_t(f)=\|T_tf\|$。上一段证明了函数族
$\{F_t:0\le t\le T\}$ 逐点有界。Banach 空间是 Baire 空间，因而定理
\ref{theorem:1-4-1-3} 给出非空开集 $U\subset B$ 和 $C<\infty$，使
\[
  \|T_tf\|\le C\qquad(f\in U, 0\le t\le T).
\]
取 $f_0\in B$ 与 $r>0$ 使 $B(f_0,r)\subset U$。对任意 $\|h\|\le1$，
$f_0$ 与 $f_0+(r/2)h$ 都在 $U$ 中，故
\[
 \frac r2\|T_th\|
 \le \|T_t(f_0+(r/2)h)\|+\|T_tf_0\|
 \le2C.
\]
对 $t\in[0,T]$ 取上确界，得 $M_T\le4C/r$。
\end{proof}

命题~\ref{proposition:4-4-3-0} 还把零时刻的强连续性推到了所有时刻。
对 $h\downarrow0$，
\[
 \|T_{t+h}f-T_tf\|\le\|T_t\|\,\|T_hf-f\|\longrightarrow0,
\]
而当 $0<h<t$ 时，
\[
 \|T_tf-T_{t-h}f\|
 =\|T_{t-h}(T_hf-f)\|
 \le M_t\|T_hf-f\|\longrightarrow0.
\]
因而 $t\mapsto T_tf$ 是 $[0,\infty)$ 到 $B$ 的连续曲线，在每个紧时间区间上
Bochner 可积。

\begin{definition}[生成元]\label{def:4-4-3-2}
$C_0$ 半群 $(T_t)$ 的\emph{生成元} $A$ 定义为
\[
 D(A)=\left\{f\in B:
   \lim_{t\downarrow0}\frac{T_tf-f}{t}
   \text{ 在 $B$ 中存在}\right\},
 \qquad
 Af=\lim_{t\downarrow0}\frac{T_tf-f}{t}.
\]
差商的线性和范数极限的唯一性说明 $D(A)$ 是线性子空间，
$A:D(A)\to B$ 是线性算子。定义没有声称 $D(A)=B$。
\end{definition}

\begin{definition}[收缩半群与 Markov 性]\label{def:4-4-3-3}
若 $\|T_t\|\le1$ 对每个 $t\ge0$ 成立，则称 $(T_t)$ 为\emph{收缩半群}。
若 $B$ 是带逐点序的函数 Banach 空间、$\1\in B$，每个 $T_t$ 保正，且
$T_t\1=\1$，则称它为\emph{Markov 半群}。

在无穷测度空间的 $L^p$（$p<\infty$）中，常数一不是该空间的向量，
所以不把 $T_t\1=\1$ 写成 $L^p$ 内的等式。此时可在 $L^p\cap L^\infty$ 上使用
\emph{sub-Markov 性}：
\[
  0\le f\le1\quad\Longrightarrow\quad0\le T_tf\le1
  \quad\text{a.e.}
\]
\end{definition}

\begin{proposition}[轨道微分]\label{proposition:4-4-3-1}
若 $f\in D(A)$，则 $T_tf\in D(A)$，并且
\[
  AT_tf=T_tAf.
\]
曲线 $t\mapsto T_tf$ 在 $t>0$ 处可微且
\[
  \frac{\dd}{\dd t}T_tf=AT_tf=T_tAf;
\]
在 $t=0$ 上式的导数按右导数理解。
\end{proposition}

\begin{proof}
对 $h>0$，半群律给出
\[
 \frac{T_hT_tf-T_tf}{h}
 =T_t\frac{T_hf-f}{h}\longrightarrow T_tAf.
\]
因为 $T_t$ 有界，极限可经 $T_t$ 传递。这证明
$T_tf\in D(A)$ 及 $AT_tf=T_tAf$，同时给出右导数。

若 $t>0$ 且 $0<h<t$，记 $q_h=(T_hf-f)/h$，则左差商为
\[
 \frac{T_tf-T_{t-h}f}{h}=T_{t-h}q_h.
\]
命题~\ref{proposition:4-4-3-0} 与 $q_h\to Af$ 给出
\[
 \|T_{t-h}(q_h-Af)\|\le M_t\|q_h-Af\|\longrightarrow0.
\]
另一方面，
\[
 \|T_{t-h}Af-T_tAf\|
 \le M_t\|T_hAf-Af\|\longrightarrow0.
\]
所以左差商也趋于 $T_tAf$。
\end{proof}

生成元域不是人为挑出的光滑函数类。下一个平均操作表明，它至少稠密到足以
近似每个向量。

\begin{proposition}[生成元稠密性的 Yosida 平均]\label{proposition:4-4-3-2}
对 $f\in B$ 与 $t>0$，定义
\[
 f_t=\frac1t\int_0^tT_sf\,\dd s.
\]
则 $f_t\in D(A)$、$f_t\to f$ 于 $B$，并且
\[
 Af_t=\frac{T_tf-f}{t}.
\]
\end{proposition}

\begin{proof}
轨道连续性与命题~\ref{proposition:4-4-3-0} 保证定义中的 Bochner 积分存在。
由 Bochner 积分的基本估计，
\[
 \|f_t-f\|
 \le\frac1t\int_0^t\|T_sf-f\|\,\dd s
 \le\sup_{0\le s\le t}\|T_sf-f\|\longrightarrow0.
\]
对 $h>0$，有界算子与 Bochner 积分的交换性
（命题~\ref{proposition:3-4-3-2}）和半群律给出
\begin{align*}
 \frac{T_hf_t-f_t}{h}
 &=\frac1t\left\{
   \frac1h\int_t^{t+h}T_sf\,\dd s
   -\frac1h\int_0^hT_sf\,\dd s\right\}.
\end{align*}
连续 $B$ 值函数在端点附近的平均趋于端点值；例如
\[
 \left\|\frac1h\int_t^{t+h}T_sf\,\dd s-T_tf\right\|
 \le\sup_{t\le s\le t+h}\|T_sf-T_tf\|\longrightarrow0.
\]
对另一个积分使用同样的估计，便得
$(T_hf_t-f_t)/h\to(T_tf-f)/t$。
\end{proof}

对一个只在 $D(A)$ 上定义的线性算子，“闭”是指：只要
$f_n\in D(A)$、$f_n\to f$ 且 $Af_n\to g$ 于 $B$，就有
$f\in D(A)$ 与 $Af=g$。这是其图像在 $B\times B$ 中为闭集的序列表达。

\begin{proposition}[生成元的闭性]\label{proposition:4-4-3-5}
每个 $C_0$ 半群的生成元都稠密定义且为闭算子。
\end{proposition}

\begin{proof}
稠密性已由命题~\ref{proposition:4-4-3-2} 证明。先建立后面需要的轨道积分公式。
对 $f\in D(A)$ 固定 $t>0$，令 $h=t/N$ 且
$q_h=(T_hf-f)/h$。望远镜求和给出严格的代数等式
\[
 T_tf-f
 =\sum_{k=0}^{N-1}T_{kh}(T_hf-f)
 =h\sum_{k=0}^{N-1}T_{kh}q_h.
\]
它与 $h\sum_{k=0}^{N-1}T_{kh}Af$ 之差的范数不超过
$tM_t\|q_h-Af\|$，因而趋于零。连续轨道
$s\mapsto T_sAf$ 的 Riemann 和收敛于其 Bochner 积分，所以
\begin{equation}\label{eq:4-4-3-orbit-integral}
 T_tf-f=\int_0^tT_sAf\,\dd s
 \qquad(f\in D(A)).
\end{equation}

现设 $f_n\in D(A)$、$f_n\to f$ 且 $Af_n\to g$。把
\eqref{eq:4-4-3-orbit-integral} 用于 $f_n$，则
\[
 T_tf_n-f_n=\int_0^tT_sAf_n\,\dd s.
\]
左边趋于 $T_tf-f$，而
\[
 \left\|\int_0^tT_s(Af_n-g)\,\dd s\right\|
 \le tM_t\|Af_n-g\|\longrightarrow0.
\]
因此
\[
 T_tf-f=\int_0^tT_sg\,\dd s.
\]
除以 $t$ 并令 $t\downarrow0$，轨道连续性给出
\[
 \frac{T_tf-f}{t}=\frac1t\int_0^tT_sg\,\dd s\longrightarrow g.
\]
所以 $f\in D(A)$ 且 $Af=g$。
\end{proof}

现在识别前述平移半群与热半群的生成元。

\begin{proposition}[平移生成元的双向识别]
\label{proposition:4-4-3-translation-generator}
设 $1\le p<\infty$，$B=L^p(\R)$，并令 $T_tf(x)=f(x+t)$。则 $(T_t)$ 是等距
$C_0$ 半群，它的生成元满足
\[
  D(A)=W^{1,p}(\R),\qquad Af=f'.
\]
若改用 $S_tf(x)=f(x-t)$，则生成元为 $-f'$，域仍为
$W^{1,p}(\R)$。
\end{proposition}

\begin{proof}
半群律和等距性由平移不变性得到，强连续性是定理
\ref{theorem:4-4-1-1}。

先设 $f\in W^{1,p}(\R)$。由命题
\ref{proposition:4-3-3-ac-representative}，在每个 $I_m=(-m,m)$ 上选取 $f$ 的绝对连续代表
$\widetilde f_m$。在 $I_m$ 上，$\widetilde f_m$ 与 $\widetilde f_{m+1}$ a.e. 相等；两者的差连续，
所以其零集的补集不可含非空开区间，从而它们在 $I_m$ 上处处相等。因此这些代表粘合成
$f$ 的一个局部绝对连续代表，并且对 a.e. $x$
\[
 f(x+t)-f(x)=\int_0^t f'(x+s)\,\dd s.
\]
将这个等式视为 $L^p$ 中的 Bochner 积分，由积分型 Minkowski
不等式和 $L^p$ 平移连续性，
\[
 \left\|\frac{T_tf-f}{t}-f'\right\|_p
 \le\frac1t\int_0^t\|T_sf'-f'\|_p\,\dd s
 \longrightarrow0.
\]
因而 $f\in D(A)$ 且 $Af=f'$。

反过来，设 $f\in D(A)$ 且
$(T_tf-f)/t\to g$ 于 $L^p$。对任意 $\varphi\in C_c^\infty(\R)$，变量代换给出
\begin{align*}
 \int_\R\frac{f(x+t)-f(x)}t\,\varphi(x)\,\dd x
 &=\int_\R f(y)\frac{\varphi(y-t)-\varphi(y)}t\,\dd y.
\end{align*}
右边括号中的差商在一个固定紧集上支撑，并在
$L^{p'}$ 中趋于 $-\varphi'$；当 $p=1$时，这里的 $p'=\infty$
收敛来自 $\varphi'$ 的一致连续性。\Holder{} 不等式于是允许两边取极限，得
\[
 \int_\R g\varphi\,\dd x=-\int_\R f\varphi'\,\dd x.
\]
这正是 $g$ 为 $f$ 的弱导数的定义。故 $f\in W^{1,p}(\R)$且
$f'=g$。对 $S_t$ 把 $t$ 换成 $-t$，差商的极限符号相反。
\end{proof}

\begin{proposition}[热半群在测试函数上的生成元]
\label{proposition:4-4-3-heat-test-generator}
对 $1\le p<\infty$，$(H_t)_{t\ge0}$ 是 $L^p(\R^n)$ 上的保正收缩
$C_0$ 半群，并具有 sub-Markov 性。若 $f\in C_c^\infty(\R^n)$，则
$f\in D(A)$ 且 $Af=\Delta f$。
\end{proposition}

\begin{proof}
$p_t\ge0$ 且 $\int p_t=1$，所以 $H_t$ 保正；基础型 Young 不等式给出
$\|H_tf\|_p\le\|f\|_p$。若 $0\le f\le1$，则
$0\le p_t*f\le\int p_t=1$，故 sub-Markov 性成立。半群律由命题
\ref{proposition:4-4-2-3} 给出，强连续性由定理
\ref{theorem:4-4-2-1} 给出。

现设 $f\in C_c^\infty$。由命题~\ref{proposition:4-4-2-4}，对 $s>0$
$\frac{\dd}{\dd s}H_sf=\Delta(H_sf)=(\Delta p_s)*f$。对每个坐标 $j$，
$\partial_{x_j}^2p_s(x-y)=\partial_{y_j}^2p_s(x-y)$。由于 $f$ 光滑且具紧支撑，
在 $y_j$ 上连续作两次分部积分，边界项均为零，从而
\begin{align*}
 ((\Delta p_s)*f)(x)
 &=\sum_{j=1}^n\int_{\R^n}
   \partial_{y_j}^2p_s(x-y)f(y)\,\dd y\\
 &=\sum_{j=1}^n\int_{\R^n}
   p_s(x-y)\partial_{y_j}^2f(y)\,\dd y
 =(H_s\Delta f)(x).
\end{align*}
所以对 $s>0$
\[
 \frac{\dd}{\dd s}H_sf=H_s\Delta f
\]
在 $L^p$ 中成立。固定 $0<\varepsilon<t$，对连续的 Banach 值导函数作 Riemann 和
并用有限和的望远镜等式，得
\[
 H_tf-H_\varepsilon f=\int_\varepsilon^tH_s\Delta f\,\dd s.
\]
令 $\varepsilon\downarrow0$。近似恒等定理分别用于 $f$ 和 $\Delta f$，给出
\begin{equation}\label{eq:4-4-3-heat-difference-integral}
 H_tf-f=\int_0^tH_s\Delta f\,\dd s.
\end{equation}
因此
\[
 \left\|\frac{H_tf-f}{t}-\Delta f\right\|_p
 \le\frac1t\int_0^t\|H_s\Delta f-\Delta f\|_p\,\dd s
 \longrightarrow0.
\]
这只识别了 $C_c^\infty$ 上的生成元，不包含对 $D(A)$ 的完整描述。
\end{proof}

上述两例也显示了为什么生成元通常无界。例如取非零
$\eta\in C_c^\infty(\R)$ 与 $f_k(x)=e^{ikx}\eta(x)$，则 $\|f_k\|_p=\|\eta\|_p$，而
\[
 \|f_k'\|_p\ge k\|\eta\|_p-\|\eta'\|_p\longrightarrow\infty.
\]
“$A$ 是线性的”不能与“$A$ 在整个 $B$ 上有界”混为一谈。

\begin{theorem}[热能量恒等式]\label{theorem:4-4-3-3}
设 $f\in C_c^\infty(\R^n)$、$u(t)=H_tf$。对每个 $t>0$，
\[
 \frac{\dd}{\dd t}\|u(t)\|_2^2
 =-2\|\nabla u(t)\|_2^2\le0.
\]
更一般地，若 $f\in L^2(\R^n)$ 且 $b>0$，则
\begin{equation}\label{eq:4-4-3-integrated-energy}
 \|H_bf\|_2^2+2\int_0^b\|\nabla H_sf\|_2^2\,\dd s
 =\|f\|_2^2.
\end{equation}
后一式是时间积分恒等式；它不声称粗糙 $f$ 在 $t=0$ 具有点态时间导数。
\end{theorem}

\begin{proof}
先取 $f\in C_c^\infty$。命题~\ref{proposition:4-4-2-4} 给出
$\partial_tu=\Delta u$，并且对每个 $t>0$，$u(t)$ 及其所需导数都由 Gaussian
尾部支配。因此 $t\mapsto u(t)$ 在 $L^2$ 中可微，且可将空间积分先限制到
$B(0,R)$ 再令 $R\to\infty$。对复值函数计算得
\begin{align*}
 \frac{\dd}{\dd t}\|u(t)\|_2^2
 &=2\operatorname{Re}\int_{\R^n}(\partial_tu)\overline u\,\dd x\\
 &=2\operatorname{Re}\int_{\R^n}(\Delta u)\overline u\,\dd x\\
 &=-2\sum_{j=1}^n\int_{\R^n}|\partial_ju|^2\,\dd x
 =-2\|\nabla u(t)\|_2^2.
\end{align*}
第三行可先在 $B(0,R)$ 上分部积分。若 $\supp f\subset B(0,R_0)$，则对固定 $t>0$ 及
$|x|>2R_0$，由 Gaussian 核及其一阶导数的显式公式可取常数 $C_t,c_t>0$，使
\[
 |u(t,x)|+|\nabla u(t,x)|\leq C_te^{-c_t|x|^2}.
\]
因此球面边界项的绝对值不超过
$C_tR^{n-1}e^{-2c_tR^2}\to0$，令 $R\to\infty$ 即得到第三行。
在 $[a,b]\subset(0,\infty)$ 上积分即得
\begin{equation}\label{eq:4-4-3-energy-ab-smooth}
 \|H_bf\|_2^2+2\int_a^b\|\nabla H_sf\|_2^2\,\dd s
 =\|H_af\|_2^2.
\end{equation}

对一般 $f\in L^2$，由光滑稠密性取 $f_k\in C_c^\infty$ 使
$f_k\to f$ 于 $L^2$。收缩性给出
$H_tf_k\to H_tf$ 于 $L^2$，并且对 $s\in[a,b]$，Young 不等式给出
\[
 \|\nabla H_s(f_k-f)\|_2
 \le\|\nabla p_s\|_1\|f_k-f\|_2
 \le a^{-1/2}\|\nabla p_1\|_1\|f_k-f\|_2.
\]
最后一步使用了缩放恒等式
$\nabla p_s(x)=s^{-(n+1)/2}(\nabla p_1)(x/\sqrt s)$。上式在 $s\in[a,b]$ 上一致；
同时 $(\|f_k\|_2)$ 有界，所以存在 $C<\infty$ 使
\[
 \sup_{k,\,s\in[a,b]}\|\nabla H_sf_k\|_2
 +\sup_{s\in[a,b]}\|\nabla H_sf\|_2\le C.
\]
于是
\begin{align*}
 &\left|\int_a^b\|\nabla H_sf_k\|_2^2\,\dd s
 -\int_a^b\|\nabla H_sf\|_2^2\,\dd s\right|\\
 &\qquad\le
 \int_a^b\bigl(\|\nabla H_sf_k\|_2+\|\nabla H_sf\|_2\bigr)
 \|\nabla H_s(f_k-f)\|_2\,\dd s\\
 &\qquad\le 2C(b-a)a^{-1/2}\|\nabla p_1\|_1\|f_k-f\|_2
 \longrightarrow0.
\end{align*}
因此可在 \eqref{eq:4-4-3-energy-ab-smooth} 中令 $k\to\infty$，得到对 $f$ 成立的同一等式。
最后令 $a\downarrow0$。近似恒等定理给出 $H_af\to f$ 于 $L^2$，而非负函数
$s\mapsto\|\nabla H_sf\|_2^2$ 的积分由单调收敛定理趋于 $(0,b)$ 上的积分。
这就得到 \eqref{eq:4-4-3-integrated-energy}。
\end{proof}

半群记录了所有正时间的演化，生成元记录了零时刻的速度。还有第三种记录方式：
以指数权把整条轨道平均起来。它会把无界算子问题转成一个定义在整个 $B$ 上的有界算子。

\begin{proposition}[预解算子的 Laplace 表示]\label{proposition:4-4-3-4}
设 $(T_t)$ 是 $B$ 上的收缩 $C_0$ 半群，$A$ 是其生成元。对 $\lambda>0$，定义
\[
 R_\lambda f=\int_0^\infty e^{-\lambda t}T_tf\,\dd t.
\]
则 $R_\lambda\in\mathcal L(B)$、$\|R_\lambda\|\le\lambda^{-1}$，并且
\[
 R_\lambda B\subset D(A),\qquad
 (\lambda I-A)R_\lambda f=f\quad(f\in B),
\]
\[
 R_\lambda(\lambda I-A)g=g\quad(g\in D(A)).
\]
因而 $\lambda I-A:D(A)\to B$ 是双射，其有界逆是 $R_\lambda$。

若 $B$ 是复 Banach 空间，定义
\[
 \rho(A)=\{z\in\C:zI-A:D(A)\to B
 \text{ 为双射且逆算子在 $B$ 上有界}\}.
\]
则 $(0,\infty)\subset\rho(A)$，且
\[
 R(\lambda,A)=(\lambda I-A)^{-1}=R_\lambda,
 \qquad\|R(\lambda,A)\|\le\frac1\lambda.
\]
\end{proposition}

\begin{proof}
连续轨道使 $t\mapsto e^{-\lambda t}T_tf$ 强可测，而收缩性给出
\[
 \int_0^\infty\|e^{-\lambda t}T_tf\|\,\dd t
 \le\|f\|\int_0^\infty e^{-\lambda t}\,\dd t
 =\frac{\|f\|}{\lambda}.
\]
所以广义 Bochner 积分存在，基本估计给出
$\|R_\lambda\|\le1/\lambda$。

由命题~\ref{proposition:3-4-3-2}，对 $h>0$
\begin{align*}
 T_hR_\lambda f
 &=\int_0^\infty e^{-\lambda t}T_{t+h}f\,\dd t\\
 &=e^{\lambda h}\int_h^\infty e^{-\lambda s}T_sf\,\dd s\\
 &=e^{\lambda h}\left(R_\lambda f-
     \int_0^h e^{-\lambda s}T_sf\,\dd s\right).
\end{align*}
因此
\[
 \frac{T_hR_\lambda f-R_\lambda f}{h}
 =\frac{e^{\lambda h}-1}{h}R_\lambda f
 -e^{\lambda h}\frac1h\int_0^he^{-\lambda s}T_sf\,\dd s.
\]
第一项趋于 $\lambda R_\lambda f$；第二项中的平均趋于 $f$，因为
$e^{-\lambda s}T_sf\to f$ 于 $B$。故 $R_\lambda f\in D(A)$ 且
\[
 AR_\lambda f=\lambda R_\lambda f-f,
\]
这就是 $(\lambda I-A)R_\lambda=I$。

再取 $g\in D(A)$。对 $R>0$ 记
\[
 u_R=\int_0^R e^{-\lambda t}T_tg\,\dd t,
 \qquad
 v_R=\int_0^R e^{-\lambda t}T_tAg\,\dd t.
\]
由命题~\ref{proposition:4-4-3-1}，$T_tg\in D(A)$ 且
$AT_tg=T_tAg$。对上述两个连续轨道取同一列带权 Riemann 和
$u_{R,N}$ 与 $v_{R,N}$。每个 $u_{R,N}$ 都属于 $D(A)$，且线性性给出
\[
 Au_{R,N}=v_{R,N}.
\]
轨道的一致连续性使 $u_{R,N}\to u_R$、$v_{R,N}\to v_R$ 于 $B$。
命题~\ref{proposition:4-4-3-5} 的闭性于是给出
\begin{equation}\label{eq:4-4-3-resolvent-commutes-generator-finite}
 u_R\in D(A),\qquad Au_R=v_R.
\end{equation}
当 $R\to\infty$ 时，收缩性与指数可积性给出
\[
 u_R\longrightarrow R_\lambda g,
 \qquad v_R\longrightarrow R_\lambda Ag
 \quad\text{于 }B.
\]
再用一次闭性于 \eqref{eq:4-4-3-resolvent-commutes-generator-finite}，得
\[
 AR_\lambda g=R_\lambda Ag.
\]
已证的左逆恒等式应用于 $g$ 后为
$\lambda R_\lambda g-AR_\lambda g=g$，代入上式便得
\[
 R_\lambda(\lambda I-A)g=g.
\]
轨道微分命题还给出
\[
 \frac{\dd}{\dd t}\bigl(e^{-\lambda t}T_tg\bigr)
 =-e^{-\lambda t}T_t(\lambda g-Ag).
\]
因而上述右逆恒等式正是这条微分式在 $[0,\infty)$ 上的积分形式；
前面的闭性与 Riemann 和论证已经为该积分形式给出了无需借助形式计算的证明。
两个乘积恒等式分别给出满射性和单射性，其余结论随之得到。
\end{proof}

\begin{example}[标量半群]
若 $T_t=e^{-ct}I$，其中 $c\ge0$，则
\[
 \frac{T_tf-f}{t}=\frac{e^{-ct}-1}{t}f\longrightarrow-cf,
\]
所以 $A=-cI$ 且 $D(A)=B$。直接积分得
\[
 R_\lambda=\int_0^\infty e^{-(\lambda+c)t}\,\dd t\,I
 =\frac1{\lambda+c}I.
\]
因此预解算子在两侧都消去 $\lambda I-A=(\lambda+c)I$，并且
$\|R_\lambda\|=(\lambda+c)^{-1}\le\lambda^{-1}$。
\end{example}

\begin{remark}[与 Hille--Yosida 定理的关系]
本小节从一个已给定的收缩 $C_0$ 半群出发，证明了它的生成元稠密、闭，且
\[
  \|\lambda R(\lambda,A)\|\le1\qquad(\lambda>0).
\]
这只是 Hille--Yosida 必要条件的第一层，不包含高次预解式估计。
反过来，从一个闭算子的预解式条件构造收缩 $C_0$ 半群，是一般 Hille--Yosida 生成定理的困难方向，
需要高次预解式估计。
\end{remark}

\begin{remark}[实空间与复预解集]
命题~\ref{proposition:4-4-3-4} 中对正实参数 $\lambda$ 的逆公式同样适用于实 Banach 空间。
但 $\rho(A)\subset\C$ 是复线性理论的记号；在没有先作复化时，不把实空间上的算子
直接代入 $zI-A$。
\end{remark}

\begin{figure}[htbp]
\centering
\begin{tikzpicture}
  \draw[book line,-{Stealth[length=2.4mm]}] (-0.3,1.8)--(12.0,1.8)
    node[right,book label]{时间};
  \foreach \x/\lab in {0/$0$,4/$t$,10.5/$t+s$}{
    \draw[book line] (\x,1.68)--(\x,1.92);
    \node[book label,above] at (\x,1.95) {\lab};
  }
  \draw[book accent arrow] (0,1.35)--node[below,book label]{$T_t$}(4,1.35);
  \draw[book accent arrow] (4,1.35)--node[below,book label]{$T_s$}(10.5,1.35);
  \draw[book dashed arrow] (0,0.75)--node[below,book label]{$T_{t+s}=T_sT_t$}(10.5,0.75);

  \node[book node] (all) at (0.8,-0.65) {$f\in B$};
  \node[book strong node] (avg) at (5.1,-0.65)
    {$f_h=\dfrac1h\int_0^hT_rf\,\dd r\in D(A)$};
  \node[book warm node] (vel) at (11.1,-0.65)
    {$Af_h=\dfrac{T_hf-f}{h}$};
  \draw[book accent arrow] (all)--node[above,book label,align=center,font=\scriptsize]{Yosida\\平均}(avg);
  \draw[book arrow] (avg)--node[above,book label,align=center,font=\scriptsize]{短时\\速度}(vel);
\end{tikzpicture}
\caption{半群律拼接时间段；Yosida 平均则把任意向量送入生成元域，且保留可计算的差商。}
\label{fig:4-4-3-semigroup-yosida}
\end{figure}

\begin{remark}[Markov 性与转移核]
第~6 章引入过程与转移核以后，概率核 $K_t(x,\dd y)$ 将产生算子
\[
 T_tf(x)=\int f(y)K_t(x,\dd y).
\]
保正性和质量一分别对应保正与常数不变，Chapman--Kolmogorov 公式对应半群律。
坐标函数 $x_i$ 和二次函数 $x_ix_j$ 未必属于所考察的 $L^p$ 或 $C_0$ 空间；短时矩计算需要先截断，
再在矩的一致控制下撤去截断。热核 $p_t$ 与 $p_{t/2}$ 对应的生成元分别为
$\Delta$ 与 $\tfrac12\Delta$，这与上一小节的二阶矩计算一致。
\end{remark}

\begin{exercise}
\begin{enumerate}
  \item 对 $S_tf(x)=f(x-t)$ 重做命题~\ref{proposition:4-4-3-translation-generator}
        的双向论证，并在测试函数配对中跟踪负号。
  \item 不引用命题~\ref{proposition:4-4-3-1}，从半群律直接证明
        $T_tD(A)\subset D(A)$ 与轨道微分公式。
  \item 若 $f_t=t^{-1}\int_0^tT_sf\,\dd s$，证明
        $\|f_t-f\|\le\sup_{0\le s\le t}\|T_sf-f\|$，并逐步计算
        $(T_hf_t-f_t)/h$ 的极限。
  \item 对 $f\in C_c^\infty(\R^n)$ 从 $\partial_tp_t=\Delta p_t$ 出发验证
        \eqref{eq:4-4-3-heat-difference-integral}，再完成热能量恒等式中的分部积分边界估计。
  \item 从 Laplace 积分直接证明
        $\|R_\lambda\|\le\lambda^{-1}$ 与两个逆恒等式，并说明这是
        Hille--Yosida 理论中从半群到生成元的方向。
\end{enumerate}
\end{exercise}

至此，测度、积分、核与时间演化已经放在同一张图景中。第~5 章将把可测集读作事件、
把积分读作期望，并开始追问“已知部分信息”如何改变我们对同一无穷对象的观测。

